% Les révolutions scientifiques et techniques depuis la fin de la Deuxième Guerre mondiale — Histoire Tle Tech — Cours signé OnBuch+
% Programme officiel MINESEC. Compiler avec : tectonic N05-revolutions-scientifiques-techniques-depuis-fin-deuxiem.tex
\newcommand{\DOCMATIERE}{Histoire}
\newcommand{\DOCNIVEAU}{Tle Tech}
\newcommand{\DOCMODULE}{Histoire – classe de Terminale technique}
\newcommand{\DOCLECON}{N05}
\newcommand{\DOCTITRE}{Les révolutions scientifiques et techniques depuis la fin de la Deuxième Guerre mondiale}
% OnBuch+ — préambule « cours signés » ENRICHI (compilé avec tectonic / XeLaTeX)
% Système visuel « Pop » d'OnBuch+ : fond crème (jamais blanc), bordures encre
% épaisses, ombres « dures » décalées, titres Archivo Black en capitales.
% Version MATHÉMATIQUES : graphes (pgfplots/tikz), boîtes pédagogiques
% (définition, propriété, méthode, attention, le savais-tu, exercice résolu,
% exercice, TP, bilan), page de garde et signature OnBuch+ ancrée en bas.
%
% Macros attendues dans main.tex AVANT \input{preamble} :
%   \DOCMATIERE  \DOCNIVEAU  \DOCMODULE  \DOCLECON (numéro)  \DOCTITRE
\documentclass[11pt]{article}

\usepackage{fontspec}
\usepackage[french]{babel}
\usepackage[a4paper,margin=2cm,top=3.4cm,bottom=2.8cm,headheight=1.4cm,footskip=1.3cm]{geometry}
\usepackage{amsmath,amssymb,amsfonts}
\usepackage{mathtools} % \xmapsto, \coloneqq... souvent utilisés par les modèles
\usepackage{cancel} % simplifications barrées (bilans de forces)
% \abs{z} (module, valeur absolue) et \norm{u} (norme), utilisés par les modèles
\providecommand{\abs}{}\let\abs\relax\DeclarePairedDelimiter{\abs}{\lvert}{\rvert}
\providecommand{\norm}{}\let\norm\relax\DeclarePairedDelimiter{\norm}{\lVert}{\rVert}
\usepackage[table]{xcolor} % [table] : fournit aussi \rowcolors (alternance de lignes)
\usepackage{pagecolor}
\usepackage{tikz}
\usetikzlibrary{arrows.meta,positioning,calc,shapes.geometric,shapes.misc,shapes.arrows,decorations.pathmorphing,decorations.pathreplacing,decorations.markings,patterns,shadows,mindmap,trees,fit,backgrounds,matrix,intersections,angles,quotes}
\usepackage{pgfplots}
\usepackage[version=4]{mhchem} % équations nucléaires \ce{^{235}_{92}U}
\usepackage[european,siunitx]{circuitikz} % schémas électriques
\pgfplotsset{compat=1.18}
\usepgfplotslibrary{fillbetween,groupplots} % aires entre courbes (calcul intégral)
\usepackage{enumitem}
\usepackage{fancyhdr}
% [most] charge toutes les bibliothèques tcolorbox (listings, minted, raster,
% documentation...) et gonfle inutilement la mémoire TeX sur un document long
% (~300 boîtes) : on ne charge que ce qui est réellement utilisé.
\usepackage[skins,breakable]{tcolorbox}
\usepackage{graphicx}
\usepackage{colortbl}
\usepackage{booktabs}
\usepackage{tabularx}
\usepackage{multirow}
\usepackage{diagbox} % case d'en-tête diagonale des tableaux à double entrée
% Colonnes extensibles centrée / gauche / droite (C, L, R), souvent utilisées par les modèles
\newcolumntype{C}{>{\centering\arraybackslash}X}
\newcolumntype{L}{>{\raggedright\arraybackslash}X}
\newcolumntype{R}{>{\raggedleft\arraybackslash}X}
\usepackage{array}
\usepackage{multicol}
\usepackage{siunitx}
% parse-numbers=false : accepte aussi \SI{2,5 \times 10^{-3}}{...}
\sisetup{locale=FR,output-decimal-marker={,},group-minimum-digits=4,parse-numbers=false}
\DeclareSIUnit\ohm{\ensuremath{\symup{\Omega}}} % U+2126 absent de la police
\DeclareSIUnit\division{div} % réglages d'oscilloscope (ms/div, V/div)
\DeclareSIUnit\tr{tr} % tours (vitesse de rotation, ex. \SI{1500}{\tr\per\minute})
\DeclareSIUnit\molar{M} % concentration molaire (ex. \SI{3}{\molar}, \SI{1}{\micro\molar})
\DeclareSIUnit\an{an} % année (le modèle confond parfois avec \year, un compteur TeX) — ex. \SI{5}{\kilo\meter\per\an}
\DeclareSIUnit\FCFA{FCFA} % franc CFA (ex. \SI{500}{\FCFA\per\kilo\gram})
\DeclareSIUnit\degreeBrix{\textdegree Brix} % teneur en sucre d'un sirop/jus (réfractométrie)
\DeclareSIUnit\month{mois} % le modèle utilise parfois \month, un compteur TeX, comme unité de durée
\DeclareSIUnit\microliter{\micro\liter} % le modèle écrit parfois \microliter en un seul token
\DeclareSIUnit\million{millions} % dénombrement (ex. \SI{15}{\million\per\mL} spermatozoïdes)
\usepackage{qrcode}
% \degree : commande fréquemment utilisée par les modèles pour un angle ou une
% température en dehors de \SI{}{} (non fournie par les paquets chargés).
\providecommand{\degree}{^{\circ}}
% \qed (fin de démonstration), utilisé par les modèles sans amsthm
\providecommand{\qed}{\hfill$\square$}
% \barre : double barre des valeurs interdites dans un tableau de variations (tkz-tab)
\providecommand{\barre}{\ensuremath{\|}}
% environnement proof (amsthm non chargé) utilisé par les modèles
\ifdefined\proof\else\newenvironment{proof}[1][Démonstration]{\par\noindent\textit{#1.}\ }{\qed\par}\fi
\usepackage{lastpage}
\usepackage{hyperref}
\hypersetup{hidelinks}

% ============================================================
% Palette « Pop » OnBuch+ — mêmes hex que la classe P (plus_ui.dart)
% ============================================================
\definecolor{poporange}{HTML}{F59321}
\definecolor{popdark}{HTML}{E07A0C}
\definecolor{popcream}{HTML}{FFF6E8}
\definecolor{popink}{HTML}{1C1714}
\definecolor{poppurple}{HTML}{7A5AE0}
\definecolor{popgreen}{HTML}{2E9E6A}
\definecolor{poppink}{HTML}{E0457B}
\definecolor{popblue}{HTML}{2F7DE1}
\definecolor{popgold}{HTML}{C98A12}
\definecolor{popmuted}{HTML}{8A7F76}
\definecolor{popline}{HTML}{EADFD2}
\definecolor{popgray}{HTML}{8A7F76}
\colorlet{popgrey}{popgray}
% Alias tolérants (le modèle nomme parfois les couleurs différemment)
\colorlet{popwhite}{white}
\colorlet{popblack}{popink}
\colorlet{popred}{poppink}
\colorlet{popyellow}{popgold}
% Teintes claires (fonds de tableaux/encadrés) — le modèle nomme parfois
% les couleurs avec un suffixe L (ex. popblueL) sans les avoir définies.
\colorlet{poporangeL}{poporange!20}
\colorlet{popdarkL}{popdark!20}
\colorlet{poppurpleL}{poppurple!20}
\colorlet{popgreenL}{popgreen!20}
\colorlet{poppinkL}{poppink!20}
\colorlet{popblueL}{popblue!20}
\colorlet{popgoldL}{popgold!20}
\colorlet{popredL}{popred!20}
\colorlet{popgrayL}{popgray!20}
% Teintes claires pour fonds de tableaux / zones
\colorlet{popblueL}{popblue!10!white}
\colorlet{poporangeL}{poporange!14!white}
\colorlet{popgreenL}{popgreen!12!white}
\colorlet{poppurpleL}{poppurple!12!white}
\colorlet{poppinkL}{poppink!12!white}
\colorlet{popgoldL}{popgold!14!white}

\pagecolor{popcream}

% ============================================================
% Typographie
% ============================================================
\defaultfontfeatures{Ligatures=TeX}
\setmainfont{PlusJakartaSans-Medium.ttf}[Path=../../../fonts/,BoldFont=PlusJakartaSans-Bold.ttf]
% Maths en sans-serif (Fira Math) avec chiffres et lettres droites de Plus
% Jakarta, pour que les formules se fondent dans le texte.
\usepackage[math-style=ISO,bold-style=ISO]{unicode-math}
\setmathfont{FiraMath-Regular.otf}[Path=../../../fonts/]
\setmathfont{PlusJakartaSans-Medium.ttf}[Path=../../../fonts/,range={up/num,up/{Latin,latin},\mathup}]
\usepackage{icomma} % virgule décimale sans espace en mode math
% Caractères mathématiques Unicode absents des polices de texte (Plus Jakarta,
% Archivo Black) : rendus par leur équivalent mathématique, en texte comme en
% formule — sinon ils s'affichent en carré vide (ex. « Arithmétique dans ℤ »).
\usepackage{newunicodechar}
\newunicodechar{ℝ}{\ensuremath{\mathbb{R}}}
\newunicodechar{ℤ}{\ensuremath{\mathbb{Z}}}
\newunicodechar{ℂ}{\ensuremath{\mathbb{C}}}
\newunicodechar{ℕ}{\ensuremath{\mathbb{N}}}
\newunicodechar{ℚ}{\ensuremath{\mathbb{Q}}}
\newunicodechar{𝔻}{\ensuremath{\mathbb{D}}}
\newunicodechar{α}{\ensuremath{\alpha}}
\newunicodechar{β}{\ensuremath{\beta}}
\newunicodechar{γ}{\ensuremath{\gamma}}
\newunicodechar{δ}{\ensuremath{\delta}}
\newunicodechar{ε}{\ensuremath{\varepsilon}}
\newunicodechar{θ}{\ensuremath{\theta}}
\newunicodechar{λ}{\ensuremath{\lambda}}
\newunicodechar{μ}{\ensuremath{\mu}}
\newunicodechar{σ}{\ensuremath{\sigma}}
\newunicodechar{τ}{\ensuremath{\tau}}
\newunicodechar{φ}{\ensuremath{\varphi}}
\newunicodechar{ψ}{\ensuremath{\psi}}
\newunicodechar{ω}{\ensuremath{\omega}}
\newunicodechar{Σ}{\ensuremath{\Sigma}}
% majuscules produites par \MakeUppercase dans les titres (α-aminés -> Α)
\newunicodechar{Α}{\ensuremath{\alpha}}
\newunicodechar{Β}{\ensuremath{\beta}}
\newunicodechar{Γ}{\ensuremath{\Gamma}}
\newunicodechar{Δ}{\ensuremath{\Delta}}
\newunicodechar{Ε}{\ensuremath{\varepsilon}}
\newunicodechar{Θ}{\ensuremath{\Theta}}
\newunicodechar{Λ}{\ensuremath{\Lambda}}
\newunicodechar{Μ}{\ensuremath{\mu}}
\newunicodechar{Τ}{\ensuremath{\tau}}
\newunicodechar{Φ}{\ensuremath{\Phi}}
\newunicodechar{Ψ}{\ensuremath{\Psi}}
\newunicodechar{Ω}{\ensuremath{\Omega}}
\newunicodechar{↦}{\ensuremath{\mapsto}}
\newunicodechar{∈}{\ensuremath{\in}}
\newunicodechar{∉}{\ensuremath{\notin}}
\newunicodechar{∧}{\ensuremath{\wedge}}
\newunicodechar{⇔}{\ensuremath{\Leftrightarrow}}
\newunicodechar{⟺}{\ensuremath{\Longleftrightarrow}}
\newunicodechar{⇒}{\ensuremath{\Rightarrow}}
\newunicodechar{⟹}{\ensuremath{\Longrightarrow}}
\newunicodechar{∘}{\ensuremath{\circ}}
\newunicodechar{∪}{\ensuremath{\cup}}
\newunicodechar{∩}{\ensuremath{\cap}}
\newunicodechar{≡}{\ensuremath{\equiv}}
\newunicodechar{⊂}{\ensuremath{\subset}}
\newunicodechar{⊥}{\ensuremath{\perp}}
\newunicodechar{∃}{\ensuremath{\exists}}
\newunicodechar{∀}{\ensuremath{\forall}}
\newunicodechar{∛}{\ensuremath{\sqrt[3]{\,}}}
\newunicodechar{∜}{\ensuremath{\sqrt[4]{\,}}}
\newunicodechar{⃗}{\ensuremath{{}^{\to}}}
\newunicodechar{ⁿ}{\ifmmode^{n}\else\textsuperscript{\ensuremath{n}}\fi}
\newunicodechar{ˣ}{\ifmmode^{x}\else\textsuperscript{\ensuremath{x}}\fi}
\newunicodechar{ᵇ}{\ifmmode^{b}\else\textsuperscript{\ensuremath{b}}\fi}
\newunicodechar{ʳ}{\ifmmode^{r}\else\textsuperscript{\ensuremath{r}}\fi}
\newunicodechar{ᵘ}{\ifmmode^{u}\else\textsuperscript{\ensuremath{u}}\fi}
\newunicodechar{ʸ}{\ifmmode^{y}\else\textsuperscript{\ensuremath{y}}\fi}
\newunicodechar{ᵃ}{\ifmmode^{a}\else\textsuperscript{\ensuremath{a}}\fi}
\newunicodechar{ᵉ}{\ifmmode^{e}\else\textsuperscript{\ensuremath{e}}\fi}
\newunicodechar{ᵏ}{\ifmmode^{k}\else\textsuperscript{\ensuremath{k}}\fi}
\newunicodechar{ᵖ}{\ifmmode^{p}\else\textsuperscript{\ensuremath{p}}\fi}
\newunicodechar{ᴺ}{\ifmmode^{N}\else\textsuperscript{\ensuremath{N}}\fi}
\newunicodechar{ᶜ}{\ifmmode^{c}\else\textsuperscript{\ensuremath{c}}\fi}
\newunicodechar{ʰ}{\ifmmode^{h}\else\textsuperscript{\ensuremath{h}}\fi}
\newunicodechar{ᵠ}{\ifmmode^{q}\else\textsuperscript{\ensuremath{q}}\fi}
\newunicodechar{ₙ}{\ifmmode_{n}\else\textsubscript{\ensuremath{n}}\fi}
\newunicodechar{ₐ}{\ifmmode_{a}\else\textsubscript{\ensuremath{a}}\fi}
\newunicodechar{ᵢ}{\ifmmode_{i}\else\textsubscript{\ensuremath{i}}\fi}
\newunicodechar{ᵣ}{\ifmmode_{r}\else\textsubscript{\ensuremath{r}}\fi}
\newunicodechar{ₖ}{\ifmmode_{k}\else\textsubscript{\ensuremath{k}}\fi}
\newunicodechar{ᵦ}{\ifmmode_{\beta}\else\textsubscript{\ensuremath{\beta}}\fi}
\newunicodechar{ₓ}{\ifmmode_{x}\else\textsubscript{\ensuremath{x}}\fi}
\newfontfamily\popdisplay{ArchivoBlack-Regular.ttf}[Path=../../../fonts/]
\newfontfamily\poplogo{SpaceGrotesk-Bold.ttf}[Path=../../../fonts/]
\newfontfamily\popsemi{PlusJakartaSans-SemiBold.ttf}[Path=../../../fonts/]
\newfontfamily\popxbold{PlusJakartaSans-ExtraBold.ttf}[Path=../../../fonts/]
\newfontfamily\popxboldls{PlusJakartaSans-ExtraBold.ttf}[Path=../../../fonts/,LetterSpace=3.0]

\setlist[enumerate]{leftmargin=1.7em,itemsep=5pt,topsep=4pt}
\setlist[itemize]{leftmargin=1.7em,itemsep=4pt,topsep=4pt,label={\color{poporange}\textbullet}}
\setlength{\parindent}{0pt}
\setlength{\parskip}{5pt}
\renewcommand{\arraystretch}{1.25}

% pgfplots : style Pop par défaut
\pgfplotsset{
  every axis/.append style={axis line style={popink,line width=1pt},tick style={popink},
    grid style={popline,line width=0.5pt},label style={font=\small},tick label style={font=\footnotesize}},
  popcurve/.style={poporange,line width=1.8pt,no markers,smooth},
}
% Même style disponible dans un \draw TikZ simple (hors pgfplots)
\tikzset{popcurve/.style={poporange,line width=1.8pt}}
\tikzset{popgrid/.style={popline,very thin}}
\colorlet{popcurve}{poporange} % le nom du style est parfois utilisé comme couleur

% ============================================================
% Pill / squiggle
% ============================================================
\newcommand{\poppill}[3]{%
  \tikz[baseline=-0.55ex]\node[draw=popink,line width=1.2pt,rounded corners=2.6pt,%
    fill=#2,inner xsep=6.5pt,inner ysep=3pt]{%
    \popxboldls\fontsize{7}{7}\selectfont\color{#3}\MakeUppercase{#1}};%
}
\newcommand{\popsquiggle}[1]{%
  \tikz[baseline=-0.4ex]{
    \draw[#1,line width=2.6pt,line cap=round]
      plot [smooth] coordinates {(0,0.09) (0.34,0.01) (0.68,0.21) (1.02,0.01) (1.36,0.10)};
  }%
}
% Mot-clé surligné (marqueur orange sous le texte)
\newcommand{\cle}[1]{{\popxbold\color{popdark}#1}}

% ============================================================
% En-tête / pied
% ============================================================
\pagestyle{fancy}
\fancyhf{}
\renewcommand{\headrulewidth}{0pt}
\renewcommand{\footrulewidth}{0pt}
\lhead{\raisebox{-0.15ex}{\poplogo\fontsize{13}{13}\selectfont\color{popink}OnBuch\color{poporange}+}}
\rhead{\poppill{\DOCMATIERE\ \textbullet\ \DOCNIVEAU\ \textbullet\ Leçon \DOCLECON}{white}{popink}}
\lfoot{\popxbold\fontsize{7.6}{7.6}\selectfont\color{popmuted}\MakeUppercase{Cours signé OnBuch+}\ \textbullet\ {\popsemi\DOCTITRE}}
\rfoot{\tikz[baseline=-0.6ex]\node[rounded corners=8.5pt,fill=popink,minimum height=17pt,minimum width=17pt,inner xsep=5pt,inner sep=0]{\popxbold\fontsize{8}{8}\selectfont\color{popcream}\thepage\,/\,\pageref*{LastPage}};}

% ============================================================
% Titres
% ============================================================
\newcommand{\chaptitre}[2]{%
  \begin{tcolorbox}[enhanced,colback=poporange,colframe=popink,boxrule=2.6pt,arc=11pt,
      drop shadow=popink,left=15pt,right=15pt,top=12pt,bottom=11pt,before skip=3pt,after skip=18pt]
    \poppill{#2}{popink}{popcream}\par\vspace{9pt}
    {\raggedright\hyphenpenalty=10000\exhyphenpenalty=10000\popdisplay\fontsize{22}{24}\selectfont\color{popcream}\MakeUppercase{#1}\par}
  \end{tcolorbox}
}
% Section numérotée : pastille orange avec le numéro + titre Archivo
\newcounter{popsec}
\newcommand{\coursec}[1]{%
  \stepcounter{popsec}\setcounter{popsub}{0}%
  \par\vspace{16pt}\noindent
  \begin{minipage}{\linewidth}
  \tikz[baseline=-0.7ex]\node[draw=popink,line width=1.6pt,fill=poporange,rounded corners=4pt,minimum size=22pt,inner sep=0]{\popdisplay\fontsize{12}{12}\selectfont\color{popcream}\thepopsec};%
  \hspace{8pt}{\popdisplay\fontsize{15}{17}\selectfont\color{popink}\MakeUppercase{#1}}\par
  \vspace{4pt}\noindent{\color{poporange}\rule{36pt}{3.6pt}}
  \end{minipage}\par\nopagebreak\vspace{8pt}
}
\newcounter{popsub}
\newcommand{\courssub}[1]{%
  \stepcounter{popsub}%
  \par\vspace{9pt}\noindent{\popxbold\fontsize{11.5}{13}\selectfont\color{popink}{\color{poporange}\thepopsec.\thepopsub}\hspace{6pt}#1}\par\nopagebreak\vspace{4pt}
}

% ============================================================
% Boîtes pédagogiques — même recette : fond blanc, bordure épaisse colorée,
% ombre dure encre, badge plein. Titre optionnel [#1].
% ============================================================
\tcbset{popbase/.style={breakable,enhanced,colback=white,boxrule=2.3pt,arc=9pt,
  shadow={3pt}{-3pt}{0pt}{popink},left=14pt,right=14pt,top=11pt,bottom=11pt,before skip=11pt,after skip=15pt,
  fontupper=\popsemi\fontsize{10.6}{14.5}\selectfont\color{popink},
  fontlower=\popsemi\fontsize{10.6}{14.5}\selectfont\color{popink}}}
% Coche dessinée (le glyphe ✓ n'existe pas dans les polices du document)
\newcommand{\popcheck}[1]{\tikz[baseline=-0.5ex]\draw[#1,line width=1.6pt,line cap=round,line join=round](-0.13,0.0)--(-0.04,-0.09)--(0.13,0.1);}
\providecommand{\coche}{\popcheck{popgreen}}
\providecommand{\resultat}[1]{\par\smallskip\noindent\fcolorbox{popgreen}{popgreenL}{\parbox{\dimexpr\linewidth-2\fboxsep-2\fboxrule\relax}{\textbf{Résultat.} #1}}\par\smallskip}
\newcommand{\popboxhead}[3]{\poppill{#1}{#2}{white}\ifx\relax#3\relax\else\hspace{6pt}{\popxbold\fontsize{10.6}{12}\selectfont\color{popink}#3}\fi\par\vspace{7pt}}

\newtcolorbox{aretenir}[1][]{popbase,colframe=popgreen,before upper={\popboxhead{À retenir}{popgreen}{#1}}}
\newtcolorbox{exemplebox}[1][]{popbase,colframe=poporange,before upper={\popboxhead{Exemple}{poporange}{#1}}}
\newtcolorbox{definition}[1][]{popbase,colframe=popblue,before upper={\popboxhead{Définition}{popblue}{#1}}}
\newtcolorbox{propriete}[1][]{popbase,colframe=poppurple,before upper={\popboxhead{Propriété}{poppurple}{#1}}}
\newtcolorbox{methode}[1][]{popbase,colframe=popgold,before upper={\popboxhead{Méthode}{popgold}{#1}}}
\newtcolorbox{attention}[1][]{popbase,colframe=poppink,before upper={\popboxhead{Attention}{poppink}{#1}}}
\newtcolorbox{experience}[1][]{popbase,colframe=popblue,colback=popblueL,before upper={\popboxhead{Expérience}{popblue}{#1}}}
% « Le savais-tu ? » : carte violette pleine (contexte, culture, Cameroun)
\newtcolorbox{savaistu}[1][]{popbase,colback=poppurple,colframe=popink,
  fontupper=\popsemi\fontsize{10.4}{14.2}\selectfont\color{white},
  before upper={\poppill{Le savais-tu\,?}{white}{poppurple}\ifx\relax#1\relax\else\hspace{6pt}{\popxbold\color{white}#1}\fi\par\vspace{7pt}}}
% Exercice résolu : énoncé en haut, \tcblower puis la solution sur fond crème
\newcounter{popexo}\newcounter{popexr}
\newtcolorbox{exoresolu}[1][]{popbase,colframe=poporange,colbacklower=popcream,
  segmentation style={popink,line width=1.2pt,dashed},
  before upper={\stepcounter{popexr}\popboxhead{Exercice résolu \thepopexr}{poporange}{#1}},
  before lower={\poppill{Solution}{popink}{popcream}\par\vspace{6pt}}}
% Exercice (non résolu en place) — la correction est dans la partie Corrigés
\newtcolorbox{exercice}[1][]{popbase,colframe=popink,
  before upper={\stepcounter{popexo}\popboxhead{Exercice \thepopexo}{popink}{#1}}}
% Corrigé
\newtcolorbox{corrige}[1][]{popbase,colframe=popgreen,colback=popgreenL,
  before upper={\popboxhead{Corrigé}{popgreen}{#1}}}
% Objectifs / prérequis / bilan
\newtcolorbox{objectifs}{popbase,colframe=popink,colback=poporangeL,
  before upper={\popboxhead{Objectifs de la leçon}{popink}{}}}
\newtcolorbox{prerequis}{popbase,colframe=popink,colback=popblueL,
  before upper={\popboxhead{Prérequis}{popblue}{}}}
\newtcolorbox{bilan}{popbase,colframe=popink,colback=white,boxrule=2.8pt,
  before upper={\popboxhead{Fiche bilan}{poporange}{L'essentiel à connaître pour l'examen}}}
% Figure encadrée avec légende
\newtcolorbox{popfigure}[1][]{enhanced,breakable=false,colback=white,colframe=popline,boxrule=1.4pt,arc=8pt,
  left=10pt,right=10pt,top=10pt,bottom=8pt,before skip=10pt,after skip=12pt,halign=center,
  lower separated=false,halign lower=center,
  fontlower=\popsemi\fontsize{9}{11.5}\selectfont\color{popmuted},
  #1}
\newcommand{\legende}[1]{\par\vspace{6pt}{\popsemi\fontsize{9}{11.5}\selectfont\color{popmuted}#1}}

% ============================================================
% Signature OnBuch+
% ============================================================
\newcommand{\poplogoinline}[1]{{\poplogo\fontsize{#1}{#1}\selectfont\color{popink}OnBuch\color{poporange}+}}
\newcommand{\signatureonbuch}{%
  \begin{tcolorbox}[enhanced,colback=popink,colframe=popink,boxrule=2.2pt,arc=10pt,
      drop shadow=poporange,left=14pt,right=14pt,top=10pt,bottom=10pt,before skip=0pt,after skip=0pt]
    \tikz[baseline=-0.7ex]\node[circle,fill=popgreen,draw=popcream,line width=1.3pt,minimum size=22pt,inner sep=0]{\popcheck{white}};%
    \hspace{9pt}\begin{minipage}[c]{0.62\linewidth}
      {\popxbold\fontsize{12}{13}\selectfont\color{popcream}Signé\ \textbullet\ {\poplogo OnBuch\color{poporange}+}}\par\vspace{2pt}
      {\raggedright\popsemi\fontsize{8}{10.5}\selectfont\color{popline}\MakeUppercase{Équipe pédagogique OnBuch+}\\\MakeUppercase{Conforme au programme officiel MINESEC}\par}
    \end{minipage}\hfill
    {\poplogo\fontsize{22}{22}\selectfont\color{popcream}OnBuch\color{poporange}+}
  \end{tcolorbox}%
}

% ============================================================
% Page de garde
% #1 titre · #2 matière · #3 niveau · #4 module · #5 n° leçon · #6 durée ·
% #7 description / objectifs (texte ou itemize)
% ============================================================
\newcommand{\pagegarde}[7]{%
  \thispagestyle{empty}
  \newgeometry{a4paper,margin=1.7cm,top=1.6cm,bottom=1.4cm}%
  \begin{tikzpicture}[remember picture,overlay]
    \fill[poporange!18] ([xshift=-3.2cm,yshift=-3.6cm]current page.north east) circle (4.6cm);
    \fill[poppurple!14] ([xshift=2.4cm,yshift=7.6cm]current page.south west) circle (3.8cm);
  \end{tikzpicture}%
  {\poplogo\fontsize{30}{30}\selectfont\color{popink}OnBuch\color{poporange}+}\hfill
  \raisebox{0.6ex}{\poppill{Cours signé \textbullet\ Premium}{popink}{popcream}}\par
  \vspace{1pt}\popsquiggle{poppurple}\par
  \vspace{8pt}
  \begin{tcolorbox}[enhanced,colback=poporange,colframe=popink,boxrule=2.8pt,arc=13pt,
      drop shadow=popink,left=18pt,right=18pt,top=12pt,bottom=13pt,after skip=10pt]
    \poppill{#2 \textbullet\ #3}{popink}{popcream}\hspace{5pt}\poppill{#4}{white}{popink}\par\vspace{8pt}
    {\popdisplay\fontsize{14}{14}\selectfont\color{popink}LEÇON #5}\par\vspace{4pt}
    {\raggedright\hyphenpenalty=10000\exhyphenpenalty=10000\popdisplay\fontsize{25}{27}\selectfont\color{popcream}\MakeUppercase{#1}\par}
    \vspace{8pt}
    {\popxbold\fontsize{9.5}{9.5}\selectfont\color{popink}\MakeUppercase{Durée conseillée : #6}}
  \end{tcolorbox}
  \begin{tcolorbox}[enhanced,colback=white,colframe=popink,boxrule=2.2pt,arc=10pt,
      drop shadow=popink,left=16pt,right=16pt,top=10pt,bottom=10pt,after skip=8pt]
    \poppill{Dans cette leçon}{poppurple}{white}\par\vspace{5pt}
    \popsemi\fontsize{9.3}{12.6}\selectfont\color{popink}#7
  \end{tcolorbox}
  \noindent\begin{minipage}[t]{0.487\linewidth}
    \begin{tcolorbox}[enhanced,colback=popgreen,colframe=popink,boxrule=2.2pt,arc=10pt,
        drop shadow=popink,left=11pt,right=11pt,top=10pt,bottom=10pt,height=3.1cm,valign=center]
      \tcbox[colback=white,colframe=popink,boxrule=1.3pt,arc=5pt,boxsep=0pt,nobeforeafter,
        left=3pt,right=3pt,top=3pt,bottom=3pt,valign=center]{\qrcode[height=1.9cm]{https://play.google.com/store/apps/details?id=com.luvvix.onbuch&hl=fr}}%
      \hspace{8pt}\begin{minipage}[c]{0.5\linewidth}
        {\popdisplay\fontsize{11}{12.5}\selectfont\color{white}\MakeUppercase{Télécharge OnBuch}}\par\vspace{4pt}
        {\raggedright\popsemi\fontsize{8.4}{11}\selectfont\color{white}Gratuit \textbullet\ Google Play\\Tuteur IA, annales, concours, résultats\par}
      \end{minipage}
    \end{tcolorbox}
  \end{minipage}\hfill
  \begin{minipage}[t]{0.487\linewidth}
    \begin{tcolorbox}[enhanced,colback=poppurple,colframe=popink,boxrule=2.2pt,arc=10pt,
        drop shadow=popink,left=11pt,right=11pt,top=10pt,bottom=10pt,height=3.1cm,valign=center]
      \tikz[baseline=-0.7ex]\node[circle,fill=white,draw=popink,line width=1.3pt,minimum size=1.6cm,inner sep=0]{\popdisplay\fontsize{24}{24}\selectfont\color{poppurple}+};%
      \hspace{8pt}\begin{minipage}[c]{0.6\linewidth}
        {\popdisplay\fontsize{11}{12.5}\selectfont\color{white}\MakeUppercase{Passe à OnBuch+}}\par\vspace{4pt}
        {\raggedright\popsemi\fontsize{8.4}{11}\selectfont\color{white}Tous les cours signés, Tuteur IA illimité, fiches PDF — onglet OnBuch+ de l'app\par}
      \end{minipage}
    \end{tcolorbox}
  \end{minipage}
  \vspace{8pt}
  \popcontenu
  \vfill
  \signatureonbuch
  \restoregeometry
  \newpage
}

% Rangée « Ce cours contient » (page de garde)
\newcommand{\poptile}[3]{% #1 couleur #2 gros texte #3 légende
  \begin{tcolorbox}[enhanced,colback=#1,colframe=popink,boxrule=2pt,arc=9pt,shadow={2.5pt}{-2.5pt}{0pt}{popink},
      left=6pt,right=6pt,top=9pt,bottom=9pt,halign=center,valign=center,height=1.75cm,width=\linewidth,nobeforeafter]
    {\popdisplay\fontsize{12.5}{13}\selectfont\color{white}\MakeUppercase{#2}}\par\vspace{3pt}
    {\centering\popsemi\fontsize{7.8}{9.5}\selectfont\color{white}#3\par}
  \end{tcolorbox}}
\newcommand{\popcontenu}{%
  \par\noindent{\popxboldls\fontsize{7.5}{7.5}\selectfont\color{popmuted}CE COURS SIGNÉ CONTIENT}\par\vspace{7pt}
  \noindent\begin{minipage}[t]{0.235\linewidth}\poptile{popblue}{Cours}{complet, illustré, pas à pas}\end{minipage}\hfill
  \begin{minipage}[t]{0.235\linewidth}\poptile{poporange}{Méthodes}{et exercices résolus}\end{minipage}\hfill
  \begin{minipage}[t]{0.235\linewidth}\poptile{poppink}{Exercices}{type examen + corrigés}\end{minipage}\hfill
  \begin{minipage}[t]{0.235\linewidth}\poptile{popgreen}{Fiche bilan}{l'essentiel à retenir}\end{minipage}\par}

% Sommaire
\newtcolorbox{sommaire}{enhanced,colback=white,colframe=popink,boxrule=2.2pt,arc=10pt,
  drop shadow=popink,left=18pt,right=18pt,top=13pt,bottom=13pt,before skip=6pt,after skip=20pt,
  before upper={\poppill{Sommaire}{poppurple}{white}\par\vspace{9pt}\popsemi\fontsize{10.6}{14.5}\selectfont\color{popink}}}

% Signature de fin de cours — poussée tout en bas de la dernière page
\newcommand{\findecours}{%
  \par\vfill\nopagebreak
  \begin{center}\popsquiggle{poporange}\end{center}\vspace{4pt}
  \signatureonbuch
}
\AtBeginDocument{\def\llbracket{\mathopen{[\mkern-3.5mu[}}\def\rrbracket{\mathclose{]\mkern-3.5mu]}}} % intervalles d'entiers (glyphe ⟦ absent de la police)
% Glyphes absents de Fira Math : redéfinis à partir de glyphes disponibles
\AtBeginDocument{%
  \def\setminus{\mathbin{\backslash}}\let\smallsetminus\setminus
  \def\vdots{\mathord{\vbox{\baselineskip3.2pt\lineskiplimit0pt\kern4pt\hbox{.}\hbox{.}\hbox{.}}}}%
  \def\ddots{\mathinner{\mkern1mu\raise6.5pt\hbox{.}\mkern2mu\raise3.6pt\hbox{.}\mkern2mu\raise0.7pt\hbox{.}\mkern1mu}}%
  \def\triangle{\mathord{\scalebox{0.9}{$\symup{\Delta}$}}}%
  \def\nabla{\mathord{\raisebox{\depth}{\scalebox{1}[-1]{$\symup{\Delta}$}}}}%
  \def\checkmark{\popcheck{popdark}}%
  \def\star{\mathbin{\ast}}\def\bigstar{\mathord{\ast}}%
  \def\diamond{\mathbin{\scalebox{0.6}{$\Diamond$}}}%
  \def\bigcup{\mathop{\mathchoice{\raisebox{-0.3em}{\scalebox{1.7}{$\cup$}}}{\scalebox{1.25}{$\cup$}}{\cup}{\cup}}\displaylimits}%
  \def\bigcap{\mathop{\mathchoice{\raisebox{-0.3em}{\scalebox{1.7}{$\cap$}}}{\scalebox{1.25}{$\cap$}}{\cap}{\cap}}\displaylimits}%
  \def\lesseqqgtr{\mathrel{\lessgtr}}\def\bigoplus{\mathop{\oplus}\displaylimits}%
}
\providecommand{\ssi}{si et seulement si} % abréviation fréquente
\AtBeginDocument{\providecommand{\wideparen}{\overparen}} % arc de cercle

%% FIGFIX-BEGIN v4 — correctifs globaux des figures (docs/onbuch-generation/tools/figfix)
\usepackage{adjustbox}\usepackage{etoolbox}
% toute figure TikZ/pgfplots plus large que la ligne est réduite à la largeur disponible
\BeforeBeginEnvironment{tikzpicture}{\begin{adjustbox}{max width=\linewidth}}
\AfterEndEnvironment{tikzpicture}{\end{adjustbox}}
% légendes pgfplots lisibles par-dessus les courbes
\pgfplotsset{every axis/.append style={legend style={fill=white,fill opacity=0.92,text opacity=1,draw=black!15,inner sep=3pt}}}
% étiquettes d'arêtes (syntaxe edge["..."]) avec fond blanc
\tikzset{every edge quotes/.append style={fill=white,inner sep=1.2pt}}
%% FIGFIX-END
\begin{document}

\pagegarde{Les révolutions scientifiques et techniques depuis la fin de la Deuxième Guerre mondiale}{Histoire}{Tle Tech}{Histoire – classe de Terminale technique}{N05}{2 séances (fiche de progression harmonisée)}{Cette leçon étudie les grandes découvertes scientifiques (atomique, espace, biologie, médecine) et les grandes innovations techniques (informatique, transports, télécommunications, énergie) qui ont transformé le monde depuis 1945. Elle montre leurs effets concrets sur la vie quotidienne, y compris au Cameroun, et leurs limites.
\vspace{4pt}
\begin{multicols}{2}\small
\begin{itemize}
\item À la fin de cette leçon, je sais définir les notions de révolution scientifique et de révolution technique.
\item À la fin de cette leçon, je sais citer et dater les grandes découvertes scientifiques depuis 1945 (atome, ADN, espace, médecine).
\item À la fin de cette leçon, je sais identifier les grandes innovations techniques (informatique, internet, transports, énergie) et leurs domaines d'application.
\item À la fin de cette leçon, je sais expliquer les causes de l'accélération scientifique et technique après 1945.
\item À la fin de cette leçon, je sais analyser des documents (photographies, données chiffrées, cartes) pour établir une notion historique.
\item À la fin de cette leçon, je sais construire une frise chronologique et un diagramme à partir de données sur les innovations.
\end{itemize}\end{multicols}}

\begin{sommaire}
\begin{multicols}{2}
\begin{enumerate}
\item Introduction : notions et contexte de l'accélération scientifique (1945-...)
\item Les grandes découvertes scientifiques : maîtrise de l'atome et conquête de l'espace
\item Les grandes découvertes scientifiques : révolution de la biologie et de la médecine
\item Les grandes innovations techniques : informatique, internet et télécommunications
\item Les grandes innovations techniques : transports, énergie et agriculture
\item Bilan : les effets des RST sur le monde et le Cameroun
\item Activité d'intégration
\item Méthodes et exercices résolus
\item Exercices
\item Corrigés des exercices
\item Fiche bilan
\end{enumerate}
\end{multicols}
\end{sommaire}

\begin{objectifs}
\begin{itemize}
\item À la fin de cette leçon, je sais définir les notions de révolution scientifique et de révolution technique.
\item À la fin de cette leçon, je sais citer et dater les grandes découvertes scientifiques depuis 1945 (atome, ADN, espace, médecine).
\item À la fin de cette leçon, je sais identifier les grandes innovations techniques (informatique, internet, transports, énergie) et leurs domaines d'application.
\item À la fin de cette leçon, je sais expliquer les causes de l'accélération scientifique et technique après 1945.
\item À la fin de cette leçon, je sais analyser des documents (photographies, données chiffrées, cartes) pour établir une notion historique.
\item À la fin de cette leçon, je sais construire une frise chronologique et un diagramme à partir de données sur les innovations.
\item À la fin de cette leçon, je sais appliquer ces notions à des situations concrètes de la vie courante au Cameroun (téléphonie mobile, Mobile Money, santé).
\item À la fin de cette leçon, je sais rédiger et justifier une conclusion argumentée sur les bienfaits et les limites du progrès scientifique et technique.
\end{itemize}
\end{objectifs}

\begin{prerequis}
\begin{itemize}
\item La Seconde Guerre mondiale et ses conséquences (étudiée en classe de Première)
\item La course aux armements et la Guerre froide (rivalité USA-URSS)
\item La révolution industrielle et les premières inventions (vapeur, électricité) vues en 4e-3e
\item La décolonisation et l'entrée du Cameroun dans le monde moderne (1960)
\item Notions de base : innovation, progrès, mondialisation
\end{itemize}
\end{prerequis}

\newpage

\coursec{Introduction : notions et contexte de l'accélération scientifique (1945-...)}

\courssub{Situation d'accroche et problématique}

En 2024, un lycéen de Yaoundé commande un téléphone fabriqué en Chine, paie par \cle{Mobile Money}, suit la livraison par \cle{GPS} et consulte ses résultats scolaires en ligne. Pourtant, son grand-père, en 1950, n'avait jamais vu ni télévision ni ordinateur. Comment expliquer qu'en moins de 80 ans, le monde — et le Cameroun — ait connu une telle transformation ? Quelles découvertes scientifiques et quelles innovations techniques, nées après 1945, ont rendu cela possible ? C'est toute l'histoire des \cle{révolutions scientifiques et techniques} (RST) depuis la fin de la Deuxième Guerre mondiale.

\courssub{Notions clés : découverte, innovation, révolution scientifique et technique}

\begin{definition}[Découverte scientifique vs Innovation technique]
On distingue deux concepts fondamentaux :
\begin{itemize}
    \item La \textbf{découverte scientifique} : c'est l'accès à une connaissance nouvelle sur la nature, sans application immédiate prévue. \emph{Exemple :} la détermination de la structure en double hélice de l'\cle{ADN} par Watson, Crick et Franklin (1953).
    \item L'\textbf{innovation technique} : c'est la mise en œuvre concrète et industrialisée d'une découverte (ou d'une combinaison de savoirs) pour répondre à un besoin. \emph{Exemple :} les tests de paternité ou le dépistage génétique, applications directes de la connaissance de l'ADN.
\end{itemize}
La \cle{recherche fondamentale} vise la découverte ; la \cle{recherche appliquée} et le \cle{développement} (R\&D) visent l'innovation.
\end{definition}

\begin{definition}[Révolution scientifique et technique (RST)]
Une \cle{révolution scientifique et technique} est un ensemble de découvertes et d'innovations qui transforment rapidement et profondément la production, les transports, la communication et la vie quotidienne. Contrairement aux révolutions industrielles classiques (vapeur, électricité), la RST post-1945 se caractérise par :
\begin{itemize}
    \item l'interpénétration science/technique (la science précède et guide la technique) ;
    \item l'accélération des cycles (quelques années entre la découverte et le marché) ;
    \item la portée systémique (économie, société, géopolitique, environnement).
\end{itemize}
\end{definition}

\begin{attention}[Piège fréquent : confusion des révolutions industrielles]
Ne pas confondre la \textbf{3e révolution industrielle} (numérique, électronique, nucléaire, après 1945) avec les deux premières :
\begin{itemize}
    \item 1re (fin XVIIIe–XIXe) : machine à vapeur, charbon, textile, chemin de fer.
    \item 2e (fin XIXe–1939) : électricité, pétrole, chimie, automobile, sidérurgie.
    \item 3e (depuis 1945) : nucléaire, électronique, informatique, biotechnologies, spatial, télécommunications.
\end{itemize}
Au Baccalauréat technique, précisez toujours la période et les secteurs moteurs.
\end{attention}

\courssub{Le point de départ : les acquis de la Deuxième Guerre mondiale (1939–1945)}

La guerre totale a agi comme un formidable accélérateur. Les États-Unis, le Royaume-Uni, l'URSS et l'Allemagne nazie ont mobilisé des moyens financiers, humains et industriels sans précédent pour la recherche militaire. Quatre avancées majeures structurent l'après-guerre :

\begin{itemize}
    \item \textbf{Le radar et l'électronique :} la détection radioélectrique (radar) impose la maîtrise des hyperfréquences et des tubes électroniques, base de l'électronique grand public et des télécommunications.
    \item \textbf{Les fusées V2 (Allemagne) :} premier engin balistique opérationnel (Wernher von Braun). Elle ouvre la voie aux lanceurs spatiaux et aux missiles intercontinentaux (ICBM).
    \item \textbf{La bombe atomique (Projet Manhattan, USA) :} maîtrise de la fission nucléaire (Hiroshima/Nagasaki, août 1945). Point de départ de l'énergie nucléaire civile et de la dissuasion stratégique.
    \item \textbf{La pénicilline (production de masse, USA/UK) :} premier antibiotique industrialisé, symbole de la révolution pharmaceutique.
    \item \textbf{Les premiers calculateurs électroniques :} \cle{ENIAC} (USA, 1946), \cle{Colossus} (UK, 1943, secret), \cle{Z3} (Allemagne, 1941). Naissance de l'informatique.
\end{itemize}

\begin{savaistu}[Hiroshima, 6 août 1945 : tragédie et basculement]
La bombe \emph{Little Boy} (uranium) sur Hiroshima, puis \emph{Fat Man} (plutonium) sur Nagasaki (9 août), font 200 000 morts immédiats et différés. Au-delà de l'horreur, cet événement marque l'entrée dans l'\cle{ère nucléaire} : course aux armements (URSS 1949, UK 1952, France 1960, Chine 1964), doctrine de la \cle{dissuasion nucléaire}, mais aussi espoirs de l'énergie « trop bon marché pour être mesurée » (centrales civiles dès 1954 à Obninsk, URSS).
\end{savaistu}

\begin{popfigure}
\begin{tikzpicture}[
    timeline/.style={very thick, popblue},
    event/.style={circle, fill=popblue, inner sep=2pt},
    eventlabel/.style={anchor=west, font=\footnotesize, text width=3.5cm, align=left},
    major/.style={rectangle, fill=poporange, rounded corners=2pt, inner sep=3pt, font=\footnotesize\bfseries, text width=4cm, align=center},
    yearmark/.style={font=\small\bfseries, popdark}
]
% Axe principal
\draw[timeline] (0,0) -- (16,0);
% Années
\foreach \x/\year in {0/1945, 2/1955, 4/1965, 6/1975, 8/1985, 10/1995, 12/2005, 14/2015, 16/2025} {
    \draw (\x,0.15) -- (\x,-0.15) node[below, yearmark] {\year};
}
% Événements majeurs (année approximative -> position x = (année-1945)/5)
% 1945: Fin WWII, Bombe A, ENIAC
\node[event] at (0,0.6) {};
\node[eventlabel] at (0.3,0.6) {Fin WWII, Bombe A, ENIAC};
% 1947: Transistor
\node[event] at (0.4, -0.7) {};
\node[eventlabel] at (0.7,-0.7) {Transistor (Bell Labs)};
% 1953: ADN
\node[event] at (1.6, 0.9) {};
\node[eventlabel] at (1.9,0.9) {Structure ADN};
% 1957: Spoutnik
\node[event] at (2.4, -1.0) {};
\node[eventlabel] at (2.7,-1.0) {Spoutnik 1 (URSS)};
% 1958: NASA
\node[event] at (2.6, 0.6) {};
\node[eventlabel] at (2.9,0.6) {Création NASA};
% 1969: Apollo 11
\node[event] at (4.8, -0.8) {};
\node[eventlabel] at (5.1,-0.8) {Apollo 11 (Lune)};
% 1971: Microprocesseur Intel 4004
\node[event] at (5.2, 0.9) {};
\node[eventlabel] at (5.5,0.9) {Microprocesseur 4004};
% 1983: TCP/IP (Internet)
\node[event] at (7.6, -0.7) {};
\node[eventlabel] at (7.9,-0.7) {TCP/IP standard};
% 1990: Web (Berners-Lee)
\node[event] at (9.0, 0.8) {};
\node[eventlabel] at (9.3,0.8) {World Wide Web};
% 2001: Séquençage génome humain
\node[event] at (11.2, -0.9) {};
\node[eventlabel] at (11.5,-0.9) {Génome humain};
% 2007: iPhone / Smartphone
\node[event] at (12.4, 0.7) {};
\node[eventlabel] at (12.7,0.7) {iPhone / Smartphones};
% 2020: Vaccins ARN messager (Covid)
\node[event] at (15.0, -0.8) {};
\node[eventlabel] at (15.3,-0.8) {Vaccins ARNm};
% Légende
\node[major, anchor=south west] at (0, -2.2) {Frise chronologique : jalons majeurs des RST (1945–2025)};
\end{tikzpicture}
\legende{Frise chronologique simplifiée des grandes découvertes et innovations depuis 1945. Les positions horizontales sont proportionnelles aux années (échelle ~5 ans/cm).}
\end{popfigure}

\courssub{Les moteurs de l'accélération : Guerre froide, reconstruction, État stratège}

\begin{definition}[Facteurs structurels de l'accélération post-1945]
Trois moteurs principaux expliquent la vitesse sans précédent des RST :
\begin{enumerate}
    \item \textbf{La rivalité Est–Ouest (Guerre froide, 1947–1991) :} la confrontation idéologique et militaire entre le bloc occidental (OTAN, mené par les USA) et le bloc de l'Est (Pacte de Varsovie, mené par l'URSS) transforme la science en enjeu de puissance. Deux courses emblématiques :
    \begin{itemize}
        \item \cle{Course à l'espace} : Spoutnik (1957), Gagarine (1961), Apollo (1969). Retombées : télécoms par satellite, matériaux, informatique embarquée, GPS.
        \item \cle{Course aux armements nucléaires} : bombe H (1952/1953), missiles balistiques (ICBM), dissuasion. Financement massif de la physique des hautes énergies, du calcul numérique, de l'aéronautique.
    \end{itemize}
    \item \textbf{Les besoins de la reconstruction et de la croissance (les « Trente Glorieuses », 1945–1973) :} l'Europe et le Japon dévastés doivent reconstruire infrastructures, logement, industrie. La demande massive en énergie (pétrole, puis nucléaire), en chimie (plastiques, engrais), en transports (autoroutes, aviation civile) tire l'innovation.
    \item \textbf{L'État stratège (planification, financement, commande publique) :} aux USA (\cle{DARPA} 1958, budgets défense), en URSS (Académies des sciences, \cle{Gosplan}), en France (Commissariat à l'énergie atomique \cle{CEA} 1945, Plan Calcul 1966), au Japon (MITI), l'État oriente, finance les laboratoires, passe commande (ex. : internet né d'ARPANET, projet militaire US).
\end{enumerate}
\end{definition}

\begin{popfigure}
\begin{tikzpicture}[
    mindmap,
    concept color=popblue,
    level 1/.append style={level distance=3.5cm, sibling angle=90, concept color=popblue!70},
    level 2/.append style={level distance=2.8cm, sibling angle=45, concept color=popgreen!70},
    level 3/.append style={level distance=2.2cm, sibling angle=30, concept color=poporange!70},
    every node/.style={font=\small},
    edge from parent/.style={draw, thick, popline}
]
\node[concept, align=center]{Facteurs de l'accélération\\scientifique (1945-...)}
    child[concept color=poppurple] {node[concept, align=center]{Guerre froide\\Rivalité Est-Ouest}
        child {node[concept, align=center]{Course à l'espace\\Spoutnik, Apollo, Satellites}}
        child {node[concept, align=center]{Course aux armements\\Bombe H, ICBM, DARPA}}
        child {node[concept, align=center]{Physique hautes énergies\\CERN (1954)}}
    }
    child[concept color=popgreen] {node[concept, align=center]{Reconstruction\\Croissance (30 Glorieuses)}
        child {node[concept, align=center]{Énergie : Pétrole,\\Nucléaire civil}}
        child {node[concept, align=center]{Chimie : Plastiques,\\Engrais, Pharmaco}}
        child {node[concept, align=center]{Transports : Auto,\\Aviation civile, TGV}}
        child {node[concept, align=center]{Agriculture : Révolution\\verte (semences, irrigation)}}
    }
    child[concept color=poporange] {node[concept, align=center]{État stratège\\Planification \& Financement}
        child {node[concept, align=center]{USA : DARPA, NASA,\\Budget Défense}}
        child {node[concept, align=center]{France : CEA, Plan Calcul,\\CNES (1961)}}
        child {node[concept, align=center]{Japon : MITI,\\Keiretsu R\&D}}
        child {node[concept, align=center]{Chine : 863 Program\\ (1986), CNSA}}
    }
    child[concept color=poppink] {node[concept, align=center]{Mondialisation\\Diffusion Sud}
        child {node[concept, align=center]{Transfert tech.\\(licences, IDE)}}
        child {node[concept, align=center]{Chaînes de valeur\\mondiales}}
        child {node[concept, align=center]{Cameroun : Mobile Money,\\Télécoms, Santé num.}}
    };
\end{tikzpicture}
\legende{Carte mentale des quatre moteurs structurels de l'accélération scientifique et technique depuis 1945.}
\end{popfigure}

\courssub{Les grandes puissances et les lieux de la recherche}

\begin{itemize}
    \item \textbf{États-Unis :} hégémonie scientifique (budget R\&D \textasciitilde 3\% du PIB), universités d'élite (MIT, Stanford, Caltech), laboratoires fédéraux (Los Alamos, Bell Labs), agences (NASA 1958, DARPA 1958, NIH). Modèle « triple hélice » État–Université–Entreprise (Silicon Valley).
    \item \textbf{URSS / Russie :} système académique centralisé (Académie des sciences), villes closes (Sarov, Dubna), exploits spatiaux et nucléaires, mais diffusion civile limitée. Héritage : excellence en maths/physique, fragilité du transfert technologique.
    \item \textbf{Europe :} reconstruction via \cle{CECA} (1951) puis \cle{CEE} (1957). Création du \cle{CERN} (1954, Genève) pour la physique des particules (LHC, boson de Higgs 2012), \cle{ESA} (1975, spatial), \cle{Airbus} (1970, aéronautique), \cle{ESRF} (synchrotron). Programme-cadre UE pour la R\&D collaborative.
    \item \textbf{Japon (dès 1950s) :} rattrapage par imitation-amélioration (électronique grand public, automobile, robotique), coordination MITI/keiretsu, effort R\&D privé massif (Sony, Toyota, Hitachi).
    \item \textbf{Chine (depuis 1978/1986) :} « Programme 863 » (haute technologie), retour des cerveaux, investissements massifs (R\&D > 2,4\% PIB), leadership actuel : 5G (Huawei), spatial (CNSA, station Tiangong), IA, quantique, biotechnologies.
\end{itemize}

\begin{aretenir}[Lieux emblématiques de la recherche post-1945]
\begin{center}
\begin{tabularx}{\linewidth}{|l|X|X|}\hline
\rowcolor{popblueL}
\textbf{Structure} & \textbf{Création} & \textbf{Rôle majeur} \\\hline
CERN (Genève) & 1954 & Physique des particules (LHC, Web inventé ici par Berners-Lee 1990) \\\hline
NASA (USA) & 1958 & Programme spatial civil (Apollo, Navette, ISS, Mars, Artemis) \\\hline
DARPA (USA) & 1958 & Ruptures militaires (ARPANET$\to$Internet, GPS, IA, drones) \\\hline
CEA (France) & 1945 & Énergie nucléaire (civile/militaire), recherche fondamentale \\\hline
CNES (France) & 1961 & Spatial (Ariane, Kourou, sciences de l'univers) \\\hline
Bell Labs (USA) & 1925 (pic 1947-80) & Transistor (1947), Laser (1960), Unix, C, cellule solaire \\\hline
\end{tabularx}
\end{center}
\end{aretenir}

\courssub{Mondialisation et diffusion vers les pays du Sud : le cas du Cameroun}

La diffusion des innovations suit une géographie inégale : brevetage (OCDE), licences, investissements directs étrangers (IDE), chaînes de valeur mondiales. Le \cle{transfert de technologie} n'est pas automatique ; il suppose capacités d'absorption (éducation, infrastructures), cadre institutionnel et demande locale.

Au \textbf{Cameroun}, la réception des RST se lit à plusieurs échelles :
\begin{itemize}
    \item \textbf{Télécommunications :} saut direct au mobile (années 2000) sans passer par le fixe filaire massif. Taux de pénétration des abonnements mobiles > 90\% (2023). \cle{Mobile Money} (MTN MoMo, Orange Money) : inclusion financière, paiement scolaire, épargne.
    \item \textbf{Santé :} vaccins (PEV), biotechnologies (diagnostic PCR, thérapies ARV), télémédecine naissante.
    \item \textbf{Agriculture :} variétés améliorées (IRAD/IITA : maïs, manioc, plantain), engrais, mécanisation légère, information marchés par SMS.
    \item \textbf{Éducation/Administration :} numérisation concours (MINFOPRA), résultats en ligne, plateformes d'e-learning (COVID-19), fibre optique (backbone national géré par Camtel).
    \item \textbf{Énergie :} barrages (Lom Pangar, Nachtigal), solaire rural, enjeu nucléaire civil (études AIEA).
\end{itemize}

\begin{exemplebox}[Le « saut de grenouille » (leapfrogging) camerounais]
Le Cameroun n'a pas déployé massivement le téléphone fixe cuivre (coût, entretien). L'arrivée de la 2G/3G/4G a permis une couverture rapide du territoire. De même, le Mobile Money a supplanté le réseau bancaire classique pour les populations non bancarisées. C'est un exemple typique de \cle{diffusion différentielle} : le Sud adopte les technologies matures, standards, peu coûteuses, sans reproduire l'historique technique du Nord.
\end{exemplebox}

\begin{methode}[Analyser l'impact d'une innovation au Cameroun (démarche Baccalauréat)]
\begin{enumerate}
    \item \textbf{Identifier l'innovation} (ex. : smartphone 4G, Mobile Money, variété de manioc résistante).
    \item \textbf{Dater l'introduction} et situer l'origine (importation, transfert, adaptation locale).
    \item \textbf{Décrire les usages} (acteurs : jeunes, paysans, commerçants, administration).
    \item \textbf{Évaluer les effets} : économiques (revenus, productivité), sociaux (lien, inclusion), spatiaux (désenclavement), environnementaux (déchets électroniques, consommation énergie).
    \item \textbf{Nuancer} : fracture numérique (urbain/rural, genre, revenu), dépendance technique, souveraineté des données.
\end{enumerate}
\end{methode}

\begin{exoresolu}[Exemple : l'introduction du Mobile Money au Cameroun]
\textbf{Énoncé :} En vous appuyant sur vos connaissances, montrez comment le Mobile Money illustre la diffusion d'une innovation technique post-1945 dans un pays du Sud.

\textbf{Solution détaillée :}
\begin{enumerate}
    \item \textbf{Origine scientifique/technique :} cryptographie à clé publique (années 1970, Diffie-Hellman, RSA), réseaux cellulaires GSM (norme européenne 1987, 2G), protocoles USSD/SMS, bases de données distribuées. Innovation \emph{financière} (modèle M-Pesa Kenya 2007) adaptée au contexte camerounais (MTN 2010, Orange 2011).
    \item \textbf{Diffusion :} adoption exponentielle (proximité agents, pas de compte bancaire requis, interopérabilité progressive). 10 millions de comptes actifs (2023).
    \item \textbf{Effets structurels :} inclusion financière (femmes, ruraux), réduction coûts de transaction, formalisation informel, levier pour e-commerce, paiement salaires/bourse élèves.
    \item \textbf{Limites :} dépendance opérateurs privés, cybersécurité (fraude SIM swap), fiscalité (prélèvement 0,2\%), exclusion des zones sans couverture (zones blanches).
\end{enumerate}
\emph{Conclusion :} le Mobile Money est une innovation de \cle{recombinaison} (technos existantes + modèle d'affaires) qui transforme en profondeur l'économie camerounaise, illustrant la mondialisation asymétrique des RST.
\end{exoresolu}

\courssub{Bilan de l'introduction}

Depuis 1945, la conjugaison unique de \textbf{besoins militaires} (Guerre froide), de \textbf{demande civile massive} (reconstruction, croissance) et d'un \textbf{Etat investisseur} a produit une série de ruptures — nucléaire, électronique, informatique, spatiale, biologique — qui constituent la \cle{3e révolution industrielle}. Leur diffusion, inégale mais accélérée par la mondialisation, atteint aujourd'hui le Cameroun, où elles recomposent les rapports au travail, à l'argent, à la santé et au savoir. Les sections suivantes détailleront ces révolutions sectorielles : atome et espace (2), biologie-médecine (3), numérique-télécoms (4), transports-énergie-agriculture (5), avant un bilan global (6).

\coursec{Les grandes découvertes scientifiques : maîtrise de l'atome et conquête de l'espace}

Après avoir posé le cadre général de l'accélération scientifique d'après-guerre — contexte géopolitique de la Guerre froide, investissements massifs des États, structuration de la recherche en « \cle{Big Science} » —, nous entrons dans le cœur des deux chantiers qui ont le plus profondément redessiné le rapport de l'humanité à l'énergie et à l'espace : la maîtrise de l'atome et la conquête spatiale. Ces deux aventures, nées dans les laboratoires de physique fondamentale, basculent très vite dans l'ingénierie à grande échelle, avec des retombées civiles majeures et des enjeux stratégiques planétaires.

\courssub{L'énergie nucléaire : de la fission à la première centrale civile}

\textbf{1. L'intuition physique : la fission, source d'énergie colossale.}
Rappelons le principe découvert en 1938 par Otto Hahn et Fritz Strassmann (expérience), interprété par Lise Meitner et Otto Frisch : le noyau d'un isotope lourd (uranium-235, plutonium-239), bombardé par un neutron, se scinde en deux noyaux plus légers, libérant \emph{de nouveaux neutrons} et une énergie considérable (la perte de masse $\Delta m$ se convertit en énergie selon $E = \Delta m \, c^2$). Chaque fission libère environ 200 MeV ; la \cle{réaction en chaîne} — chaque fission déclenchant plusieurs nouvelles fissions — permet une libération d'énergie considérable, contrôlée (réacteur) ou explosive (bombe).

\begin{definition}[Réaction en chaîne contrôlée]
Dans un \cle{réacteur nucléaire}, la réaction de fission est maintenue à l'état critique : chaque fission entraîne en moyenne une seule fission suivante. Le contrôle s'obtient par des \cle{barres de contrôle} (matériaux absorbant les neutrons : bore, hafnium, cadmium) et un \cle{modérateur} (eau légère, eau lourde ou graphite) qui ralentit les neutrons pour augmenter leur probabilité de provoquer de nouvelles fissions de l'uranium-235.
\end{definition}

\textbf{2. Du projet Manhattan à Obninsk (1954) : la bascule civile.}
La Seconde Guerre mondiale accélère la maîtrise technique (projet Manhattan, réacteurs de Hanford produisant le plutonium). Dès 1945, des physiciens (Szilard, Fermi) envisagent l'usage civil : produire de la chaleur pour transformer de l'eau en vapeur, actionner une turbine et générer de l'électricité.
\begin{itemize}
    \item \textbf{1951, EBR-I (Idaho, États-Unis) :} première production d'électricité d'origine nucléaire (quatre ampoules allumées par un réacteur expérimental à neutrons rapides).
    \item \textbf{27 juin 1954, Obninsk (URSS) :} mise en service de la \cle{première centrale nucléaire raccordée à un réseau électrique} (réacteur de 5 MW électriques, modéré au graphite, refroidi à l'eau). Elle alimente la région de Moscou.
    \item \textbf{1956, Calder Hall (Royaume-Uni) :} première centrale de puissance à vocation commerciale (filière Magnox, environ 50 MW par réacteur), produisant aussi du plutonium à usage militaire.
    \item \textbf{1957, Shippingport (États-Unis) :} première centrale américaine entièrement civile (réacteur à eau pressurisée, environ 60 MW).
\end{itemize}

\begin{popfigure}
\begin{tikzpicture}[scale=0.75, every node/.style={font=\small}]
    % REACTEUR (enceinte de confinement)
    \draw[thick, fill=popcream] (0,0) rectangle (6,6);
    \node[align=center] at (3,6.4) {\textbf{Enceinte de confinement}};
    
    % Coeur (cuve)
    \draw[thick, fill=poppinkL] (1.5,1.5) rectangle (4.5,4.5);
    \node[poppink, align=center] at (3,3) {\textbf{CŒUR}\\Combustible (UO$_2$)\\Modérateur / Caloporteur};
    % Barres de controle
    \foreach \y in {1.8, 2.4, 3.0, 3.6, 4.2} {
        \draw[thick, fill=popdark] (1.6,\y) rectangle (1.9,\y+0.15);
        \draw[thick, fill=popdark] (4.1,\y) rectangle (4.4,\y+0.15);
    }
    \node[popdark, rotate=90] at (1.2,3) {Barres de contrôle};
    
    % Circuit primaire (eau sous pression)
    \draw[popblue, very thick, ->] (4.5,3) -- (7,3) node[midway, above] {Eau primaire chaude};
    \draw[popblue, thick] (7,3) -- (7,4);
    
    % Generateur de vapeur
    \draw[thick, fill=popblueL] (7,4) rectangle (8.5,5.5);
    \node[align=center] at (7.75,4.75) {Générateur\\de vapeur};
    % Fleche vapeur secondaire
    \draw[poporange, very thick, ->] (8.5,4.75) -- (11,4.75) node[midway, above] {Vapeur haute pression};
    
    % Turbine
    \draw[thick, fill=popmuted!30] (11,3.5) rectangle (13.5,6);
    \node[popdark, align=center] at (12.25,4.75) {\textbf{TURBINE}};
    \draw[->, thick] (12.25,3.5) -- (12.25,2.5);
    
    % Alternateur
    \draw[thick, fill=popgoldL] (11,1.5) rectangle (13.5,3.5);
    \node[popdark, align=center] at (12.25,2.5) {\textbf{ALTERNATEUR}};
    \draw[popgold, very thick, ->] (13.5,2.5) -- (15,2.5) node[right] {Électricité};
    
    % Condenseur / Circuit tertiaire
    \draw[poporange, thick, ->] (12.25,3.5) -- (12.25,1.5) node[midway, right] {Vapeur basse pression};
    \draw[thick, fill=popblueL!50] (10,0) rectangle (14.5,1.5);
    \node[popblue, align=center] at (12.25,0.75) {CONDENSEUR};
    \draw[popblue, thick, ->] (10,0.75) -- (7,0.75) node[midway, below] {Eau de refroidissement};
    \draw[popblue, thick, ->] (7,0.75) -- (7,1.5) -- (7,3);
    
    % Legende circuits
    \node[align=left, anchor=north west] at (0.2, -0.5) {
        \textcolor{popblue}{\textbf{Circuit primaire :}} eau pressurisée, traverse le cœur, cède sa chaleur au générateur de vapeur.\\
        \textcolor{poporange}{\textbf{Circuit secondaire :}} eau transformée en vapeur, actionne la turbine couplée à l'alternateur.\\
        \textcolor{popblue!60}{\textbf{Circuit tertiaire :}} refroidit le condenseur (tour aéroréfrigérante, rivière ou mer).
    };
\end{tikzpicture}
\legende{Schéma simplifié d'une centrale nucléaire à eau pressurisée (REP), filière la plus répandue dans le monde. Le cœur chauffe l'eau du circuit primaire ; celle-ci vaporise l'eau du circuit secondaire dans le générateur de vapeur ; la vapeur entraîne la turbine couplée à l'alternateur qui produit l'électricité.}
\end{popfigure}

\begin{aretenir}[Les trois circuits d'une centrale REP]
\begin{enumerate}
    \item \textbf{Primaire (radioactif) :} eau maintenue sous forte pression (environ 155 bars) pour rester liquide vers 320 °C ; elle transporte la chaleur du cœur.
    \item \textbf{Secondaire (normalement non radioactif) :} eau $\to$ vapeur $\to$ turbine $\to$ alternateur $\to$ électricité.
    \item \textbf{Tertiaire (environnement) :} refroidissement du condenseur (tour aéroréfrigérante ou prélèvement en mer/rivière).
\end{enumerate}
La séparation stricte entre primaire et secondaire est une barrière majeure de sûreté (confinement de la radioactivité).
\end{aretenir}

\textbf{3. Applications de l'atome : électricité, médecine, armement.}
\begin{itemize}
    \item \textbf{Électricité :} au début des années 2020, environ 440 réacteurs en service dans une trentaine de pays fournissent environ 10 \% de l'électricité mondiale (environ 25 \% dans l'Union européenne, plus de 60 \% en France). Avantages : production massive, faible émission de CO$_2$. Problèmes : gestion des déchets radioactifs à vie longue et risque d'accident majeur (Three Mile Island 1979, Tchernobyl 1986, Fukushima 2011).
    \item \textbf{Médecine nucléaire :}
        \begin{itemize}
            \item \emph{Diagnostic :} scintigraphie (traceurs radioactifs comme le technétium-99m), tomographie par émission de positons (TEP).
            \item \emph{Thérapie :} \cle{radiothérapie} (irradier les tumeurs pour détruire les cellules cancéreuses), curiethérapie (sources placées au contact de la tumeur), traitements métaboliques (iode-131 contre certains cancers de la thyroïde).
        \end{itemize}
    \item \textbf{Armement :} bombe A (\emph{fission}, Hiroshima et Nagasaki, 1945), bombe H (\emph{fusion thermonucléaire}, premiers essais américain en 1952 et soviétique en 1953). Course aux armements pendant la Guerre froide (environ 70 000 têtes nucléaires dans le monde au milieu des années 1980), puis maîtrise des armements : traité de non-prolifération (TNP, 1968), traités de réduction (SALT, START).
\end{itemize}

\begin{attention}[Nucléaire civil et militaire : la dualité technologique]
Les mêmes techniques d'enrichissement de l'uranium produisent du combustible civil (environ 3 à 5 \% d'uranium-235) ou de la matière fissile militaire (plus de 90 \% d'uranium-235). Le retraitement du combustible usé permet d'extraire du plutonium utilisable dans une arme. C'est le cœur du problème de la \cle{non-prolifération} : le « droit à l'énergie nucléaire » reconnu par le TNP peut ouvrir la voie à une capacité militaire. Au Bac, distingue toujours les \textbf{centrales électriques} (usage civil) des \textbf{réacteurs de production} de plutonium (usage militaire).
\end{attention}

\begin{exoresolu}[Analyse de document : le discours « Atoms for Peace » (Eisenhower, ONU, 8 décembre 1953)]
\textbf{Énoncé :} « Les États-Unis [...] souhaitent [...] consacrer leur force à servir la paix [...] en fournissant des matériaux de fission à d'autres nations pour des usages pacifiques. » Comment ce discours illustre-t-il la stratégie américaine face à l'URSS en 1953 ?
\tcblower
\textbf{Solution détaillée :}
\begin{enumerate}
    \item \textbf{Contexte :} mort de Staline (mars 1953) et dégel relatif ; l'URSS vient de réussir un essai thermonucléaire (août 1953) ; la peur nucléaire domine l'opinion mondiale.
    \item \textbf{Stratégie :} « Atoms for Peace » (« Des atomes pour la paix ») est une diplomatie scientifique. Objectifs : conserver le leadership technologique américain, encadrer la diffusion du nucléaire civil (ce qui mènera à la création de l'AIEA en 1957) et légitimer l'arsenal américain en opposant l'« atome pacifique » occidental à l'« atome menaçant » du bloc de l'Est.
    \item \textbf{Conséquence :} diffusion de réacteurs de recherche et de puissance vers les alliés des États-Unis (Europe, Japon, Iran du Shah...), ce qui pose les bases du futur régime de non-prolifération.
\end{enumerate}
\end{exoresolu}

\courssub{La conquête spatiale : de Spoutnik à l'ISS, enjeux et retombées}

\textbf{1. L'intuition orbitale : Newton, Tsiolkovski, Goddard, von Braun.}
Placer un objet en orbite suppose de lui donner une vitesse horizontale d'environ \SI{7,9}{km/s} en orbite basse, pour que sa chute gravitationnelle « rate » en permanence la surface de la Terre. Les travaux théoriques de Tsiolkovski (équation de la fusée, 1903) et les fusées à propergol liquide de Goddard (1926) en posent les bases. La fusée \cle{V2} allemande (1944) prouve la faisabilité technique ; après 1945, les ingénieurs (Wernher von Braun aux États-Unis, Sergueï Korolev en URSS) transforment le missile militaire en lanceur spatial.

\textbf{2. La frise des grands jalons (1957–1998).}

\begin{popfigure}
\begin{tikzpicture}[scale=0.9, every node/.style={font=\small}]
    % Axe principal
    \draw[thick, ->] (0,0) -- (16,0) node[right] {Années};
    \foreach \x/\year in {0/1957, 2/1961, 4.5/1969, 8/1975, 11/1981, 14/1998} {
        \draw (\x,0.1) -- (\x,-0.1) node[below] {\year};
    }
    
    % Evenements HAUT
    \node[align=center, popblue, anchor=south] (s1) at (0, 1.5) {\textbf{1957}\\\textbf{Spoutnik 1}\\(URSS)\\\textit{Premier satellite}};
    \draw[popblue, thick] (0,0) -- (s1.south);
    
    \node[align=center, popblue, anchor=south] (s2) at (2, 2.5) {\textbf{1961}\\\textbf{Youri Gagarine}\\(Vostok 1)\\\textit{Premier homme dans l'espace}};
    \draw[popblue, thick] (2,0) -- (s2.south);
    
    \node[align=center, poporange, anchor=south] (s3) at (4.5, 1.5) {\textbf{1969}\\\textbf{Apollo 11}\\(États-Unis)\\\textit{Premiers pas sur la Lune}};
    \draw[poporange, thick] (4.5,0) -- (s3.south);
    
    \node[align=center, poppurple, anchor=south] (s4) at (8, 2.5) {\textbf{1975}\\\textbf{Apollo-Soyouz}\\\textit{Premier rendez-vous}\\\textit{USA-URSS}};
    \draw[poppurple, thick] (8,0) -- (s4.south);
    
    \node[align=center, popgreen, anchor=south] (s5) at (11, 1.5) {\textbf{1981}\\\textbf{Navette Columbia}\\\textit{Premier vol spatial}\\\textit{réutilisable}};
    \draw[popgreen, thick] (11,0) -- (s5.south);
    
    \node[align=center, poppink, anchor=south] (s6) at (14, 2.5) {\textbf{1998}\\\textbf{ISS (Zarya/Unity)}\\\textit{Station spatiale}\\\textit{internationale}};
    \draw[poppink, thick] (14,0) -- (s6.south);
    
    % Evenements BAS
    \node[align=center, popmuted, anchor=north] (b1) at (1, -1.5) {\textbf{1958}\\\textit{Explorer 1 (USA)}\\\textit{Ceintures de Van Allen}};
    \draw[popmuted, thick] (1,0) -- (b1.north);
    
    \node[align=center, popmuted, anchor=north] (b2) at (3.5, -2.5) {\textbf{1962}\\\textit{Telstar 1}\\\textit{Première TV transatlantique}};
    \draw[popmuted, thick] (3.5,0) -- (b2.north);
    
    \node[align=center, popmuted, anchor=north] (b3) at (6, -1.5) {\textbf{1977}\\\textit{Voyager 1 et 2}\\\textit{Grande tournée des planètes}};
    \draw[popmuted, thick] (6,0) -- (b3.north);
    
    \node[align=center, popmuted, anchor=north] (b4) at (9, -2.5) {\textbf{1990}\\\textit{Hubble}\\\textit{Télescope spatial}};
    \draw[popmuted, thick] (9,0) -- (b4.north);
    
    \node[align=center, popmuted, anchor=north] (b5) at (12.5, -1.5) {\textbf{1997}\\\textit{Pathfinder / Sojourner}\\\textit{Premier rover sur Mars}};
    \draw[popmuted, thick] (12.5,0) -- (b5.north);
    
    % Legende couleurs
    \node[align=left, anchor=north west] at (0.2, -4.2) {
        \textcolor{popblue}{■ URSS / premières mondiales} \quad \textcolor{poporange}{■ États-Unis / Lune} \quad \textcolor{poppurple}{■ Coopération} \quad \textcolor{popgreen}{■ Réutilisable} \quad \textcolor{poppink}{■ Station permanente} \quad \textcolor{popmuted}{■ Satellites et sondes}
    };
\end{tikzpicture}
\legende{Frise chronologique des grands jalons de la conquête spatiale (1957–1998). La période 1957–1969 est marquée par la compétition (« course à l'espace » URSS/USA) ; après 1975, la coopération s'installe (Apollo-Soyouz, puis ISS).}
\end{popfigure}

\begin{definition}[Satellite artificiel]
Un \cle{satellite artificiel} est un engin spatial placé volontairement en orbite autour de la Terre (ou d'un autre corps céleste) pour y remplir des missions de \textbf{télécommunications}, d'\textbf{observation de la Terre} (météorologie, environnement, renseignement), de \textbf{navigation} (GPS, Galileo, Glonass, Beidou) ou d'\textbf{exploration scientifique}.
\end{definition}

\begin{exemplebox}[Exemple chiffré : l'explosion du nombre de satellites]
En 1957, un seul objet : \textbf{Spoutnik 1} (environ 84 kg, émettant un simple signal radio « bip-bip »).
Au milieu des années 2020 : plusieurs milliers de satellites actifs (dont une grande partie appartenant à la constellation Starlink de SpaceX) et plus de 30 000 débris de plus de 10 cm suivis depuis le sol. L'orbite basse se sature : risque de \cle{syndrome de Kessler} (collisions en cascade). La gestion du trafic spatial devient un enjeu réglementaire majeur (COPUOS, comité des Nations unies pour les utilisations pacifiques de l'espace).
\end{exemplebox}

\textbf{3. Les missions Apollo et le premier pas sur la Lune (1969).}

\begin{popfigure}
\begin{tikzpicture}[scale=1.1, every node/.style={font=\small}]
    % Fond lunaire
    \fill[popmuted!20] (-3,-1.5) rectangle (3,0);
    \draw[thick, popmuted] (-3,0) -- (3,0);
    % Module lunaire (LEM) simplifié
    \draw[thick, fill=popcream] (-0.5,0) -- (-0.3,0.8) -- (0.3,0.8) -- (0.5,0) -- cycle;
    \draw[thick, fill=popgoldL] (-1,0) rectangle (1,0.3);
    \draw[thick] (-0.8,0) -- (-1.2,-0.3);
    \draw[thick] (0.8,0) -- (1.2,-0.3);
    \draw[thick] (-0.2,0) -- (-0.5,-0.3);
    \draw[thick] (0.2,0) -- (0.5,-0.3);
    
    % Astronaute
    \draw[thick, fill=white] (-0.1,0.3) circle (0.15);
    \draw[thick] (-0.1,0.15) -- (-0.1,0);
    \draw[thick] (-0.1,0.1) -- (-0.3,0);
    \draw[thick] (-0.1,0.1) -- (0.2,0.2);
    \draw[thick] (-0.1,0) -- (-0.2,-0.3);
    \draw[thick] (-0.1,0) -- (0.05,-0.3);
    
    % Drapeau
    \draw[thick] (0.6,0) -- (0.6,0.8);
    \draw[thick, fill=popblue!50] (0.6,0.5) -- (0.9,0.65) -- (0.6,0.8) -- cycle;
    
    % Terre au loin
    \draw[thick, fill=popblue!30] (2.2, 1.5) circle (0.4);
    \node at (2.2, 1.5) {\tiny Terre};
    
    % Annotations
    \node[align=left, anchor=south west] at (-2.8, 1.2) {
        \textbf{Apollo 11 -- 21 juillet 1969 (2 h 56 UTC)}\\
        Neil Armstrong descend l'échelle du LEM \textit{Eagle}.\\
        « C'est un petit pas pour l'homme,\\
        un bond de géant pour l'humanité. »\\
        Buzz Aldrin le rejoint 19 minutes plus tard.\\
        Environ 2 h 30 de sortie, 21,5 kg d'échantillons.
    };
    \node[align=left, anchor=north west] at (-2.8, -1.8) {
        \textbf{Analyse de document (photo NASA) :}\\
        1. \textit{Source :} photographie prise sur la Lune par l'équipage d'Apollo 11.\\
        2. \textit{Contenu :} astronaute, LEM, drapeau, sol gris cratérisé, ciel noir.\\
        3. \textit{Portée :} victoire américaine dans la course à la Lune (annonce de Kennedy, 1961).\\
        4. \textit{Limites :} coût énorme du programme ; dernier alunissage en 1972 (Apollo 17).
    };
\end{tikzpicture}
\legende{Croquis commenté d'après les photographies de la mission Apollo 11 (Mer de la Tranquillité, juillet 1969). L'absence d'atmosphère (ciel noir en plein jour, ombres nettes), la faible gravité (démarche bondissante) et le drapeau « flottant » (maintenu par une tige horizontale rigide) sont des marqueurs de l'environnement lunaire.}
\end{popfigure}

\textbf{4. Les satellites : télécommunications, météo, observation, navigation.}
\begin{itemize}
    \item \textbf{Télécommunications :} l'écrivain et ingénieur Arthur C. Clarke imagine dès 1945 le satellite géostationnaire (orbite à environ 36 000 km d'altitude, fixe au-dessus d'un point de l'équateur). \textit{Syncom 3} retransmet les Jeux olympiques de Tokyo (1964) ; \textit{Intelsat} lance le service commercial (« Early Bird », 1965). Aujourd'hui : constellations en orbite basse (Starlink, OneWeb) pour l'internet haut débit mondial.
    \item \textbf{Météorologie :} \textit{TIROS-1} (1960) fournit les premières images de nuages vues de l'espace ; suivent \textit{Meteosat} (Europe) et \textit{GOES} (États-Unis). Applications : prévision du temps, suivi des cyclones, vigilance contre les crues.
    \item \textbf{Observation de la Terre (télédétection) :} \textit{Landsat} (depuis 1972), \textit{SPOT} (France, 1986), \textit{Sentinel} (programme européen Copernicus, depuis 2014). Applications : suivi des cultures, de l'urbanisation, de la déforestation, des ressources en eau et des risques naturels.
    \item \textbf{Navigation par satellite (GNSS) :} \textit{GPS} (États-Unis, opérationnel en 1995), \textit{Glonass} (Russie), \textit{Galileo} (Union européenne), \textit{Beidou} (Chine). Principe : le récepteur calcule sa position à partir du temps de trajet des signaux radio émis par au moins quatre satellites.
\end{itemize}

\begin{savaistu}[Le Cameroun et l'espace : des yeux sur le territoire]
Le Cameroun ne lance pas de satellites, mais \textbf{utilise les données spatiales} :
\begin{itemize}
    \item \textbf{Désertification (Extrême-Nord) :} suivi par imagerie satellitaire de l'avancée des sables et de la régression spectaculaire du lac Tchad (qui a perdu l'essentiel de sa superficie depuis les années 1960), au service des programmes de reboisement (Grande Muraille verte).
    \item \textbf{Déforestation (Sud-Est, bassin du Congo) :} cartographie de l'exploitation forestière et de l'agro-industrie, protection des parcs nationaux (Boumba-Bek, Nki, Lobéké).
    \item \textbf{Agriculture et sécurité alimentaire :} suivi de l'état des cultures (mil, sorgho, maïs) et alerte précoce en cas de mauvaises récoltes.
    \item \textbf{Gestion urbaine :} cartographie de l'étalement de Yaoundé et Douala (risques d'inondations et de glissements de terrain).
\end{itemize}
Ces applications montrent qu'une innovation technique mondiale peut servir directement le développement d'un pays africain.
\end{savaistu}

\textbf{5. Stations spatiales et sondes : présence permanente et exploration robotique.}
\begin{itemize}
    \item \textbf{Stations :} Saliout (URSS, 1971), Skylab (États-Unis, 1973), Mir (URSS puis Russie, 1986–2001), puis l'\textbf{ISS} (Station spatiale internationale, depuis 1998, coopération de 15 pays, environ 420 tonnes, à environ 400 km d'altitude) : recherches en microgravité, biologie, matériaux, observation de la Terre. La Chine exploite sa propre station, Tiangong, depuis 2021.
    \item \textbf{Sondes interplanétaires :} \textit{Voyager 1 et 2} (1977, aujourd'hui dans l'espace interstellaire), \textit{Cassini-Huygens} (Saturne et Titan), \textit{Curiosity} et \textit{Perseverance} (Mars), télescope spatial \textit{James Webb} (lancé en 2021, observation de l'univers primordial en infrarouge).
\end{itemize}

\textbf{6. Retombées civiles de la recherche spatiale (transferts de technologie).}
Les contraintes extrêmes de l'espace (masse limitée, fiabilité absolue, vide, rayonnement, températures extrêmes) obligent à innover ; ces innovations redescendent ensuite vers la vie quotidienne :
\begin{center}
\begin{tabularx}{\linewidth}{|l|X|X|}\hline
\rowcolor{popblueL}
\textbf{Domaine spatial} & \textbf{Innovation contrainte} & \textbf{Application civile} \\ \hline
Électronique & Miniaturisation, résistance aux radiations & Circuits intégrés, capteurs d'appareils photo et d'imagerie médicale \\ \hline
Matériaux & Légèreté, résistance thermique & Mousses à mémoire de forme, fibres textiles très résistantes, revêtements anti-rayures \\ \hline
Énergie & Panneaux solaires, batteries performantes & Photovoltaïque terrestre, véhicules électriques \\ \hline
Télécommunications & Liaisons fiables, correction d'erreurs & GPS grand public, télévision par satellite, internet mobile \\ \hline
Médecine & Surveillance médicale des astronautes & Télémédecine, progrès de l'imagerie médicale, prothèses \\ \hline
Eau et environnement & Recyclage de l'eau et de l'air à bord des stations & Purification de l'eau dans les zones arides, capteurs de qualité de l'air \\ \hline
\end{tabularx}
\end{center}

\begin{attention}[Course à l'espace = course aux armements (Guerre froide)]
Ne pas réduire la conquête spatiale à l'exploration scientifique.
\begin{itemize}
    \item \textbf{Lanceurs et missiles sont cousins :} la fusée R-7 qui lance Spoutnik est aussi le premier missile balistique intercontinental soviétique ; les fusées américaines Atlas et Titan sont des missiles adaptés.
    \item \textbf{Reconnaissance :} les satellites espions permettent de surveiller l'adversaire et de vérifier l'application des traités de désarmement.
    \item \textbf{Traité de l'espace (1967) :} il interdit le déploiement d'armes de destruction massive dans l'espace et l'appropriation de la Lune par un État, et affirme l'usage pacifique de l'espace.
\end{itemize}
Au Probatoire et au Bac, on attend que vous reliiez \textbf{progrès technique}, \textbf{puissance des États} et \textbf{droit international}.
\end{attention}

\begin{methode}[Commenter un document iconographique spatial (type Bac)]
\begin{enumerate}
    \item \textbf{Identifier :} nature du document (photo, schéma, carte), source (NASA, agence spatiale, presse), date, auteur si connu.
    \item \textbf{Décrire objectivement :} ce qui est visible (objets, personnes, paysage, couleurs, ombres), avec un vocabulaire technique précis (module lunaire, casque, panneau solaire, antenne).
    \item \textbf{Contextualiser :} mission, date, enjeu géopolitique (Guerre froide), prouesse technique.
    \item \textbf{Analyser la portée :} pourquoi cette image a-t-elle été prise et diffusée ? Propagande, preuve scientifique, symbole pour l'humanité ?
    \item \textbf{Conclure et ouvrir :} héritage (technologies, droit de l'espace, retour vers la Lune, exploration de Mars).
\end{enumerate}
\end{methode}

\begin{exercice}[Entraînement : document iconographique]
\textbf{Document :} photographie du « lever de Terre » (\textit{Earthrise}), prise par William Anders depuis la capsule Apollo 8 en orbite lunaire, le 24 décembre 1968.
\textbf{Consigne :} en suivant la méthode ci-dessus, rédigez un commentaire structuré (introduction, description, analyse, conclusion) d'environ vingt lignes montrant en quoi cette image est un « document fondateur » de la conscience planétaire.
\end{exercice}

\begin{corrige}[Exercice : Earthrise (Apollo 8)]
\textbf{Introduction :} photographie prise le 24 décembre 1968 par l'astronaute William Anders, lors d'Apollo 8, première mission habitée à orbiter autour de la Lune. Contexte : Guerre froide et année 1968 tragique (guerre du Vietnam, printemps de Prague écrasé, assassinats de Martin Luther King et de Robert Kennedy).
\textbf{Description :} au premier plan, la surface lunaire grise, cratérisée, sans vie, à l'horizon net (absence d'atmosphère). À l'arrière-plan, la Terre apparaît comme une bille bleue, blanche et brune, partiellement éclairée, suspendue dans le noir absolu. Le contraste est saisissant : la vie, l'eau et l'atmosphère face au vide minéral.
\textbf{Analyse :} 1. \emph{Portée scientifique :} pour la première fois, des humains voient la Terre entière « de l'extérieur ». 2. \emph{Portée symbolique :} les frontières sont invisibles ; la biosphère apparaît fragile, entourée d'une fine couche d'atmosphère. Cette image nourrit le mouvement écologiste naissant (premier Jour de la Terre en 1970, rapport du Club de Rome en 1972). 3. \emph{Portée géopolitique :} victoire américaine dans la course à la Lune, mais message universel porté par la mission.
\textbf{Conclusion :} \textit{Earthrise} fait basculer le regard : la Terre devient un « vaisseau spatial » unique à gérer collectivement. Cette image est l'une des matrices de la prise de conscience environnementale mondiale, qui mènera aux grandes conférences internationales sur l'environnement (Sommet de Rio, 1992 ; Accord de Paris, 2015).
\end{corrige}

\coursec{Les grandes découvertes scientifiques : révolution de la biologie et de la médecine}

Après avoir maîtrisé l'énergie de l'atome et ouvert la voie de l'espace, les savants du monde d'après 1945 se sont tournés vers un territoire encore plus intime : le \cle{vivant} lui-même. Comprendre comment un être vivant transmet ses caractères à ses descendants, soigner des maladies autrefois mortelles, voir à l'intérieur du corps humain sans l'ouvrir : telle est l'ambition de la \cle{révolution de la biologie et de la médecine}, qui transforme profondément l'existence des hommes depuis le milieu du \textsc{xx}\textsuperscript{e} siècle. Cette révolution repose sur une découverte fondatrice, celle de la structure de l'\cle{ADN}, puis sur des applications techniques de plus en plus puissantes, qui soulèvent aujourd'hui de grandes questions éthiques.

\courssub{La découverte de la structure de l'ADN (1953) : naissance de la biologie moléculaire}

\textbf{Une énigme vieille comme la science.} Depuis longtemps, les savants savaient que les caractères (couleur des yeux, groupe sanguin, certaines maladies) se transmettent des parents aux enfants. Mais par quel support matériel ? Au début du \textsc{xx}\textsuperscript{e} siècle, les biologistes ont identifié dans le noyau des cellules une molécule candidate : l'\cle{ADN}. Restait à comprendre sa forme, car c'est la forme d'une molécule qui explique son fonctionnement.

\begin{definition}[L'ADN]
L'\cle{ADN} (acide désoxyribonucléique) est la molécule qui porte l'\cle{information génétique} de tout être vivant. Elle est formée de deux brins enroulés l'un autour de l'autre en \cle{double hélice}. Chaque brin est un enchaînement de \cle{nucléotides} portant quatre types de \cle{bases} : adénine (A), thymine (T), guanine (G) et cytosine (C). Les bases s'associent toujours deux à deux : A avec T, G avec C. L'ordre des bases constitue le \cle{code génétique}, véritable « texte » de la vie.
\end{definition}

\textbf{La démarche d'investigation : de l'image au modèle.} En 1953, à Cambridge (Grande-Bretagne), deux jeunes chercheurs, l'Américain James \cle{Watson} et le Britannique Francis \cle{Crick}, s'appuient sur les travaux de \cle{cristallographie aux rayons X} menés à Londres par Rosalind \cle{Franklin} et Maurice Wilkins. La célèbre « photographie 51 » de Franklin montre une figure en croix : c'est la signature d'une structure en \cle{hélice}. Watson et Crick construisent alors un modèle en tiges et en plaques métalliques, qu'ils ajustent jusqu'à ce qu'il soit compatible avec toutes les données : deux brins complémentaires, enroulés en double hélice. Ils publient leur découverte en avril 1953 dans la revue \emph{Nature}. Watson, Crick et Wilkins reçoivent le prix Nobel de médecine en 1962 ; Rosalind Franklin, morte en 1958, ne peut être récompensée, le Nobel n'étant pas attribué à titre posthume.

\begin{attention}[Piège fréquent]
Ne pas attribuer la découverte de l'ADN lui-même à Watson et Crick : la molécule était connue depuis le \textsc{xix}\textsuperscript{e} siècle. Leur mérite est d'avoir découvert sa \cle{structure} en double hélice, ce qui a permis de comprendre comment l'information génétique est stockée et copiée. Autre piège : oublier le rôle essentiel de Rosalind Franklin.
\end{attention}

\textbf{Pourquoi est-ce une révolution ?} La structure en double hélice explique immédiatement deux choses : comment l'information est codée (par la suite des bases) et comment elle est copiée (chaque brin sert de modèle pour fabriquer le brin complémentaire lors de la division cellulaire). Cette découverte fonde la \cle{biologie moléculaire}, science qui étudie le vivant à l'échelle des molécules, et ouvre la voie à toutes les biotechnologies modernes.

\begin{popfigure}
\begin{center}
\begin{tikzpicture}[scale=1.0]
% Brin gauche (sinusoïde décalée)
\draw[very thick,popblue] plot[domain=0:6.28,samples=100] ({0.9*sin(\x r)}, \x*0.55);
% Brin droit
\draw[very thick,poporange] plot[domain=0:6.28,samples=100] ({-0.9*sin(\x r)}, \x*0.55);
% Barreaux (bases)
\foreach \t in {0.4,0.9,1.4,1.9,2.4,2.9,3.4,3.9,4.4,4.9,5.4,5.9}{
  \draw[popmuted] ({0.9*sin(\t r)}, \t*0.55) -- ({-0.9*sin(\t r)}, \t*0.55);
}
% Étiquettes de bases
\node[popdark,font=\small] at (0, 0.4*0.55) {A--T};
\node[popdark,font=\small] at (0, 1.4*0.55) {G--C};
\node[popdark,font=\small] at (0, 2.4*0.55) {T--A};
\node[popdark,font=\small] at (0, 3.4*0.55) {C--G};
% Légendes des brins
\node[popblue,font=\small\bfseries] at (2.6, 3.2) {Brin 1};
\draw[->,popblue] (2.3,3.1) -- (0.95,3.0);
\node[poporange,font=\small\bfseries] at (-2.6, 1.6) {Brin 2};
\draw[->,poporange] (-2.3,1.5) -- (-0.95,1.4);
\node[popmuted,font=\small,align=center] at (3.4, 0.6) {Paires de\\bases};
\draw[->,popmuted] (2.9,0.55) -- (0.5,0.22);
\end{tikzpicture}
\end{center}
\legende{Schéma simplifié de la double hélice d'ADN : deux brins complémentaires reliés par des paires de bases (A--T et G--C), comme les barreaux d'une échelle enroulée sur elle-même.}
\end{popfigure}

\courssub{Le génie génétique : manipuler le vivant}

Comprendre le code génétique, c'est bien ; pouvoir le modifier, c'est une puissance inédite. Le \cle{génie génétique} est l'ensemble des techniques qui permettent de lire, découper, copier et modifier l'ADN.

\begin{itemize}
\item \textbf{Années 1970 : les premières manipulations.} En 1973, les Américains Stanley Cohen et Herbert Boyer réalisent la première \cle{transgenèse} : ils insèrent un fragment d'ADN étranger dans une bactérie, qui le reproduit. C'est la naissance du \cle{génie génétique} et des \cle{OGM} (organismes génétiquement modifiés). En 1978, la première \cle{insuline} humaine produite par des bactéries modifiées révolutionne le traitement du diabète.
\item \textbf{1996 : le clonage de la brebis Dolly.} En Écosse, l'équipe d'Ian Wilmut obtient la brebis \cle{Dolly}, premier mammifère \cle{cloné} à partir d'une cellule adulte : le noyau d'une cellule de mamelle est transféré dans un ovule vidé de son propre noyau. Dolly prouve qu'une cellule adulte contient encore tout le programme génétique de l'individu.
\item \textbf{2003 : le génome humain séquencé.} Le programme international \cle{Génome humain}, lancé en 1990, achève en 2003 le \cle{séquençage} complet du génome humain : la lecture des quelque 3 milliards de paires de bases qui constituent notre ADN. Cette « carte » du vivant humain permet d'identifier les gènes responsables de maladies héréditaires et ouvre l'ère de la \cle{médecine personnalisée}.
\end{itemize}

\begin{aretenir}[Repères chronologiques]
1953 : structure de l'ADN (Watson, Crick, à partir des données de Franklin). 1973 : première transgenèse. 1978 : insuline humaine par génie génétique. 1996 : clonage de la brebis Dolly. 2003 : séquençage complet du génome humain.
\end{aretenir}

\courssub{Les progrès médicaux : vaccins, antibiotiques, greffes}

La révolution biologique se traduit d'abord par des victoires concrètes contre la maladie et la mort.

\textbf{Les vaccins.} La \cle{vaccination} consiste à présenter à l'organisme un agent infectieux tué, affaibli ou un simple fragment, pour que le corps fabrique des défenses (anticorps) sans tomber malade. Après 1945, les campagnes de vaccination changent la face du monde : le vaccin contre la \cle{poliomyélite} de Jonas Salk (1955), puis celui d'Albert Sabin (1961, par voie orale), font reculer une maladie qui paralysait des centaines de milliers d'enfants ; le vaccin contre la \cle{rougeole} (1963) sauve des millions de vies. La \cle{variole} est même totalement \cle{éradiquée} de la planète en 1980, grâce à une campagne mondiale de l'Organisation mondiale de la santé (OMS) : c'est la première maladie éliminée par l'action humaine.

\textbf{Les antibiotiques.} Découverte en 1928 par Alexander Fleming, la \cle{pénicilline} n'est produite en masse que pendant la Deuxième Guerre mondiale. Après 1945, les antibiotiques se \cle{généralisent} : pneumonies, tuberculose, infections des plaies cessent d'être des condamnations à mort. Mais leur usage excessif fait apparaître des bactéries \cle{résistantes}, nouveau défi de la médecine actuelle.

\textbf{Les greffes d'organes.} Remplacer un organe malade par un organe sain : longtemps impossible à cause du \cle{rejet} (le système immunitaire attaque l'organe étranger), la greffe devient réalité grâce aux progrès de la chirurgie et aux médicaments \cle{immunosuppresseurs}. La première greffe de rein réussie date de 1954 (États-Unis). Le 3 décembre 1967, le chirurgien sud-africain Christiaan \cle{Barnard} réalise au Cap la première \cle{greffe du cœur} de l'histoire : le patient survit 18 jours. L'exploit, relayé par la presse mondiale, fait date. Aujourd'hui, les greffes de cœur, de rein, de foie ou de cornée sauvent chaque année des dizaines de milliers de malades.

\courssub{L'imagerie médicale : voir sans ouvrir}

Jusqu'au milieu du \textsc{xx}\textsuperscript{e} siècle, pour savoir ce qui se passe dans le corps d'un malade, il faut souvent opérer. L'\cle{imagerie médicale} rend le corps transparent.

\begin{itemize}
\item \textbf{L'échographie} (diffusée à partir des années 1960) utilise des \cle{ultrasons} réfléchis par les tissus. Sans danger, elle permet de suivre la grossesse et d'examiner le cœur ou le foie.
\item \textbf{Le scanner} (1971) : l'ingénieur britannique Godfrey \cle{Hounsfield} met au point le \cle{tomodensitomètre}, qui combine des rayons X et un ordinateur pour reconstruire des images en coupes du corps. La première image de cerveau est obtenue en 1971. Hounsfield reçoit le prix Nobel en 1979.
\item \textbf{L'IRM} (imagerie par résonance magnétique, années 1970-1980) utilise un puissant champ magnétique et des ondes radio : elle donne des images très précises des tissus mous (cerveau, muscles, articulations) sans rayons X.
\end{itemize}

Ces techniques, qui doivent beaucoup à l'informatique, permettent des \cle{diagnostics précoces} : détecter une tumeur de quelques millimètres, c'est souvent la guérir.

\courssub{La lutte contre les grandes pandémies : sida et Covid-19}

La science médicale est aussi mise à l'épreuve par de nouvelles épidémies mondiales, les \cle{pandémies}.

\textbf{Le sida.} En 1981, des médecins américains décrivent une maladie nouvelle qui détruit les défenses immunitaires : le \cle{sida} (syndrome d'immunodéficience acquise). En 1983, l'équipe française de Luc Montagnier (Institut Pasteur) identifie le virus responsable, le \cle{VIH}. Il n'existe pas de vaccin, mais à partir de 1996, les \cle{trithérapies} (combinaisons d'antirétroviraux) transforment le sida en maladie chronique avec laquelle on peut vivre longtemps. L'Afrique subsaharienne, dont le Cameroun, reste la région la plus touchée : la lutte repose sur le dépistage, la prévention et l'accès aux traitements.

\textbf{La Covid-19.} Fin 2019, un nouveau coronavirus apparaît en Chine et provoque en 2020 une pandémie mondiale qui bouleverse l'économie et la vie quotidienne (confinements, fermeture des frontières). La réponse scientifique est d'une rapidité inédite : le génome du virus est séquencé dès janvier 2020, et des vaccins sont disponibles avant la fin de l'année.

\begin{savaistu}[Un record historique : les vaccins à ARN messager]
Les vaccins contre la Covid-19 (Pfizer-BioNTech, Moderna) ont été mis au point en \cle{moins d'un an}, un record absolu — il fallait auparavant 5 à 10 ans pour créer un vaccin. Ils utilisent la technologie de l'\cle{ARN messager} : au lieu d'injecter le virus affaibli, on injecte un « mode d'emploi » génétique qui apprend aux cellules à fabriquer un fragment inoffensif du virus, déclenchant la réponse immunitaire. Cette technologie, étudiée depuis les années 1990, ouvre des espoirs contre le paludisme et certains cancers.
\end{savaistu}

\courssub{Les enjeux éthiques : jusqu'où aller ?}

Pouvoir lire et modifier le vivant pose des questions que la science seule ne peut pas trancher : ce qui est techniquement possible est-il moralement acceptable ?

\begin{itemize}
\item \textbf{Les OGM.} Les plantes génétiquement modifiées (résistantes aux insectes ou à la sécheresse) peuvent augmenter les récoltes, mais suscitent des débats : risques pour l'environnement, dépendance des paysans envers les firmes qui vendent les semences, précaution pour la santé.
\item \textbf{La manipulation du génome humain.} Le clonage reproductif humain est interdit dans la plupart des pays. En 2018, un chercheur chinois annonce la naissance de bébés dont l'ADN a été modifié : l'annonce provoque une condamnation mondiale. Modifier le génome d'embryons, c'est transformer l'hérédité des générations futures.
\item \textbf{La procréation médicalement assistée (PMA).} Depuis la naissance en 1978 de Louise Brown, premier « bébé-éprouvette » (fécondation \emph{in vitro}), la PMA aide les couples stériles. Mais elle soulève des questions : tri des embryons, gestation pour autrui, accès réservé aux pays riches.
\end{itemize}

Ces débats ont conduit à la création de \cle{comités d'éthique} et à des lois de \cle{bioéthique} qui encadrent la recherche. La science avance, mais sous le regard de la société.

\courssub{Un progrès mesurable mais inégalement partagé}

Les effets de la révolution biologique et médicale se lisent dans les statistiques : les hommes vivent plus longtemps et les enfants meurent moins.

\begin{exemplebox}[L'espérance de vie en hausse]
L'\cle{espérance de vie} mondiale passe d'environ \textbf{46 ans en 1950} à plus de \textbf{72 ans en 2020}. Au Cameroun, elle passe d'environ \textbf{40 ans} à plus de \textbf{60 ans} sur la même période. Ce gain spectaculaire s'explique par les vaccins, les antibiotiques, l'hygiène et la baisse de la mortalité infantile.
\end{exemplebox}

\begin{popfigure}
\begin{center}
\begin{tikzpicture}
\begin{axis}[
  ybar, bar width=8pt,
  width=0.95\linewidth, height=6.2cm,
  ymin=0, ymax=220,
  ylabel={Décès pour 1\,000 naissances},
  symbolic x coords={1960,1980,2000,2020},
  xtick=data,
  legend style={at={(0.5,1.02)},anchor=south,legend columns=2,font=\small},
  ymajorgrids, grid style={popline!40},
  tick label style={font=\small},
  label style={font=\small}
]
\addplot+[ybar,fill=popblue,draw=popblue] coordinates {(1960,185) (1980,120) (2000,75) (2020,38)};
\addplot+[ybar,fill=poporange,draw=poporange] coordinates {(1960,180) (1980,140) (2000,150) (2020,75)};
\legend{Monde, Cameroun}
\end{axis}
\end{tikzpicture}
\end{center}
\legende{Baisse de la mortalité infantile (décès avant 5 ans pour 1\,000 naissances vivantes) dans le monde et au Cameroun, 1960-2020 (ordres de grandeur, sources ONU/Banque mondiale). La baisse est spectaculaire, mais le Cameroun reste au-dessus de la moyenne mondiale et a connu une remontée dans les années 1990 (sida, crises économiques).}
\end{popfigure}

\begin{attention}[Un progrès inégalement partagé]
Le progrès médical n'est pas le même pour tous. L'accès aux soins reste difficile dans de nombreuses \cle{zones rurales camerounaises} : manque de médecins et d'hôpitaux, distances longues, coût des médicaments, coupures d'électricité qui empêchent la conservation des vaccins. Les scanners et IRM sont concentrés à Yaoundé et Douala. Au Bac, évitez donc de présenter la révolution médicale comme un succès uniforme : montrez toujours le contraste entre les prouesses techniques et les inégalités d'accès.
\end{attention}

\begin{exoresolu}[Analyser un document statistique]
\textbf{Énoncé.} D'après le graphique ci-dessus, la mortalité infantile au Cameroun passe d'environ 180 décès pour 1\,000 naissances en 1960 à environ 75 en 2020. 1) Calculez la baisse en pourcentage. 2) Expliquez deux causes de cette baisse. 3) Pourquoi le Cameroun reste-t-il au-dessus de la moyenne mondiale ?
\tcblower
\textbf{Solution détaillée.}
1) Baisse $= \dfrac{180-75}{180}\times 100 \approx 58\,\%$. La mortalité infantile a donc diminué d'environ 58\,\% en soixante ans.
2) Causes : les campagnes de \cle{vaccination} (polio, rougeole) soutenues par l'OMS et l'UNICEF ; la diffusion des \cle{antibiotiques} et des traitements contre le paludisme ; l'amélioration de l'hygiène et de l'accès à l'eau potable.
3) Le Cameroun reste au-dessus de la moyenne mondiale à cause des \cle{inégalités d'accès aux soins} : sous-effectif médical, pauvreté, zones rurales enclavées, et poids de maladies comme le paludisme et le sida. Conclusion : la révolution médicale bénéficie au Cameroun, mais de façon incomplète et inégale.
\end{exoresolu}

\begin{exercice}[Pour s'entraîner]
1) Rédigez un paragraphe construit (10-15 lignes) montrant que la découverte de la structure de l'ADN en 1953 est à l'origine d'une révolution scientifique majeure.
2) Classez ces événements par ordre chronologique : séquençage du génome humain ; clonage de Dolly ; première greffe du cœur ; structure de l'ADN ; vaccins à ARN messager.
3) « Les progrès de la médecine depuis 1945 profitent-ils également à tous les hommes ? » Répondez en une quinzaine de lignes en vous appuyant sur des exemples précis, dont le Cameroun.
\end{exercice}

\begin{corrige}[Exercice n 2]
Ordre chronologique : structure de l'ADN (1953) ; première greffe du cœur (1967) ; clonage de la brebis Dolly (1996) ; séquençage complet du génome humain (2003) ; vaccins à ARN messager contre la Covid-19 (2020).
\end{corrige}

\begin{aretenir}[L'essentiel de la section]
La révolution de la biologie naît de la découverte de la structure de l'\cle{ADN} (1953), qui fonde la biologie moléculaire. Elle débouche sur le \cle{génie génétique} (transgenèse, clonage de Dolly en 1996, génome humain séquencé en 2003), sur des progrès médicaux immenses (\cle{vaccins}, \cle{antibiotiques}, \cle{greffes} d'organes, \cle{imagerie médicale}) et sur des réponses rapides aux pandémies (\cle{sida}, \cle{Covid-19}). L'espérance de vie s'envole, mais le progrès reste inégalement partagé et soulève des questions \cle{éthiques} nouvelles (OGM, manipulation du génome, PMA).
\end{aretenir}

\coursec{Les grandes innovations techniques : informatique, internet et télécommunications}

La révolution de la biologie et de la médecine, étudiée précédemment, a profondément transformé la santé humaine depuis 1945. Mais ces découvertes scientifiques n'auraient pas connu une telle accélération sans un bouleversement parallèle : celui des \cle{techniques de l'information et de la communication (TIC)}. L'ordinateur, internet et le téléphone mobile ont changé la manière de travailler, de commercer, de s'informer et de vivre, y compris au Cameroun. C'est l'objet de cette section.

\courssub{De l'ordinateur géant au micro-ordinateur}

\textbf{Intuition.} Un ordinateur est une machine capable d'effectuer automatiquement des calculs et des traitements d'information à partir d'un programme. En 1945, une telle machine occupe une salle entière ; soixante ans plus tard, elle tient dans une poche. Comment une telle \cle{miniaturisation} a-t-elle été possible ?

\textbf{Les étapes de la révolution informatique.}

\begin{itemize}
\item \textbf{1945 : l'ENIAC.} Construit aux États-Unis pour l'armée, l'\cle{ENIAC} (Electronic Numerical Integrator and Computer) est l'un des premiers ordinateurs électroniques à usage général. Il pèse environ 30 tonnes, occupe plus de 140 m$^2$ et fonctionne avec près de 18\,000 tubes à vide (lampes électroniques fragiles et gourmandes en énergie). Il sert à calculer des trajectoires d'artillerie.
\item \textbf{1947 : le transistor.} Inventé dans les laboratoires Bell (États-Unis) par John Bardeen, Walter Brattain et William Shockley, le \cle{transistor} est un minuscule composant électronique à semi-conducteur qui remplace les tubes à vide. Plus petit, plus fiable, moins énergivore, il permet de réduire considérablement la taille des machines.
\item \textbf{1958 : le circuit intégré.} Jack Kilby (Texas Instruments) et Robert Noyce (Fairchild) mettent au point indépendamment le \cle{circuit intégré} (ou puce) : plusieurs composants électroniques sont gravés sur une seule plaquette de silicium. On passe ainsi de circuits câblés à la main à des circuits microscopiques.
\item \textbf{1971 : le microprocesseur.} La firme Intel met au point le \cle{microprocesseur} (Intel 4004) : l'ensemble des fonctions de calcul (unité arithmétique et logique, registres) d'un ordinateur tient désormais sur une seule puce. C'est la naissance de la micro-informatique.
\end{itemize}

\begin{aretenir}[La logique de la miniaturisation]
Chaque innovation (transistor 1947, circuit intégré 1958, microprocesseur 1971) a permis de concentrer toujours plus de puissance de calcul dans un volume toujours plus petit, à un coût toujours plus faible. C'est cette logique cumulative, décrite par la \cle{loi de Moore} (1965 : doublement du nombre de transistors sur une puce tous les deux ans), qui rend possible l'ordinateur personnel puis le smartphone.
\end{aretenir}

\begin{popfigure}
\begin{center}
\begin{tikzpicture}[scale=0.9, every node/.style={transform shape}]
% ENIAC : gros bloc
\draw[fill=popblueL, draw=popblue, thick] (0,0) rectangle (2.6,2.2);
\node[align=center, font=\small] at (1.3,1.1) {\textbf{ENIAC}\\1945\\30 tonnes\\140 m$^2$};
% flèche
\draw[-{Stealth}, very thick, poporange] (2.8,1.1) -- (3.7,1.1);
% transistor
\draw[fill=popgreenL, draw=popgreen, thick] (3.9,0.5) rectangle (5.7,1.7);
\node[align=center, font=\small] at (4.8,1.1) {\textbf{Transistor}\\1947\\quelques cm$^3$};
\draw[-{Stealth}, very thick, poporange] (5.9,1.1) -- (6.8,1.1);
% circuit intégré
\draw[fill=poppurpleL, draw=poppurple, thick] (7.0,0.6) rectangle (8.6,1.6);
\node[align=center, font=\small] at (7.8,1.1) {\textbf{Circuit}\\\textbf{intégré}\\1958};
\draw[-{Stealth}, very thick, poporange] (8.8,1.1) -- (9.7,1.1);
% microprocesseur
\draw[fill=popgoldL, draw=popgold, thick] (9.9,0.7) rectangle (11.3,1.5);
\node[align=center, font=\small] at (10.6,1.1) {\textbf{Micro-}\\\textbf{processeur}\\1971};
\draw[-{Stealth}, very thick, poporange] (11.5,1.1) -- (12.4,1.1);
% smartphone
\draw[fill=poppinkL, draw=poppink, thick, rounded corners] (12.6,0.3) rectangle (14.0,1.9);
\node[align=center, font=\small] at (13.3,1.1) {\textbf{Smart-}\\\textbf{phone}\\2007\\$\sim$150 g};
\end{tikzpicture}
\end{center}
\legende{Schéma de la miniaturisation : de l'ENIAC (30 tonnes, 1945) au smartphone (environ 150 g, 2007). Chaque étape correspond à une innovation clé : transistor, circuit intégré, microprocesseur.}
\end{popfigure}

\courssub{L'ordinateur personnel et sa démocratisation}

Grâce au microprocesseur, l'informatique sort des laboratoires, des universités et des grandes administrations pour entrer dans les foyers et les PME.

\begin{itemize}
\item \textbf{1976 : Apple.} Steve Jobs et Steve Wozniak fondent Apple et commercialisent l'Apple I (kit à assembler) puis l'Apple II (1977), l'un des premiers micro-ordinateurs grand public à succès, avec clavier, boîtier et graphismes couleur.
\item \textbf{1981 : l'IBM PC.} Le géant IBM lance son \cle{Personal Computer} (PC 5150), dont l'architecture ouverte (bus ISA, BIOS) devient une norme mondiale (\emph{standard} « PC compatible »). Le système d'exploitation MS-DOS est fourni par une jeune entreprise : Microsoft, de Bill Gates et Paul Allen.
\item \textbf{Années 1980-1990 : la démocratisation.} La baisse des prix, l'apparition des interfaces graphiques (Macintosh 1984, Windows 3.0 en 1990) et de la souris rendent l'ordinateur accessible au grand public. Il devient un outil de bureautique (traitement de texte, tableur), de jeu et bientôt de communication (modem).
\end{itemize}

\begin{attention}[Piège fréquent au Probatoire]
Ne pas confondre l'invention de l'ordinateur (machine de calcul électronique, années 1940) avec celle du micro-ordinateur ou ordinateur personnel (années 1970-1980). Au Bac, situe chaque innovation à sa date exacte : ENIAC = 1945 ; microprocesseur = 1971 ; Apple II = 1977 ; IBM PC = 1981.
\end{attention}

\courssub{La naissance d'internet et du Web}

\begin{definition}[Internet et Web]
\cle{Internet} est un réseau mondial interconnectant des millions d'ordinateurs (et aujourd'hui objets connectés), qui peuvent échanger des données selon des règles communes (protocoles). Le \cle{Web} (World Wide Web) est un \emph{service} fonctionnant \emph{sur} internet : il permet de consulter des documents hypertextes (pages reliées par des liens) à l'aide d'un navigateur. Internet est l'infrastructure (les routes) ; le Web est l'un de ses usages (un service de livraison de pages).
\end{definition}

\textbf{Les grandes étapes.}

\begin{itemize}
\item \textbf{1969 : Arpanet.} Le réseau \cle{Arpanet}, financé par l'agence DARPA du ministère américain de la Défense en pleine Guerre froide, relie quatre universités américaines (UCLA, Stanford, UCSB, Utah). C'est l'ancêtre d'internet : les informations y circulent par \emph{commutation de paquets}, ce qui rend le réseau résistant aux pannes (pas de centre unique).
\item \textbf{Années 1970-1980 : interconnexion des réseaux.} Les protocoles \cle{TCP/IP} (Vint Cerf et Bob Kahn, 1974), adoptés comme standard unique sur Arpanet le 1er janvier 1983, permettent de relier des réseaux hétérogènes : internet (« \emph{inter-network} », réseau des réseaux) est né.
\item \textbf{1989-1990 : le Web.} Au CERN (Genève), l'informaticien britannique \cle{Tim Berners-Lee} invente le World Wide Web : il définit les adresses (URL), le protocole de transfert (HTTP), le langage de balisage (HTML) et le premier navigateur. Il en fait un bien public, gratuit et non breveté.
\item \textbf{Années 1990-2000 : la mondialisation du réseau.} Les navigateurs graphiques (Mosaic 1993, Netscape) et les moteurs de recherche (Google, 1998) ouvrent internet au grand public. Le réseau devient planétaire : on parle de « révolution numérique » ou de « société de l'information ».
\end{itemize}

\courssub{La téléphonie mobile : du portable au smartphone}

\begin{itemize}
\item \textbf{Années 1980 : les premiers portables (1G).} Les premiers téléphones mobiles (Motorola DynaTAC 8000X, 1983) sont lourds (près de 1 kg), chers (environ 4000 \$ de l'époque) et réservés à une élite professionnelle. Ils fonctionnent sur des réseaux analogiques (1G), peu sûrs et de faible capacité.
\item \textbf{Années 1990 : la 2G (GSM).} Le standard numérique européen \cle{GSM} (Global System for Mobile Communications) permet la compression de la voix, le chiffrement et surtout le \cle{SMS} (Short Message Service). Le mobile se démocratise rapidement.
\item \textbf{Années 2000-2010 : 3G et 4G.} La 3G (UMTS, années 2000) ajoute l'internet mobile (débit ~ 2 Mb/s), la 4G (LTE, années 2010) apporte le haut débit mobile (~ 100 Mb/s), permettant la vidéo en streaming et les usages nomades intensifs.
\item \textbf{2007 : le smartphone.} Apple lance l'\cle{iPhone} : téléphone, ordinateur de poche, appareil photo, lecteur multimédia et navigateur internet réunis dans un écran tactile multipoint sans clavier physique. Le smartphone devient en quelques années l'objet technique le plus répandu au monde.
\item \textbf{Années 2020 : la 5G.} Déploiement du très haut débit (jusqu'à 1-10 Gb/s), latence très faible, connexion massive d'objets (IoT), usages industriels et véhicules autonomes.
\end{itemize}

\begin{exemplebox}[La révolution mobile au Cameroun : un « saut technologique »]
Au Cameroun, le taux de pénétration du téléphone mobile (nombre de cartes SIM actives pour 100 habitants) dépasse \textbf{80 \%} dans les années 2020, alors qu'il était quasi nul en 1995. Le mobile a supplanté le téléphone fixe, jamais largement développé hors des grandes villes : le pays est passé directement à l'ère du portable, phénomène appelé « \cle{saut technologique} » (\emph{leapfrogging}).
\end{exemplebox}

\courssub{Intelligence artificielle, big data et objets connectés}

À partir des années 2010, la puissance de calcul exponentielle (loi de Moore), la capacité de stockage massive et la quantité de données produites sur internet ouvrent une nouvelle étape.

\begin{itemize}
\item Le \cle{big data} (« mégadonnées ») désigne le stockage, le traitement et l'exploitation d'énormes volumes de données variées (recherches, achats, déplacements, capteurs), caractérisés par les « 3 V » : Volume, Vélocité, Variété. Utile aux entreprises (marketing ciblé) comme aux États (sécurité, planification).
\item L'\cle{intelligence artificielle} (IA), notamment l'apprentissage profond (\emph{deep learning}), regroupe des programmes capables d'apprendre à partir de données : reconnaissance d'images et de parole, traduction automatique, assistants vocaux, recommandation de contenu, diagnostic médical assisté, voitures autonomes.
\item La \cle{robotique} moderne, couplée à l'IA, automatise des tâches industrielles (cobots), médicales (chirurgie assistée Da Vinci) ou domestiques (aspirateurs, tondeuses).
\item Les \cle{objets connectés} (IoT : Internet of Things) — montres, compteurs intelligents, capteurs agricoles, véhicules — communiquent par internet sans intervention humaine.
\end{itemize}

\begin{attention}[Enjeux citoyens]
Ces technologies posent des questions nouvelles majeures : protection des données personnelles (RGPD en Europe, loi au Cameroun), surveillance de masse, biais algorithmiques, fracture numérique entre pays riches et pays pauvres (accès au haut débit, équipements, formation), impact environnemental (consommation énergétique des data centers, déchets électroniques).
\end{attention}

\courssub{L'irruption du numérique au Cameroun}

Le Cameroun s'insère dans la révolution numérique à partir de la fin des années 1990, avec un retard initial rattrapé par la téléphonie mobile.

\begin{itemize}
\item \textbf{1997 : arrivée d'internet.} Le Cameroun se connecte à internet, d'abord par des liaisons satellitaires lentes, coûteuses et à forte latence.
\item \textbf{Années 2000 : explosion de la téléphonie mobile.} Les opérateurs privés \cle{MTN} (arrivé en 2000) et \cle{Orange} (2000, ex-Camtel Mobile) déploient leurs réseaux 2G puis 3G ; le mobile connaît une croissance spectaculaire, comme le montre la courbe ci-dessous.
\item \textbf{Années 2010 : le Mobile Money.} Les services de paiement par téléphone (\cle{MTN Mobile Money}, \cle{Orange Money}) permettent de payer, transférer et épargner sans compte bancaire classique. C'est un levier majeur d'inclusion financière.
\item \textbf{Les câbles à fibre optique : la fin de l'isolement.} Le Cameroun est raccordé aux câbles sous-marins \cle{SAT-3/WASC} (2002) puis \cle{WACS} (West Africa Cable System, 2012), qui relient l'Afrique de l'Ouest à l'Europe (Portugal, Royaume-Uni) et à l'Asie, multipliant le débit international disponible et faisant baisser les coûts. Des dorsales nationales en fibre optique (projet \cle{Cameroun Digital Boost}, \cle{PAVIC}) relient progressivement les grandes villes et les chefs-lieux de région.
\end{itemize}

\begin{popfigure}
\begin{center}
\begin{tikzpicture}
\begin{axis}[
width=0.9\linewidth, height=6.5cm,
xlabel={Année}, ylabel={Abonnés mobile (millions)},
xmin=2000, xmax=2023, ymin=0, ymax=26,
xtick={2000,2005,2010,2015,2020,2023},
ytick={0,5,10,15,20,25},
grid=major, grid style={popline!50},
legend style={at={(0.02,0.98)}, anchor=north west, font=\small, fill=popcream, draw=popline},
tick label style={font=\small},
label style={font=\small}
]
\addplot[popcurve, mark=*, mark size=1.8pt, line width=1.5pt] coordinates {
(2000,0.2) (2002,0.7) (2004,1.6) (2006,3.1) (2008,6.1)
(2010,9.4) (2012,13.1) (2014,16.4) (2016,18.1) (2018,19.7)
(2020,21.6) (2023,23.5)};
\legend{Nombre d'abonnements mobile (cartes SIM actives)}
\end{axis}
\end{tikzpicture}
\end{center}
\legende{Évolution du nombre d'abonnements au téléphone mobile au Cameroun, 2000-2023 (ordres de grandeur, sources : ART, UIT, Banque mondiale). Trois phases : démarrage lent (2000-2004), forte accélération (2004-2014), ralentissement vers la saturation en abonnements (après 2015). L'enjeu actuel est l'internet mobile (3G/4G) et les services financiers numériques.}
\end{popfigure}

\begin{popfigure}
\begin{center}
\begin{tikzpicture}[scale=0.85, every node/.style={transform shape}]
% Europe (bloc simplifié)
\draw[fill=popgreenL, draw=popgreen, thick] (1.5,6.0) -- (5.5,6.0) -- (5.0,4.6) -- (2.0,4.6) -- cycle;
\node[font=\small\bfseries, text=popgreen] at (3.5,5.3) {EUROPE};
% Afrique (bloc simplifié)
\draw[fill=popgoldL, draw=popgold, thick] (1.0,4.2) -- (6.0,4.2) -- (6.4,2.6) -- (5.2,0.2) -- (3.6,0.0) -- (2.2,1.2) -- (1.2,2.6) -- cycle;
\node[font=\small\bfseries, text=popgold] at (3.8,3.4) {AFRIQUE};
% Amérique du Sud (bloc)
\draw[fill=popblueL, draw=popblue, thick] (-4.6,3.2) -- (-2.6,3.2) -- (-2.8,1.4) -- (-3.6,0.2) -- (-4.4,1.0) -- cycle;
\node[font=\small\bfseries, text=popblue] at (-3.6,2.0) {AM. DU SUD};
% Cameroun (point)
\fill[poppink] (2.6,3.9) circle (0.14);
\node[font=\small\bfseries, poppink, anchor=south] at (2.6,4.1) {CAMEROUN};
% câble SAT-3 (côte ouest)
\draw[very thick, popblue, dashed] (2.0,4.6) .. controls (1.0,3.4) and (1.4,1.2) .. (3.4,-0.1);
\node[font=\small, popblue, anchor=west] at (0.2,1.6) {SAT-3 (2002)};
% câble WACS
\draw[very thick, poporange] (2.2,4.6) .. controls (0.4,3.6) and (0.6,1.0) .. (3.2,-0.3);
\node[font=\small, poporange, anchor=east] at (0.3,2.6) {WACS (2012)};
% flèche nord
\draw[-{Stealth}, thick, popdark] (6.8,5.2) -- (6.8,6.2);
\node[font=\small] at (6.8,6.45) {N};
% échelle indicative
\draw[|-|, thick, popdark] (-5.5, -0.5) -- (-3.5, -0.5);
\node[font=\tiny, anchor=north] at (-4.5, -0.55) {$\sim$ 5000 km};
\end{tikzpicture}
\end{center}
\legende{Croquis simplifié des câbles sous-marins à fibre optique desservant l'Afrique de l'Ouest et le Cameroun : SAT-3/WASC (2002) et WACS (2012) longent la côte atlantique et relient l'Afrique à l'Europe. Croquis indicatif, sans précision cartographique.}
\end{popfigure}

\begin{methode}[Analyser une courbe statistique en Histoire / Géographie]
\begin{enumerate}
\item \textbf{Lire} le titre, les axes (grandeurs et unités), l'échelle, la source et la date.
\item \textbf{Dégager les grandes phases} : démarrage, accélération, ralentissement, saturation, éventuelle baisse.
\item \textbf{Repérer les ruptures} : années où la pente change brutalement (points d'inflexion).
\item \textbf{Proposer une explication historique} pour chaque phase (innovation technique, politique publique, baisse des prix, concurrence, contexte économique, équipement concurrent).
\item \textbf{Conclure} en reliant la courbe au contexte général étudié (ici : saut technologique, inclusion numérique).
\end{enumerate}
\end{methode}

\begin{exoresolu}[Analyse de la courbe des abonnements mobiles au Cameroun]
Énoncé : À partir de la courbe ci-dessus, dégage les grandes phases de l'évolution du nombre d'abonnements au téléphone mobile au Cameroun entre 2000 et 2023 et propose une explication historique pour chacune.
\tcblower
Solution détaillée :
\begin{enumerate}
\item \textbf{Phase 1 (2000-2004) : démarrage lent.} Le nombre d'abonnements passe d'environ 0,2 à 1,6 million. \textbf{Explications :} coûts élevés des terminaux et des communications, couverture limitée aux grandes villes (Yaoundé, Douala), faible pouvoir d'achat, monopole initial de Camtel Mobile avant l'arrivée effective de MTN et Orange.
\item \textbf{Phase 2 (2004-2014) : forte accélération (croissance exponentielle).} On passe d'environ 1,6 à 16,4 millions d'abonnements. \textbf{Explications :} guerre des prix entre MTN et Orange, arrivée de téléphones bas coût (Nokia 1100, puis smartphones d'entrée de gamme), extension de la couverture rurale, prépaiement (cartes à gratter) adapté aux revenus irréguliers, absence d'alternative fixe (téléphone filaire quasi inexistant).
\item \textbf{Phase 3 (2015-2023) : ralentissement et saturation en abonnements.} La croissance se tasse (de 18 à environ 23,5 millions) : le taux de pénétration approche 100 \% (multi-SIM fréquent). \textbf{Explications :} la plupart des adultes possèdent déjà une (ou plusieurs) carte(s) SIM. L'enjeu bascule vers la data (internet mobile 3G/4G), la qualité de service et les services à valeur ajoutée (Mobile Money, e-gouvernement, e-commerce).
\end{enumerate}
Conclusion : La courbe illustre le « saut technologique » du Cameroun, passé en vingt ans d'une quasi-absence de télécommunications modernes à une société du mobile, base de l'inclusion financière et numérique.
\end{exoresolu}

\begin{savaistu}[Le Mobile Money, une innovation africaine adoptée par le monde]
Le \cle{Mobile Money} permet à des millions de Camerounais sans compte bancaire (exclus du système classique) de payer, épargner, emprunter (micro-crédit) et transférer de l'argent par simple téléphone, via MTN MoMo ou Orange Money. Cette innovation de \emph{rupture} a été adoptée \emph{plus vite en Afrique qu'en Europe} : l'Afrique de l'Est (Kenya, M-Pesa dès 2007) et l'Afrique centrale sont devenues des pionnières mondiales du paiement mobile. Le numérique n'est donc pas seulement importé : il est aussi réinventé et adapté aux réalités locales sur le continent.
\end{savaistu}

\begin{exercice}[Pour s'entraîner — type Probatoire]
\begin{enumerate}
\item Reconstitue la frise chronologique ordonnée des innovations suivantes : iPhone, ENIAC, transistor, Arpanet, Web, microprocesseur, IBM PC.
\item Explique la différence entre internet et le Web en deux phrases précises.
\item Pourquoi peut-on dire que le téléphone mobile a permis au Cameroun un « saut technologique » ? Donne deux arguments.
\item \textbf{Analyse de document :} En t'appuyant sur la courbe des abonnements mobiles (figure 1) et sur la méthode, décris l'évolution entre 2000 et 2023 et explique la phase d'accélération (2004-2014).
\end{enumerate}
\end{exercice}

\begin{corrige}[Exercice]
\begin{enumerate}
\item ENIAC (1945) ; transistor (1947) ; Arpanet (1969) ; microprocesseur (1971) ; IBM PC (1981) ; Web (1989-1990) ; iPhone (2007).
\item Internet est le réseau mondial d'infrastructures qui interconnecte des millions d'ordinateurs via les protocoles TCP/IP. Le Web (World Wide Web) est un service d'hypertexte (pages HTML liées par des URL) accessible via un navigateur, qui fonctionne \emph{au-dessus} d'internet.
\item 1) Le Cameroun, où le réseau téléphonique fixe est resté très peu développé (coût d'infrastructure, géographie), a comblé en quelques années un retard d'équipement vieux de plusieurs décennies en passant directement au mobile. 2) Le mobile a ensuite servi de plateforme à l'internet mobile et au Mobile Money, créant de nouveaux services (banque, santé, éducation) inaccessibles auparavant.
\item \textbf{Description :} Trois phases : lent démarrage (0,2 à 1,6 M, 2000-2004), forte accélération (1,6 à 16,4 M, 2004-2014), ralentissement vers saturation (~23,5 M, 2015-2023). \textbf{Explication accélération :} Concurrence MTN/Orange $\rightarrow$ baisse prix ; téléphones bas coût ; prépaiement ; couverture rurale ; absence ligne fixe.
\end{enumerate}
\end{corrige}

\begin{aretenir}[L'essentiel de la section — Fiche de révision Bac]
\begin{itemize}
\item \textbf{Révolution informatique :} Miniaturisation cumulative (Transistor 1947 $\rightarrow$ Circuit intégré 1958 $\rightarrow$ Microprocesseur 1971) $\rightarrow$ Ordinateur personnel (Apple 1977, IBM PC 1981) $\rightarrow$ Smartphone (iPhone 2007). Loi de Moore.
\item \textbf{Internet / Web :} Arpanet 1969 (militaire/universitaire) $\rightarrow$ TCP/IP 1983 (standard) $\rightarrow$ Web 1989-90 (Tim Berners-Lee, CERN, bien public) $\rightarrow$ Mondialisation années 1990-2000.
\item \textbf{Mobile :} 1G (analogique) $\rightarrow$ 2G GSM/SMS (années 1990) $\rightarrow$ 3G/4G (internet mobile) $\rightarrow$ Smartphone 2007 $\rightarrow$ 5G/IoT (années 2020).
\item \textbf{Cameroun :} Internet 1997 (satellite) $\rightarrow$ Mobile MTN/Orange (2000+) $\rightarrow$ Mobile Money (inclusion financière) $\rightarrow$ Fibre optique SAT-3 (2002), WACS (2012), dorsales nationales. Saut technologique (leapfrogging).
\item \textbf{Enjeux actuels :} IA, Big Data, Robotique, Objets connectés $\leftrightarrow$ Protection données, Fracture numérique, Environnement, Éthique.
\end{itemize}
\end{aretenir}

\coursec{Les grandes innovations techniques : transports, énergie et agriculture}

\courssub{Transition : du numérique au monde physique}

Après avoir analysé la révolution numérique — ordinateurs, internet, télécommunications — qui a transformé le traitement et la circulation de l'information, nous nous tournons maintenant vers les innovations qui ont bouleversé la circulation des \textbf{êtres humains}, des \textbf{marchandises} et de l'\textbf{énergie}, ainsi que la production \textbf{alimentaire}. Ces domaines — transports, énergie, agriculture — constituent le socle matériel de la mondialisation et de la croissance économique depuis 1945. Ils illustrent la puissance de l'ingénierie moderne, mais posent aussi, avec une acuité croissante, la question des limites planétaires et des inégalités de développement, notamment au Cameroun.

\courssub{Les transports : compression de l'espace-temps}

\begin{definition}[La révolution des transports]
L'ensemble des innovations techniques (avion à réaction, train à grande vitesse, conteneurisation, automobile de masse) qui, depuis les années 1950, ont drastiquement réduit les coûts et les durées de déplacement des personnes et des marchandises, fondant l'intégration économique mondiale.
\end{definition}

\textbf{L'aviation de masse et l'avion à réaction.}\par
Le passage de l'hélice à la \textbf{turbine à réaction} (premier vol commercial : De Havilland Comet, 1952 ; mise en service régulière Boeing 707, 1958) triple la vitesse de croisière (environ \SI{900}{km/h}) et abaisse le coût au kilomètre-passager. L'apogée symbolique est le \textbf{Boeing 747 « Jumbo Jet » (1969, premier vol ; 1970, mise en service)}, premier gros-porteur à deux couloirs, capable d'emporter 400 à 500 passagers. Aujourd'hui, l'Airbus A380 (2007) pousse la capacité au-delà de 800 passagers en configuration haute densité.

\begin{exemplebox}[Chiffres clés de l'aviation civile]
\begin{itemize}
    \item Trafic mondial (passagers-km) : multiplié par \textbf{plus de 70} entre 1960 et 2019 (avant Covid-19).
    \item Part du transport aérien dans le commerce mondial en valeur : \textbf{35~\%} (pour moins de 1~\% en tonnage).
    \item Un Paris–New York : \SI{7}{h} \SI{30}{min} aujourd'hui contre \SI{14}{h} en 1958 (Super Constellation).
\end{itemize}
\end{exemplebox}

\textbf{Le train à grande vitesse (TGV) : la réponse ferroviaire.}\par
Face à l'avion sur les moyennes distances (\SIrange{300}{800}{km}), le Japon lance le \textbf{Shinkansen} (1964, Tokyo–Osaka, \SI{210}{km/h}). La France suit avec le \textbf{TGV} (1981, Paris–Lyon, \SI{270}{km/h} puis \SI{320}{km/h}). Le TGV repose sur : lignes dédiées (grands rayons, faible pente), alimentation \SI{25}{kV} \SI{50}{Hz}, motricité répartie, freinage rhéostatique/régénératif. Il capte 90~\% du trafic Paris–Lyon et rend le train compétitif porte-à-porte face à l'avion jusqu'à \SI{3}{h} de trajet.

\textbf{La conteneurisation : le standard invisible de la mondialisation.}\par
L'idée de Malcom McLean (1956, navire \emph{Ideal X}) : charger des \textbf{conteneurs} normalisés (ISO 668 : 20 pieds = \SI{6,1}{m} ; 40 pieds = \SI{12,2}{m}) directement du camion/wagon au navire sans rupture de charge. Conséquences :
\begin{itemize}
    \item Division par 20 du coût de manutention portuaire.
    \item Division par 10 du temps d'escale.
    \item Explosion du trafic maritime : \textbf{860 millions d'EVP (équivalents 20 pieds) en 2023} contre 50 millions en 1980.
    \item Navires géants : \emph{MSC Irina} (2023), 24~346~EVP, \SI{400}{m} de long.
\end{itemize}

\begin{popfigure}
\begin{tikzpicture}[scale=0.9]
% Axes et titre
\node[align=center, font=\bfseries\large] at (0,5.5) {Évolution du trafic mondial de conteneurs (EVP millions)};
\begin{axis}[
    width=\linewidth, height=8cm,
    xlabel={Année}, ylabel={EVP (millions)},
    xmin=1980, xmax=2023, ymin=0, ymax=900,
    xtick={1980,1990,2000,2010,2020},
    ytick={0,200,400,600,800},
    grid=major, grid style={popline},
    axis lines=left, enlarge x limits=0.02,
    popcurve, thick,
]
\addplot coordinates {
    (1980,50) (1990,100) (2000,225) (2010,470) (2020,830) (2022,830) (2023,860)
};
\addplot[mark=*, poporange] coordinates {(2023,860)} node[above right, font=\small] {2023 : 860 M EVP};
\end{axis}
\end{tikzpicture}
\legende{Évolution du trafic maritime conteneurisé (EVP). Source : UNCTAD / Clarkson Research. La courbe illustre l'emballement de la mondialisation logistique.}
\end{popfigure}

\textbf{L'automobile : de la production de masse à l'électrification.}\par
Le modèle fordiste (chaîne mobile, 1913) s'étend mondialement après 1945. Le \textbf{toyotisme} (années 1970, production « juste-à-temps », qualité totale) généralise la voiture comme bien de consommation courant. Production mondiale : 40~millions (1990) $\to$ 93~millions (2017) $\to$ environ 90~millions (2023, post-Covid).
Depuis 2010, rupture avec le \textbf{véhicule électrique (VE)} et \textbf{hybride} : batteries Li-ion (densité énergétique $\approx$ \SI{250}{Wh/kg}), moteurs synchrones à aimants permanents, freinage régénératif. Part des VE dans les ventes mondiales : 14~\% (2022), 18~\% (2023). Enjeux : matières critiques (lithium, cobalt, terres rares), recyclage, mix électrique pour la recharge.

\begin{attention}[Piège : « voiture zéro émission »]
La voiture électrique n'émet pas de CO$_2$ à l'échappement, mais son \textbf{bilan carbone complet} (cycle de vie : extraction, fabrication batterie, électricité de recharge, fin de vie) dépend du mix énergétique national. En Pologne (charbon), le VE émet plus qu'un thermique efficient sur \SI{150000}{km} ; en France (nucléaire/hydro), l'équilibre est atteint vers \SI{30000}{km}. Ne jamais écrire « zéro émission » sans préciser le périmètre.
\end{attention}

\courssub{Les énergies : de l'âge du pétrole à la transition bas-carbone}

\begin{definition}[Mix énergétique]
Répartition des sources d'énergie primaire (pétrole, gaz, charbon, nucléaire, renouvelables) dans la consommation totale d'un pays ou du monde. Il évolue sous l'effet des prix, des technologies, des géopolitiques et des impératifs climatiques.
\end{definition}

\textbf{L'hégémonie des fossiles (pétrole, gaz, charbon).}\par
1945–1973 : « Trente Glorieuses » portées par le pétrole bon marché (Moyen-Orient, prix < 3 \$/baril). Le pétrole devient la 1\textsuperscript{re} énergie mondiale (1964), devant le charbon. Chocs pétroliers 1973/1979 $\to$ diversification (nucléaire, gaz, efficacité). Aujourd'hui encore, les fossiles représentent \textbf{82~\%} de l'énergie primaire mondiale (2022).

\textbf{Le nucléaire civil : la fission industrielle.}\par
Premier réacteur producteur d'électricité : Obninsk (URSS, 1954), Calder Hall (Royaume-Uni, 1956). Principe : fission de l'U-235 $\to$ chaleur $\to$ vapeur $\to$ turbine. Générations : I (prototypes), II (REP/REB, gros du parc actuel), III (EPR, AP1000, sécurité passive), IV (réacteurs à neutrons rapides, thorium, horizon 2040). Avantages : bas carbone, pilotable, densité énergétique énorme (1~g U-235 $\approx$ 1~t pétrole). Inconvénients : déchets radioactifs à vie longue, risque d'accident majeur (Tchernobyl 1986, Fukushima 2011), coût capital élevé, prolifération.

\textbf{L'essor des énergies renouvelables (EnR).}\par
Poussée par le réchauffement climatique (GIEC, Accord de Paris 2015) et la chute des coûts :
\begin{itemize}
    \item \textbf{Photovoltaïque :} coût divisé par 10 depuis 2010 (LCOE $\approx$ \SI{0,05}{\$/kWh} en 2023). Puissance installée : \SI{1,4}{TW} (2023).
    \item \textbf{Éolien :} onshore mature, offshore en forte croissance (facteur de charge 40–50~\%). Puissance : \SI{1}{TW} (2023).
    \item \textbf{Hydroélectricité :} \SI{1360}{GW} (2023), 15~\% électricité mondiale, première EnR, stockable (lacs de barrage, STEP).
\end{itemize}
Défis : intermittence $\to$ besoins de stockage (batteries, hydrogène vert, STEP), réseaux, acceptabilité, criticité des métaux (cuivre, lithium, terres rares).

\begin{popfigure}
\begin{tikzpicture}
\begin{axis}[
    width=\linewidth, height=9cm,
    xbar, xlabel={Part (\%)}, ylabel={Source},
    symbolic y coords={Autres EnR, Hydroélectricité, Nucléaire, Gaz naturel, Charbon, Pétrole},
    ytick=data,
    nodes near coords={\pgfmathprintnumber\pgfplotspointmeta\,\%},
    nodes near coords align={horizontal},
    xmin=0, xmax=35,
    grid=major, grid style={popline},
    axis lines=left, enlarge y limits=0.2,
    bar width=10pt,
]
\addplot[fill=popblue!70] coordinates {
    (14,Autres EnR) (4,Hydroélectricité) (4,Nucléaire) (23,Gaz naturel) (24,Charbon) (31,Pétrole)
};
\end{axis}
\end{tikzpicture}
\legende{Répartition de la consommation d'énergie primaire mondiale en 2022. Source : Energy Institute Statistical Review 2023. Les fossiles dominent encore à 82~\%.}
\end{popfigure}

\courssub{La révolution verte en agriculture : nourrir 8 milliards d'humains}

\begin{definition}[Révolution verte]
Ensemble des techniques agricoles modernes — variétés à haut rendement (semi-naines), engrais chimiques (NPK), irrigation maîtrisée, mécanisation, phytosanitaires — diffusées depuis les années 1960 (Norman Borlaug, CIMMYT, IRRI) qui ont multiplié les rendements céréaliers (blé, riz, maïs) et évité des famines massives.
\end{definition}

\textbf{Les piliers techniques.}\par
\begin{itemize}
    \item \textbf{Semences :} gènes « nains » (Rht1 chez le blé, sd1 chez le riz) $\to$ tiges courtes, résistantes à la verse, forte réponse à l'azote. Rendements : blé Mexique \SI{1}{t/ha} (1950) $\to$ \SIrange{5}{7}{t/ha} (aujourd'hui) ; riz IR8 « riz miracle » \SIrange{5}{10}{t/ha} contre \SI{2}{t/ha} variétés traditionnelles.
    \item \textbf{Engrais :} synthèse Haber-Bosch (NH$_3$ à partir de N$_2$ + H$_2$). Consommation mondiale N : \SI{11}{Mt} (1960) $\to$ \SI{115}{Mt} (2022). Problème : excès $\to$ eutrophisation, N$_2$O (gaz à effet de serre puissant).
    \item \textbf{Irrigation :} surfaces équipées $\times$ 3 depuis 1960 (\SI{340}{Mha}). Pompe à motopompe, goutte-à-goutte (efficacité 90~\% vs 50~\% gravitaire).
    \item \textbf{Mécanisation :} tracteurs, moissonneuses-batteuses, drones de pulvérisation. Puissance agricole mondiale : 28~millions de tracteurs (2020).
\end{itemize}

\begin{exoresolu}[Bilan chiffré de la révolution verte]
\textbf{Énoncé :} À partir des données ci-dessous, calculez l'évolution de la production céréalière mondiale et du rendement moyen entre 1961 et 2022. Commentez le rôle de la surface vs le rendement.

\begin{center}
\begin{tabularx}{\linewidth}{|l|X|X|}\hline
Indicateur & 1961 & 2022 \\ \hline
Production céréalière (Mt) & 877 & 2~900 \\ \hline
Surface récoltée (Mha) & 650 & 730 \\ \hline
Rendement moyen (t/ha) & 1,35 & 4,0 \\ \hline
Population mondiale (Md) & 3,1 & 8,0 \\ \hline
\end{tabularx}
\end{center}

\tcblower
\textbf{Solution :}
\begin{itemize}
    \item Production $\times$ 3,3 (877 $\to$ 2~900~Mt).
    \item Surface $\times$ 1,12 seulement (650 $\to$ 730~Mha) : l'essentiel de la croissance vient du \textbf{rendement} ($\times$ 3,0).
    \item Production par habitant : \SI{283}{kg} (1961) $\to$ \SI{363}{kg} (2022) : malgré la démographie, la disponibilité alimentaire moyenne a augmenté.
    \item \textbf{Conclusion :} La révolution verte a « déconnecté » la production alimentaire de l'extension des terres cultivées (épargne des forêts), mais au prix d'une forte dépendance aux intrants fossiles (engrais, carburant, pesticides).
\end{itemize}
\end{exoresolu}

\textbf{Limites et externalités négatives.}\par
Érosion des sols, perte de biodiversité (monocultures), pollution nitrates/pesticides (algues vertes, effondrement insectes), épuisement nappes phréatiques (Inde, Chine, Ogallala USA), émission GES (agriculture = 22~\% GES mondiaux avec changement d'affectation des sols). Nécessité d'une « révolution verte 2.0 » : agroécologie, agriculture de précision, variétés tolérantes à la sécheresse (CRISPR), fixation biologique de l'azote.

\courssub{Matériaux nouveaux : plastiques, composites, nanomatériaux}

\begin{definition}[Plastiques et polymères]
Macromolécules organiques synthétiques (PE, PP, PVC, PET, PS, PUR) issues de la pétrochimie (naphta, éthylène, propylène). Production : \SI{2}{Mt} (1950) $\to$ \SI{400}{Mt} (2022). Propriétés : légers, moulables, isolants, peu coûteux, durables (trop ?).
\end{definition}

\begin{itemize}
    \item \textbf{Composites :} matrice (résine époxy, polyester) + renfort (fibres de carbone, verre, aramide). Rapport résistance/masse inégalé : aéronautique (A350, B787 : 50~\% masse en composite), éoliennes, automobile, sport.
    \item \textbf{Nanomatériaux :} au moins une dimension $<$ \SI{100}{nm}. Propriétés quantiques/de surface : nanotubes de carbone (conductivité, résistance), graphène, nanoparticules (TiO$_2$ crème solaire, Ag antibactérien, catalyseurs). Enjeux : toxicologie, écotoxicité, régulation (REACH).
\end{itemize}

\begin{attention}[Pollution plastique : un marqueur géologique]
8 à 12~Mt/an rejoignent les océans. Fragmentation en \textbf{microplastiques} ($<$ \SI{5}{mm}) et \textbf{nanoplastiques} retrouvés dans l'air, l'eau, la chaîne alimentaire, le sang humain, le placenta. Le plastique est devenu un \textbf{marqueur stratigraphique de l'Anthropocène}. Le traité international contre la pollution plastique (négociations ONU 2022–2024) vise un cycle de vie complet : réduction, réemploi, recyclage chimique, interdiction produits à usage unique.
\end{attention}

\courssub{Applications au Cameroun : potentialités et réalités}

Le Cameroun, « Afrique en miniature », illustre le contraste entre dotations naturelles exceptionnelles et infrastructures encore insuffisantes.

\textbf{Énergie hydroélectrique : le « château d'eau » sous-exploité.}\par
Potentiel théorique : \textbf{\SI{23}{GW}} (2\textsuperscript{e} Afrique après RDC), potentiel économiquement exploitable $\approx$ \SI{12}{GW}. Production installée 2023 : $\approx$ \SI{1,5}{GW} (essentiellement hydro). Déficit chronique : délestages, coût élevé de l'électricité industrielle (frein à l'industrialisation).

\begin{popfigure}
\begin{tikzpicture}[scale=0.85]
% Contour très simplifié du Cameroun (polygone approximatif)
\draw[popdark, thick, fill=popcream] 
    (0,0) -- (1,0.2) -- (2.5,0.5) -- (3.5,1.5) -- (4,3) -- (3.8,5) -- (3,6.5) -- (2,7) -- (1,6.5) -- (0.2,5) -- (-0.5,3.5) -- (-0.8,2) -- (-0.5,0.5) -- cycle;
% Villes repères
\node[circle, fill=popblue, inner sep=1.5pt, label=above right:\tiny Yaoundé] (Y) at (1.5,3.5) {};
\node[circle, fill=popblue, inner sep=1.5pt, label=above left:\tiny Douala] (D) at (1.2,2.2) {};
\node[circle, fill=popblue, inner sep=1.5pt, label=above:\tiny Garoua] (G) at (2.8,5.5) {};
\node[circle, fill=popblue, inner sep=1.5pt, label=left:\tiny Bertoua] (B) at (3.2,4) {};
\node[circle, fill=popblue, inner sep=1.5pt, label=below:\tiny Ebolowa] (E) at (0.5,1.5) {};
% Barrages existants
\node[star, star points=5, star point ratio=2.25, fill=popblue, inner sep=2.5pt, label=above right:\textbf{Édéa (263 MW)}] at (0.8,1.8) {};
\node[star, star points=5, star point ratio=2.25, fill=popblue, inner sep=2.5pt, label=right:\textbf{Songloulou (388 MW)}] at (1.0,1.3) {};
\node[star, star points=5, star point ratio=2.25, fill=poporange, inner sep=2.5pt, label=above:\textbf{Lom Pangar (30 MW + régulation)}] at (2.8,4.2) {};
\node[star, star points=5, star point ratio=2.25, fill=poporange, inner sep=2.5pt, label=below left:\textbf{Memve'ele (211 MW)}] at (0.3,1.0) {};
\node[star, star points=5, star point ratio=2.25, fill=popgreen, inner sep=2.5pt, label=above right:\textbf{Nachtigal (420 MW, 2023-24)}] at (1.8,3.0) {};
% Barrages projetés / en étude
\node[diamond, fill=poppink, inner sep=2.5pt, label=above:\textbf{Grand Eweng (800-1200 MW)}] at (0.2,2.5) {};
\node[diamond, fill=poppink, inner sep=2.5pt, label=right:\textbf{Kikot (500 MW)}] at (1.2,2.8) {};
\node[diamond, fill=poppink, inner sep=2.5pt, label=above:\textbf{Chollet (600 MW, frontière Congo)}] at (-0.2,1.5) {};
% Fleuves
\draw[popblue, thick, decorate, decoration={snake, amplitude=1pt, segment length=5pt}] (0.8,1.8) -- (0.5,0.5) node[below, font=\tiny] {Sanaga};
\draw[popblue, thick, decorate, decoration={snake, amplitude=1pt, segment length=5pt}] (2.8,4.2) -- (1.8,3.0) node[left, font=\tiny] {Lom / Sanaga};
% Légende
\node[align=left, font=\small, fill=white, fill opacity=0.9, text opacity=1, rounded corners, inner sep=3pt] at (5.5,3) {
    \textbf{Légende} \\
    \textcolor{popblue}{$\star$} Barrages historiques (Édéa 1955, Songloulou 1980) \\
    \textcolor{poporange}{$\star$} Barrages récents (Lom Pangar 2021, Memve'ele 2019) \\
    \textcolor{popgreen}{$\star$} Nachtigal (2023-24, +30\% capacité) \\
    \textcolor{poppink}{$\diamond$} Projets structurants (Grand Eweng, Kikot, Chollet) \\
    \textcolor{popblue}{\textasciitilde} Fleuve Sanaga (bassin 133 000 km$^2$)
};
\node[align=center, font=\bfseries\large] at (1.5,7.2) {Principaux barrages hydroélectriques du Cameroun};
\draw[->, thick] (5.8,6.5) -- (5.8,6) node[below, font=\small] {N};
\end{tikzpicture}
\legende{Localisation schématique des aménagements hydroélectriques sur le bassin de la Sanaga et le Sud. Le Cameroun n'exploite encore qu'environ 10~\% de son potentiel hydroéconomique.}
\end{popfigure}

\textbf{Transports : routes, rails, ports.}\par
\begin{itemize}
    \item \textbf{Routes :} \SI{50000}{km} dont \SI{6000}{km} bitumés (2023). Corridors structurants : Douala–Yaoundé (2x2 voies en cours), Douala–Bafoussam, Yaoundé–Bertoua–Ngaoundéré. Ponts : Wouri (2e pont 2017, \SI{750}{m}), Logone, Mbam.
    \item \textbf{Ferroviaire :} Transcamerounais (Douala–Ngaoundéré, \SI{886}{km}, voie unique, métrique, vétuste). Projet « Ferroviaire Nord » (Ngaoundéré–Maroua, \SI{400}{km}) et liaison Kribi–Mbalam (minerai de fer).
    \item \textbf{Ports :} Douala (95~\% trafic, saturation), Kribi (eau profonde, 2018, vrac, conteneurs, gaz). Objectif : hub sous-régional (CEMAC).
\end{itemize}

\textbf{Agriculture : mécanisation naissante, défis structurels.}\par
\begin{itemize}
    \item 70~\% population active, 20~\% PIB, mais productivité faible (rendements maïs \SI{1,5}{t/ha} vs \SIrange{5}{10}{t/ha} potentiel).
    \item Mécanisation : 0,5 tracteur/1000 ha (vs 10 en Afrique du Sud, 200 en France). Essor récent : SODECOTON (coton, Nord), CDC (palmier, bananier, Sud-Ouest), fermes pilotes (Mbé, Nkoteng).
    \item Intrants : engrais subventionnés (programme PEA-Jeunes, subventions 50~\%), semences améliorées (IRAD : maïs CMS, riz NERICA, manioc TMS).
    \item Enjeux : transformation locale (cacao, café, coton, bois), chaînes de froid, accès au financement, adaptation climatique (variétés courtes, agroforesterie).
\end{itemize}

\begin{savaistu}[Le paradoxe camerounais]
Le Cameroun dispose du \textbf{deuxième potentiel hydroélectrique d'Afrique} (après la RDC) et de terres arables parmi les plus fertiles du continent. Pourtant, en 2023, \textbf{40~\% de la population n'a pas accès à l'électricité} (monde rural < 20~\%) et le pays importe encore du riz, du blé, du poisson et de l'huile pour plusieurs centaines de milliards de FCFA par an. Le progrès technique (barrages, semences, routes) ne se diffuse pas automatiquement : il nécessite \textbf{institutions efficaces, financement long terme, formation technique, gouvernance transparente} pour transformer le potentiel en bien-être partagé. C'est tout l'enjeu de la \textbf{Stratégie Nationale de Développement 2020-2030 (SND30)}.
\end{savaistu}

\courssub{Bilan : les effets des révolutions scientifiques et techniques sur le monde et le Cameroun}

\begin{aretenir}[Synthèse des RST (1945–...)]
\begin{itemize}
    \item \textbf{Transports :} compression espace-temps (avion, TGV, conteneur) $\to$ mondialisation physique, chaînes de valeur globales, tourisme de masse.
    \item \textbf{Énergie :} bascule fossile $\to$ bas-carbone (nucléaire, EnR) sous contrainte climatique. Mix 2022 : 82~\% fossiles. Défi : stockage, réseaux, métaux critiques.
    \item \textbf{Agriculture :} Révolution verte $\to$ sécurité alimentaire pour 8~Md, mais externalités (sols, eau, biodiversité, GES). Nécessité transition agroécologique.
    \item \textbf{Matériaux :} plastiques, composites, nano $\to$ performance, légèreté, mais pollution persistante (plastiques) et incertitudes sanitaires (nano).
    \item \textbf{Cameroun :} potentiels énormes (hydro, agro, mines), infrastructures en retard, électricité/accès inégaux. Le développement technique est une \textbf{condition nécessaire mais non suffisante} : il faut des politiques publiques inclusives et une appropriation locale des savoirs.
\end{itemize}
\end{aretenir}

\begin{methode}[Commenter un document sur les innovations techniques (Bac)]
\begin{enumerate}
    \item \textbf{Identifier} : nature (texte, graphique, carte, photo), auteur, date, source, thème (transport, énergie, agriculture...).
    \item \textbf{Contextualiser} : rattacher à la période (Trente Glorieuses, chocs pétroliers, transition écologique, émergence Cameroun).
    \item \textbf{Analyser le contenu} : extraire 2–3 idées forces (chiffres, tendances, acteurs, innovations citées).
    \item \textbf{Expliquer / Interpréter} : mobiliser les connaissances du cours (causes : R\&D, demande, État ; conséquences : croissance, inégalités, environnement).
    \item \textbf{Conclure / Ouvrir} : porter un jugement nuancé (progrès vs limites, mondial vs local, court vs long terme) et lier à l'actualité (COP, SND30, ZLECAf).
\end{enumerate}
\end{methode}

\begin{exercice}[Entraînement Bac — Document statistique]
\textbf{Document :} Évolution de la part des énergies renouvelables dans la production d'électricité mondiale (\%).
\begin{center}
\begin{tabular}{|c|c|c|c|c|c|c|}\hline
Année & 1990 & 2000 & 2010 & 2015 & 2020 & 2023 \\ \hline
Part EnR (\% élec.) & 19,5 & 19,0 & 20,0 & 23,5 & 29,0 & 30,3 \\ \hline
\end{tabular}
\end{center}
\textbf{Questions :}
\begin{enumerate}
    \item Représentez cette évolution sous forme de graphique adapté (choix justifié).
    \item Calculez le taux de variation moyen annuel de la part des EnR sur 1990–2023.
    \item En vous appuyant sur le cours, expliquez la stagnation 1990–2000 puis l'accélération post-2010.
    \item Au Cameroun, la part de l'hydroélectricité dans le mix \emph{électrique} est d'environ 60~\%. Pourquoi cette part est-elle plus élevée que la moyenne mondiale, et quels risques cela comporte-t-il ?
\end{enumerate}
\end{exercice}

\begin{corrige}[Exercice n°1]
\begin{enumerate}
    \item \textbf{Graphique :} courbe en lignes (évolution temporelle continue) ou histogramme (comparaison dates discrètes). Axes : abscisse = années, ordonnée = \% (échelle 0–35~\%). Titre : « Part des EnR dans la production électrique mondiale (1990–2023) ». Source indiquée.
    \item \textbf{Taux variation moyen annuel (TVMA) :} $TVMA = \left(\frac{\{30,3\}}{\{19,5\}}\right)^{\frac{1}{33}} - 1 \approx 1,0146 - 1 = 0,0146$ soit \textbf{1,46~\% par an}. (Faible car base hydro déjà haute en 1990).
    \item \textbf{Explication :} 1990–2000 : hydroélectricité mature, peu de nouveaux grands barrages (contestations sociales/environnementales), solaire/éolien naissants et chers. Post-2010 : chute coûts PV/éolien (économies d'échelle, innovation, politiques climatiques — Kyoto, Paris), objectifs nationaux contraignants, compétitivité LCOE.
    \item \textbf{Cameroun :} Hydro = 60~\% électricité car géologie favorable (Sanaga, chutes), barrages historiques (Édéa, Songloulou), faible demande industrielle historique. \textbf{Risques :} dépendance climatique (sécheresses 2021–2022 $\to$ délestages sévères), concentration géographique (bassin Sanaga), évaporation réservoirs, sédimentation, impact biodiversité. Nécessité diversification (solaire Nord, biomasse, gaz, interconnexion CEMAC).
\end{enumerate}
\end{corrige}

\coursec{Bilan : les effets des RST sur le monde et le Cameroun}

\courssub{Transition : des innovations aux conséquences globales}

Après avoir examiné les grandes innovations techniques dans les transports, l'énergie et l'agriculture, il est temps de dresser le bilan global des révolutions scientifiques et techniques (RST) depuis 1945. Ces révolutions ne sont pas de simples successions d'inventions ; elles ont restructuré en profondeur les sociétés, les économies, les rapports de force géopolitiques et le quotidien de milliards d'êtres humains. Le Cameroun, comme l'ensemble du continent africain, se trouve au cœur de ces transformations, tiraillé entre opportunités inédites et défis structurels persistants.

\courssub{Les effets positifs : progrès sanitaire, alimentaire et communicationnel}

\begin{aretenir}[Les grands acquis des RST pour l'humanité]
Depuis 1945, les RST ont permis :
\begin{itemize}
    \item \textbf{L'allongement spectaculaire de l'espérance de vie} : passage d'une moyenne mondiale d'environ 46 ans en 1950 à plus de 73 ans aujourd'hui, grâce aux vaccins, aux antibiotiques, à l'imagerie médicale et à la chirurgie moderne.
    \item \textbf{Le recul des famines} : la Révolution verte (variétés à haut rendement, engrais, irrigation) a permis de nourrir une population passée de 2,5 milliards d'habitants en 1950 à plus de 8 milliards aujourd'hui, même si la faim persiste en raison des inégalités de distribution et des conflits.
    \item \textbf{L'accélération des communications} : internet, téléphonie mobile et satellites ont aboli les distances, permettant la circulation quasi instantanée de l'information, des capitaux et des savoirs.
    \item \textbf{La mondialisation des échanges} : conteneurisation, transport aérien et logistique informatisée ont intégré les marchés mondiaux, favorisant la croissance économique (notamment en Asie de l'Est).
\end{itemize}
\end{aretenir}

\begin{exemplebox}[Le Cameroun face aux acquis sanitaires]
Au Cameroun, la vaccination a contribué à l'éradication mondiale de la variole (certifiée en 1980 par l'OMS) et a fait reculer drastiquement la poliomyélite, le tétanos néonatal et la rougeole. L'espérance de vie à la naissance est passée d'environ 40 ans en 1960 à près de 61 ans en 2022 (Banque mondiale). La télémédecine commence à relier les zones rurales aux hôpitaux de référence (Yaoundé, Douala) via des applications mobiles, illustrant la convergence entre santé et TIC.
\end{exemplebox}

\begin{popfigure}
\begin{center}
\begin{tabularx}{\linewidth}{|l|X|X|}
\hline
\rowcolor{popblueL} \textbf{Domaine} & \textbf{Bienfaits majeurs} & \textbf{Limites / Risques} \\
\hline
\textbf{Santé} & Vaccins, antibiotiques, chirurgie, imagerie (IRM, scanner), espérance de vie en hausse & Résistance aux antibiotiques, inégalités d'accès (Nord/Sud, urbain/rural), coût des traitements innovants, questions éthiques (édition du génome) \\
\hline
\textbf{Communication} & Internet, téléphonie mobile, accès au savoir (encyclopédies en ligne, MOOC), télétravail, e-administration & Fracture numérique, désinformation, cybercriminalité, addiction aux écrans, surveillance de masse \\
\hline
\textbf{Énergie} & Électrification massive, nucléaire (faible en CO$_2$), énergies renouvelables (solaire, éolien), batteries & Réchauffement climatique (énergies fossiles), déchets nucléaires, dépendance aux métaux rares, intermittence des renouvelables \\
\hline
\textbf{Agriculture} & Révolution verte (rendements en hausse), irrigation, mécanisation, agriculture de précision (drones, capteurs) & Érosion des sols, pollution par les nitrates et pesticides, perte de biodiversité, dépendance aux semenciers, gaspillage alimentaire \\
\hline
\end{tabularx}
\end{center}
\legende{Tableau de synthèse : bienfaits et limites des RST par domaine stratégique. Source : élaboration à partir de données ONU, OMS, AIEA et FAO.}
\end{popfigure}

\courssub{Les effets négatifs : menaces existentielles et externalités}

\begin{definition}[Menace nucléaire]
La menace nucléaire désigne le risque d'utilisation militaire (guerre, terrorisme) ou civile (accident industriel) de l'énergie atomique, aux conséquences humanitaires et environnementales durables à l'échelle planétaire (ex. Hiroshima et Nagasaki en 1945, Tchernobyl en 1986, Fukushima en 2011).
\end{definition}

\begin{aretenir}[Les quatre grands risques systémiques des RST]
\begin{enumerate}
    \item \textbf{Menace nucléaire} : prolifération (neuf États disposent aujourd'hui de l'arme atomique), risque terroriste, déchets radioactifs à vie longue (plusieurs centaines de milliers d'années pour certains isotopes).
    \item \textbf{Pollution et réchauffement climatique} : l'industrialisation fondée sur les énergies fossiles (charbon, pétrole, gaz) a émis environ 2 400 milliards de tonnes de CO$_2$ depuis 1850, dont plus de la moitié depuis 1990. Conséquences : élévation du niveau des mers, multiplication des événements climatiques extrêmes, effondrement de la biodiversité (les scientifiques évoquent une sixième extinction de masse).
    \item \textbf{Obsolescence et déchets électroniques (DEEE)} : environ 62 millions de tonnes de DEEE générées en 2022 (ONU), dont moins d'un quart est recyclé formellement. Une partie est exportée illégalement vers l'Afrique et l'Asie (Agbogbloshie au Ghana, Guiyu en Chine) : pollution aux métaux lourds (plomb, mercure), travail des enfants.
    \item \textbf{Manipulation de l'information} : algorithmes de recommandation optimisant l'engagement par la colère ou la peur, hypertrucages (\emph{deepfakes}) permis par l'IA générative, ingérences électorales, érosion de la confiance démocratique.
\end{enumerate}
\end{aretenir}

\begin{attention}[Piège fréquent : le déterminisme technologique]
\emph{Erreur} : « La technique résout tous les problèmes » ou « La technique est neutre ».
\emph{Nuance} : toute technique porte en elle des valeurs (efficacité, profit, contrôle) et produit des externalités. Le nucléaire civil fournit de l'électricité faible en CO$_2$ \textbf{mais} génère des déchets difficiles à gérer. L'IA optimise la logistique \textbf{mais} peut reproduire des discriminations dans les recrutements. Au Bac, nuancez toujours : \textbf{progrès technique $\neq$ progrès social automatique}.
\end{attention}

\courssub{La fracture numérique : inégalités d'accès et d'usage}

\begin{definition}[Fracture numérique]
La fracture numérique désigne les inégalités d'accès aux technologies de l'information et de la communication (TIC) — infrastructures, terminaux, connectivité, compétences — entre régions (Nord/Sud), territoires (urbain/rural), genres, générations et classes sociales.
\end{definition}

\begin{methode}[Analyser la fracture numérique en trois niveaux (modèle de l'UIT)]
\begin{enumerate}
    \item \textbf{Accès (1\textsuperscript{er} niveau)} : infrastructure (fibre optique, 4G/5G), électricité, coût de la data par rapport au revenu (l'UIT fixe comme cible un coût inférieur à 2\,\% du revenu mensuel pour un forfait d'entrée de gamme).
    \item \textbf{Usage (2\textsuperscript{e} niveau)} : usages basiques (voix, SMS) contre usages avancés (e-commerce, e-gouvernement, formation en ligne, télétravail). Dépend de l'alphabétisation numérique.
    \item \textbf{Impact (3\textsuperscript{e} niveau)} : capacité à transformer l'accès en opportunités économiques, sociales et politiques (création d'entreprise, participation citoyenne, santé).
\end{enumerate}
\end{methode}

\begin{popfigure}
\begin{tikzpicture}[scale=0.9]
% Carte mondiale très simplifiée : continents en blocs colorés selon le taux d'usage d'Internet 2023 (UIT)
\definecolor{nethigh}{HTML}{1B7837}
\definecolor{netmed}{HTML}{7FBC41}
\definecolor{netlow}{HTML}{FDE725}
\definecolor{netverylow}{HTML}{D7191C}

% Amérique du Nord
\fill[nethigh] (-9,3) rectangle (-4,5);
\node[white, font=\tiny\bfseries] at (-6.5,4) {Amérique du Nord};
\node[white, font=\tiny\bfseries] at (-6.5,3.5) {> 90\%};
% Amérique du Sud
\fill[netmed] (-7,-1) rectangle (-3,-4);
\node[font=\tiny] at (-5,-2.2) {Amérique du Sud};
\node[font=\tiny] at (-5,-2.8) {environ 80\%};
% Europe
\fill[nethigh] (0,3) rectangle (3,5.5);
\node[white, font=\tiny\bfseries] at (1.5,4.5) {Europe};
\node[white, font=\tiny\bfseries] at (1.5,4) {environ 90\%};
% Afrique
\fill[netverylow] (0,-1) rectangle (5,-5);
\node[white, font=\tiny\bfseries] at (2.5,-2.7) {Afrique};
\node[white, font=\tiny\bfseries] at (2.5,-3.3) {environ 40\%};
% Asie (simplifiée)
\fill[netlow] (5,1) rectangle (11,5);
\fill[netmed] (5,-2) rectangle (11,1);
\node[font=\tiny] at (8,3.2) {Asie};
\node[font=\tiny] at (8,2.6) {environ 65\% en moyenne};
\node[font=\tiny] at (8,-0.5) {fortes disparités internes};
% Océanie
\fill[nethigh] (11,0) rectangle (13,-2);
\node[white, font=\tiny\bfseries] at (12,-0.7) {Océanie};
\node[white, font=\tiny\bfseries] at (12,-1.3) {> 85\%};

% Légende
\begin{scope}[xshift=14.5cm, yshift=2cm]
\node[font=\small\bfseries, align=center] at (0,1.8) {Taux d'usage\\d'Internet (2023)};
\fill[nethigh] (-1.2,0.6) rectangle (-0.4,1.1); \node[right, font=\tiny] at (-0.4,0.85) {> 80\%};
\fill[netmed] (-1.2,0) rectangle (-0.4,0.5); \node[right, font=\tiny] at (-0.4,0.25) {60 -- 80\%};
\fill[netlow] (-1.2,-0.6) rectangle (-0.4,-0.1); \node[right, font=\tiny] at (-0.4,-0.35) {30 -- 60\%};
\fill[netverylow] (-1.2,-1.2) rectangle (-0.4,-0.7); \node[right, font=\tiny] at (-0.4,-0.95) {< 40\%};
\node[font=\tiny, align=left] at (0,-2) {Source : UIT, Faits et chiffres 2023.\\Carte schématique non contractuelle.};
\end{scope}

% Flèche Nord
\draw[->, thick, popdark] (-10,5) -- (-10,6) node[above, font=\tiny] {N};
\end{tikzpicture}
\legende{Carte schématique mondiale du taux d'usage d'Internet (2023). L'Afrique affiche le taux le plus bas (environ 40\,\%), creusant le fossé avec l'Amérique du Nord et l'Europe (environ 90\,\%).}
\end{popfigure}

\begin{exemplebox}[La fracture numérique au Cameroun : urbain contre rural]
\begin{itemize}
    \item \textbf{Taux d'usage d'Internet} : environ 40 à 45\,\% de la population en 2023 (UIT), mais avec un net contraste entre les grandes villes (Douala, Yaoundé) et les régions de l'Extrême-Nord, du Nord, de l'Adamaoua et de l'Est, où l'accès reste très limité.
    \item \textbf{Coût de la data} : le prix d'un gigaoctet (quelques centaines de FCFA selon les forfaits) pèse lourd dans le budget des ménages, alors que le SMIG est de 41 875 FCFA depuis 2023 ; le seuil d'abordabilité de l'UIT (moins de 2\,\% du revenu) n'est pas atteint pour une large part de la population.
    \item \textbf{Infrastructures} : le réseau fédérateur en fibre optique (Camtel) et les câbles sous-marins (SAT-3/WASC, ACE, SAIL) relient les grandes villes ; le « dernier kilomètre » fait défaut en zone rurale. Le Plan National de Développement 2020--2030 (SND30) et les projets de transformation numérique prévoient l'extension de la fibre optique sur tout le territoire.
    \item \textbf{Genre} : les femmes sont nettement moins susceptibles que les hommes de posséder un smartphone ou d'utiliser l'internet mobile (GSMA 2023), ce qui reproduit les inégalités sociales dans l'espace numérique.
\end{itemize}
\end{exemplebox}

\courssub{Mutations sociales : travail, éducation, lien social}

\begin{aretenir}[Transformations majeures du tissu social]
\begin{itemize}
    \item \textbf{Nouveaux métiers et disparition d'autres} : développeur, analyste de données, community manager, technicien de maintenance des éoliennes et des panneaux solaires, contre secrétaires dactylographes, agents de saisie, caissiers (automatisés). Selon le Forum économique mondial (2020), environ 85 millions d'emplois pourraient être déplacés par l'automatisation d'ici 2025, mais 97 millions de nouveaux emplois créés : le solde est positif, mais la transition est douloureuse pour les non-qualifiés.
    \item \textbf{Télétravail et économie de plateforme} : la pandémie de Covid-19 a accéléré de plusieurs années l'adoption du travail à distance. Au Cameroun, le travail indépendant en ligne, le commerce électronique (Jumia) et les applications de transport (Yango, Gozem) créent des revenus réels mais souvent précaires pour la jeunesse.
    \item \textbf{Réseaux sociaux et espace public} : Facebook, WhatsApp, TikTok et X (ex-Twitter) sont devenus des médias majeurs d'information et de mobilisation (ex. \#EndSARS au Nigeria en 2020, mobilisations liées à la crise anglophone au Cameroun). Risques : bulles de filtre, haine en ligne, ingérences étrangères.
    \item \textbf{Éducation : cours en ligne et ressources ouvertes} : MOOC (Coursera, edX), Khan Academy, Wikipédia, plateformes nationales (enseignement à distance pendant la crise sanitaire au Cameroun). Défis : connexion, équipement, autodiscipline, reconnaissance des certifications.
\end{itemize}
\end{aretenir}

\begin{savaistu}[Le Cameroun et l'appropriation locale des RST]
Ces dernières années, le Cameroun a lancé plusieurs projets phares :
\begin{itemize}
    \item \textbf{Villes intelligentes (\emph{Smart Cities})} : Yaoundé et Douala expérimentent la gestion intelligente du trafic, l'éclairage public connecté et des plateformes citoyennes de signalement des incidents.
    \item \textbf{E-gouvernement} : dématérialisation progressive des démarches (actes d'état civil, casier judiciaire, téléprocédures fiscales à la DGI, guichet unique de création d'entreprise au CFCE). Objectifs : réduire la corruption, les délais et les coûts.
    \item \textbf{Identité numérique} : carte nationale d'identité biométrique et projet d'identifiant unique pour faciliter l'inclusion financière et sociale.
    \item \textbf{FinTech} : le Mobile Money (Orange Money, MTN Mobile Money) compte des millions de comptes actifs et constitue le moteur de l'inclusion financière : le taux de bancarisation classique reste faible, mais plus de la moitié des adultes utilisent des services financiers mobiles.
\end{itemize}
Ces initiatives montrent que les RST ne sont pas seulement subies : elles peuvent être \textbf{appropriées, adaptées et pilotées} par des politiques publiques volontaires.
\end{savaistu}

\courssub{Défis éthiques et politiques : données, surveillance, vérité}

\begin{definition}[Protection des données personnelles]
Ensemble des règles juridiques (RGPD en Europe, cadre législatif camerounais sur la cybersécurité et la protection des données) encadrant la collecte, le traitement, la conservation et le partage des données à caractère personnel (identité, localisation, santé, biométrie, navigation). Principes : consentement, finalité, minimisation, sécurité, droit d'accès, de rectification et d'effacement.
\end{definition}

\begin{aretenir}[Trois enjeux critiques pour la citoyenneté numérique]
\begin{enumerate}
    \item \textbf{Surveillance étatique et privée} : reconnaissance faciale, logiciels espions (affaire Pegasus révélée en 2021), systèmes de crédit social. Au Cameroun, la loi antiterroriste de 2014 et la loi sur la cybercriminalité de 2010 donnent de larges pouvoirs d'interception ; d'où la nécessité d'un contrôle judiciaire indépendant.
    \item \textbf{Manipulation de l'information (désinformation)} : fermes à trolls, robots logiciels (\emph{bots}), hypertrucages (vidéos et audios générés par IA). Impact : élections (États-Unis 2016, Kenya 2017), santé (rumeurs anti-vaccins pendant la Covid-19), conflits (crise anglophone au Cameroun). Réponses : vérification des faits (Africa Check), éducation aux médias et à l'information (EMI), régulation des plateformes (règlement européen sur les services numériques, DSA, 2022).
    \item \textbf{Biais algorithmiques et justice} : une IA entraînée sur des données historiques biaisées reproduit des discriminations (recrutement, crédit bancaire, reconnaissance faciale moins fiable sur les peaux foncées). Exigences : \emph{auditabilité}, \emph{explicabilité}, \emph{diversité des données d'entraînement}.
\end{enumerate}
\end{aretenir}

\begin{exoresolu}[Sujet de dissertation type Bac : « Les révolutions scientifiques et techniques depuis 1945 : progrès ou menace pour l'humanité ? »]
\textbf{Problématique} : en quoi les RST, par leur puissance transformatrice sans précédent, constituent-elles à la fois le principal levier d'amélioration de la condition humaine et la source de risques systémiques inédits, exigeant une gouvernance éthique et politique à la hauteur des enjeux ?

\textbf{Plan détaillé}
\begin{enumerate}
    \item \textbf{Un progrès matériel et sanitaire indéniable (1945--années 2000)}
    \begin{itemize}
        \item Santé : vaccination, antibiotiques, espérance de vie en hausse (environ 46 ans en 1950 contre plus de 73 ans aujourd'hui).
        \item Alimentation : Révolution verte, nourrir plus de 8 milliards d'humains.
        \item Communication et énergie : électricité, internet, mobilité, croissance et sortie de la pauvreté (Chine, Asie de l'Est).
    \end{itemize}
    \item \textbf{Des externalités négatives devenues menaces existentielles (années 1970--aujourd'hui)}
    \begin{itemize}
        \item Nucléaire : bombe, déchets, accidents (Tchernobyl 1986, Fukushima 2011).
        \item Climat : réchauffement d'environ 1,2 \textdegree C depuis l'ère préindustrielle (GIEC), érosion accélérée de la biodiversité.
        \item Numérique : fracture, surveillance, désinformation, IA mal encadrée.
    \end{itemize}
    \item \textbf{La nécessité d'une gouvernance mondiale et locale du progrès technique}
    \begin{itemize}
        \item Droit international : TNP (1968), Accord de Paris (2015), RGPD (2016), recommandation de l'UNESCO sur l'éthique de l'IA (2021).
        \item Politiques publiques : transition énergétique, éducation numérique, régulation des plateformes.
        \item Exemple camerounais : SND30, e-gouvernement, villes intelligentes $\rightarrow$ l'appropriation souveraine est possible.
    \end{itemize}
\end{enumerate}

\textbf{Conclusion} : les RST ne sont ni « bien » ni « mal » en soi : ce sont des \textbf{amplificateurs de puissance}. Elles ont sorti des milliards d'humains de la misère et de la mort précoce, mais ont aussi mis la biosphère et l'espèce elle-même en danger. Le défi du XXI\textsuperscript{e} siècle n'est pas d'arrêter l'innovation, mais de l'\textbf{orienter} par le droit, l'éthique et la démocratie pour qu'elle serve le \emph{bien commun planétaire}. L'Afrique, et le Cameroun en particulier, ne doit pas être simple consommatrice mais \textbf{productrice} de solutions techniques adaptées à ses réalités (agritech, healthtech, fintech, technologies vertes).
\tcblower
\textbf{Corrigé commenté} :
\begin{itemize}
    \item \textbf{Problématique bien posée} : tension progrès/menace, nécessité d'une gouvernance.
    \item \textbf{Chronologie respectée} : 1945--aujourd'hui, avec le tournant des années 1970 (rapport Meadows de 1972, premier choc pétrolier de 1973).
    \item \textbf{Exemples précis et datés} : espérance de vie, GIEC, RGPD, Cameroun.
    \item \textbf{Nuance} : pas de déterminisme technologique, rôle central de l'action politique.
    \item \textbf{Ouverture} : l'Afrique comme actrice, pas seulement victime.
\end{itemize}
\end{exoresolu}

\courssub{Synthèse : le progrès technique, un outil à orienter}

\begin{propriete}[Le progrès scientifique et technique est un outil ; ses effets dépendent de l'usage que les sociétés en font]
Cette proposition résume la leçon entière. Les RST depuis 1945 ont décuplé la puissance d'agir de l'humanité (énergie, information, manipulation du vivant). Cette puissance est \textbf{indifférente aux fins} : elle peut soigner ou tuer, informer ou manipuler, nourrir ou épuiser les sols, relier ou surveiller. L'histoire montre que \textbf{la technique ne détermine pas la société} : ce sont les choix politiques, économiques et éthiques, les luttes sociales et les régulations juridiques qui orientent les usages. Au Cameroun comme ailleurs, l'enjeu est de passer du statut de \textbf{consommateur passif} de technologies importées à celui d'\textbf{acteur souverain} : former des ingénieurs, financer la recherche, légiférer sur les données, réguler les géants du numérique, inventer des modèles frugaux adaptés aux contextes locaux. C'est à cette condition que les RST seront un levier de \textbf{développement endogène} et non un facteur supplémentaire de dépendance.
\end{propriete}

\begin{methode}[Rédiger une conclusion de dissertation (Bac Histoire)]
\begin{enumerate}
    \item \textbf{Rappeler la problématique} : reprenez la question centrale sans la copier mot pour mot.
    \item \textbf{Synthétiser les acquis} : une phrase par grande partie (I, II, III) pour montrer la démonstration.
    \item \textbf{Nuancer et hiérarchiser} : ce qui pèse le plus, ce qui est réversible ou non.
    \item \textbf{Ouvrir sur une perspective} : lien avec l'actualité, une autre région du monde, une autre discipline (économie, philosophie), ou une question prospective.
    \item \textbf{Exemple d'ouverture} : « L'Afrique, par sa jeunesse et son saut technologique (mobile money, drones médicaux), peut-elle inventer une "modernité frugale" qui réconcilie progrès technique et justice sociale ? »
\end{enumerate}
\end{methode}

\begin{attention}[Derniers pièges à éviter pour l'examen]
\begin{itemize}
    \item \textbf{Lister sans analyser} : « Il y a internet, les vaccins, le nucléaire » ne rapporte aucun point. Il faut \textbf{expliquer les mécanismes} (ex. « La vaccination de masse a rompu le cycle épidémique en créant une immunité collective... »).
    \item \textbf{Oublier le Cameroun} : le programme exige des exemples camerounais à chaque chapitre. Pas de Cameroun = note plafonnée.
    \item \textbf{Confondre fracture numérique et fracture sociale} : la fracture numérique \textbf{renforce} les fractures sociales existantes (genre, revenu, territoire), mais elle a sa dynamique propre (coût de la data, infrastructures, compétences).
    \item \textbf{Négliger la chronologie} : les effets ne sont pas les mêmes en 1950 (optimisme nucléaire), en 1970 (prise de conscience écologique), en 1995 (essor d'internet) ou aujourd'hui (IA générative). Situez toujours vos arguments dans le temps.
\end{itemize}
\end{attention}

\newpage

\coursec{Activité d'intégration}

\coursec{Les révolutions scientifiques et techniques depuis la fin de la Deuxième Guerre mondiale}

\courssub{Objectifs de l'activité}
\begin{itemize}
    \item \textbf{Savoirs :} Identifier et classer les grandes découvertes scientifiques (atome, espace, biologie, médecine) et les grandes innovations techniques (informatique, télécoms, transports, énergie, agriculture) depuis 1945.
    \item \textbf{Savoir-faire :} Exploiter un corpus documentaire hétérogène (images, tableaux chiffrés, extraits de presse) ; construire une frise chronologique et un diagramme statistique ; rédiger une conclusion argumentée et nuancée en mobilisant des exemples précis.
    \item \textbf{Attitudes :} Travailler en groupe, confronter des interprétations, prendre conscience de l'impact des RST sur le quotidien des Camerounais et de la persistance de la fracture numérique.
\end{itemize}

\courssub{Situation}
Depuis la fin de la Deuxième Guerre mondiale (1945), le monde connaît une accélération sans précédent des révolutions scientifiques et techniques (RST). Le Cameroun, indépendant depuis 1960, est partie prenante de ces transformations, bien qu'à des rythmes et des échelles variables.

\medskip
\noindent \textbf{Votre mission :} En tant qu'élèves de Terminale technique, vous participez à un atelier « Histoire et Citoyenneté » organisé par le Club UNESCO du lycée. Le thème est : \emph{« Les révolutions scientifiques et techniques : quel bilan pour le Cameroun 75 ans après 1945 ? »}. Vous devez produire une affiche synthétique (frise + graphique + paragraphe conclusif) pour l'exposition de la Journée Mondiale de la Science (10 novembre).

\medskip
\noindent La classe est divisée en \textbf{quatre groupes}. Chaque groupe reçoit un dossier documentaire spécifique. À l'issue du travail de groupe, une mise en commun permet de construire les outils graphiques collectifs avant la rédaction individuelle du paragraphe final.

\medskip
\noindent \textbf{Dossier 1 -- Iconographie légendée (5 documents)}
\begin{center}
\begin{tabularx}{\linewidth}{|l|X|}\hline
\textbf{Doc} & \textbf{Description et légende} \\ \hline
1 & \textbf{Spoutnik 1 (4 octobre 1957).} Premier satellite artificiel mis en orbite par l'URSS. Début de la conquête spatiale. \\ \hline
2 & \textbf{Apollo 11 -- Neil Armstrong sur la Lune (21 juillet 1969).} Aboutissement du programme spatial américain. \\ \hline
3 & \textbf{Première greffe cardiaque humaine (3 décembre 1967, Le Cap).} Pr Christian Barnard. Révolution de la chirurgie et de l'immunologie. \\ \hline
4 & \textbf{ENIAC (1946, Université de Pennsylvanie).} Premier calculateur électronique à tubes, programmable, 30 tonnes. Origine de l'informatique moderne. \\ \hline
5 & \textbf{Smartphone (iPhone 2007 / Android 2008).} Convergence téléphonie, informatique, internet mobile. Objet emblématique de la société de l'information. \\ \hline
\end{tabularx}
\end{center}

\medskip
\noindent \textbf{Dossier 2 -- Tableau de données chronologiques (Repères mondiaux 1945--2020)}
\begin{center}
\begin{tabularx}{\linewidth}{|c|l|X|}\hline
\textbf{Année} & \textbf{Événement} & \textbf{Nature (à classer)} \\ \hline
1945 & Bombardements atomiques (Hiroshima/Nagasaki) & \ldots \\ \hline
1953 & Structure de l'ADN (Watson, Crick, Franklin) & \ldots \\ \hline
1957 & Lancement Spoutnik 1 & \ldots \\ \hline
1969 & Premier pas sur la Lune (Apollo 11) & \ldots \\ \hline
1971 & Premier microprocesseur Intel 4004 & \ldots \\ \hline
1973 & Premier appel téléphone mobile (Motorola DynaTAC) & \ldots \\ \hline
1983 & Protocole TCP/IP (naissance d'Internet) & \ldots \\ \hline
1990 & Projet Génome Humain lancé / World Wide Web (CERN) & \ldots \\ \hline
1996 & Clonage de la brebis Dolly & \ldots \\ \hline
2003 & Séquençage complet du génome humain & \ldots \\ \hline
2012 & Découverte boson de Higgs (CERN) & \ldots \\ \hline
2020 & Vaccins ARN messager (Covid-19) & \ldots \\ \hline
\end{tabularx}
\end{center}

\medskip
\noindent \textbf{Dossier 3 -- Statistiques camerounaises (Évolution 2000--2023)}
\begin{center}
\begin{tabular}{|c|c|c|c|}\hline
\textbf{Année} & \textbf{Abonnés Mobile (millions)} & \textbf{Prod. Hydroélectrique (GWh)} & \textbf{Espérance de vie à la naissance (ans)} \\ \hline
2000 & 0,15 & 2\,850 & 50,2 \\ \hline
2005 & 2,10 & 3\,120 & 51,8 \\ \hline
2010 & 10,50 & 3\,450 & 54,1 \\ \hline
2015 & 18,80 & 4\,100 & 56,7 \\ \hline
2020 & 24,50 & 4\,850 & 59,3 \\ \hline
2023 & 26,80 & 5\,200 & 60,5 \\ \hline
\end{tabular}
\end{center}
\legende{Sources : ART (Agence de Régulation des Télécommunications), AES-SONEL/ENEO, Banque Mondiale / INS. Données arrondies.}

\medskip
\noindent \textbf{Dossier 4 -- Extraits de presse camerounaise (2022--2024)}
\begin{itemize}
    \item \textit{Extrait A (Mutations, mars 2023) :} « Mobile Money : plus de 10 millions de comptes actifs au Cameroun. L'inclusion financière progresse dans les zones rurales non bancarisées. »
    \item \textit{Extrait B (Cameroon Tribune, juin 2024) :} « Fracture numérique : seulement 35\% des ménages ruraux ont accès à l'internet haut débit contre 78\% à Yaoundé et Douala. Le coût du data reste un frein. »
    \item \textit{Extrait C (Le Messager, janvier 2023) :} « E-santé : l'application \emph{GiftedMom} permet le suivi de grossesse par SMS dans le Nord. La télémédecine réduit la mortalité maternelle. »
\end{itemize}

\courssub{Consignes}
\medskip
\noindent \textbf{Phase 1 -- Analyse et classification (Groupe, 20 min)}
\begin{enumerate}
    \item À partir du \textbf{Dossier 1}, classer chaque document dans l'une des deux catégories : « \cle{Découverte scientifique} » (avancée de la connaissance fondamentale) ou « \cle{Innovation technique} » (application pratique, brevet, mise sur le marché). Justifier chaque choix en une phrase.
    \item Compléter la colonne « Nature » du \textbf{Dossier 2} en utilisant les deux catégories ci-dessus.
\end{enumerate}

\medskip
\noindent \textbf{Phase 2 -- Construction graphique collective (Classe entière, 25 min)}
\begin{enumerate}
    \setcounter{enumi}{2}
    \item Au tableau, construire une \textbf{frise chronologique collective (1945--2025)}. Y reporter les 12 événements du Dossier 2 + 3 événements camerounais clés (ex: Indépendance 1960, libéralisation télécoms 2000, arrivée 3G/4G). Utiliser un code couleur : \textcolor{popblue}{Bleu = Sciences}, \textcolor{poporange}{Orange = Techniques}, \textcolor{popgreen}{Vert = Cameroun}.
    \item À partir du \textbf{Dossier 3 (colonne Abonnés Mobile)}, tracer un \textbf{diagramme en barres} (années en abscisse, millions d'abonnés en ordonnée). Identifier et nommer \textbf{trois phases de croissance} (ex: démarrage, explosion, maturation/saturation).
\end{enumerate}

\medskip
\noindent \textbf{Phase 3 -- Synthèse et argumentation (Individuel, 25 min)}
\begin{enumerate}
    \setcounter{enumi}{4}
    \item Rédiger un \textbf{paragraphe conclusif unique (15--20 lignes)} répondant à la question : \emph{« Les révolutions scientifiques et techniques ont-elles transformé la vie des Camerounais depuis 1945 ? »}.
    \item \textbf{Contraintes obligatoires :}
    \begin{itemize}
        \item Mobiliser \textbf{au moins trois exemples précis} (dates, chiffres, noms propres) issus des dossiers.
        \item Intégrer \textbf{une nuance explicite} sur la \cle{fracture numérique} (territoriale, sociale, économique) en s'appuyant sur le Dossier 4.
        \item Structurer le paragraphe : Thèse (affirmation) -- Arguments chiffrés/exemples -- Nuance -- Conclusion ouverte.
    \end{itemize}
\end{enumerate}

\courssub{Ressources à mobiliser}
\begin{methode}[Construire un diagramme en barres statistique]
\begin{enumerate}
    \item Choisir l'échelle des ordonnées (ex: 1 cm pour 5 millions d'abonnés) et des abscisses (1 cm par année ou par pas de 5 ans).
    \item Tracer les axes orthonormés, légender : « Évolution du nombre d'abonnés mobile au Cameroun (2000--2023) ».
    \item Représenter chaque valeur par un rectangle de largeur constante, hauteur proportionnelle à l'effectif.
    \item Colorier les barres selon les phases identifiées (ex: clair, moyen, foncé).
\end{enumerate}
\end{methode}

\begin{methode}[Rédiger un paragraphe argumenté en Histoire]
\begin{enumerate}
    \item \textbf{Phrase d'amorce (Thèse) :} Réponse directe à la question (Oui, transformation majeure mais inégale).
    \item \textbf{Arguments (Preuves) :} Enchaîner 3 exemples datés/chiffrés (ex: « En 2023, 26,8 millions d'abonnés mobiles... » ; « Depuis 2010, l'espérance de vie gagne 6 ans... » ; « Le Mobile Money (10 M comptes) bancarise le rural... »).
    \item \textbf{Nuance (Limite) :} Introduire par « Cependant / Toutefois / Néanmoins » + fait précis (ex: « Toutefois, l'extrait B révèle une fracture numérique : 35\% d'accès rural vs 78\% urbain »).
    \item \textbf{Conclusion (Ouverture) :} Phrase finale résumant le bilan et ouvrant sur un défi futur (ex: « Le défi reste l'inclusion numérique totale »).
\end{enumerate}
\end{methode}

\begin{aretenir}[Critères de distinction Sciences / Techniques]
\textbf{Découverte scientifique :} Nouveau savoir sur la nature (loi, structure, particule). Ex: ADN (1953), Boson de Higgs (2012).\\
\textbf{Innovation technique :} Nouvel objet, procédé ou service mis à disposition de la société. Ex: Microprocesseur (1971), Internet (1983), Smartphone (2007).\\
\emph{Articulation :} La science alimente la technique (ex: physique quantique $\to$ transistor) ; la technique outille la science (ex: télescope spatial $\to$ astrophysique).
\end{aretenir}

\courssub{Démarche et corrigé commenté}
\begin{exoresolu}[Corrigé guidé de l'activité d'intégration]
\textbf{Énoncé de l'activité :} À partir des 4 dossiers fournis, classer les événements (Sciences/Techniques), construire une frise chronologique 1945-2025 et un diagramme en barres (abonnés mobile Cameroun 2000-2023), puis rédiger un paragraphe argumenté sur la transformation de la vie des Camerounais.

\tcblower

\textbf{Phase 1 -- Classification (Dossiers 1 \& 2)}

\medskip
\noindent \textbf{Dossier 1 :}
\begin{itemize}
    \item Doc 1 (Spoutnik) : \textbf{Innovation technique} (premier objet artificiel en orbite, prouesse d'ingénierie fusée).
    \item Doc 2 (Apollo 11) : \textbf{Innovation technique} (système complexe lanceur/module/vaisseau, aboutissement programme Apollo).
    \item Doc 3 (Greffe cardiaque) : \textbf{Innovation technique} (acte chirurgical nouveau, protocole immunosuppression). \emph{Note : repose sur découvertes scientifiques (immunologie, circulation sanguine).}
    \item Doc 4 (ENIAC) : \textbf{Innovation technique} (machine calculatoire électronique, brevet, usage militaire puis civil).
    \item Doc 5 (Smartphone) : \textbf{Innovation technique} (produit grand public intégrant multiples brevets : tactile, GPS, 3G, OS).
\end{itemize}

\medskip
\noindent \textbf{Dossier 2 -- Colonne « Nature » complétée :}
\begin{center}
\begin{tabularx}{\linewidth}{|c|l|X|}\hline
\textbf{Année} & \textbf{Événement} & \textbf{Nature} \\ \hline
1945 & Bombardements atomiques & \textbf{Technique} (application militaire de la fission) \\ \hline
1953 & Structure de l'ADN & \textbf{Scientifique} (modèle double hélice) \\ \hline
1957 & Spoutnik 1 & \textbf{Technique} \\ \hline
1969 & Apollo 11 & \textbf{Technique} \\ \hline
1971 & Microprocesseur Intel 4004 & \textbf{Technique} (circuit intégré VLSI) \\ \hline
1973 & Premier mobile (DynaTAC) & \textbf{Technique} \\ \hline
1983 & TCP/IP / Internet & \textbf{Technique} (protocole/réseau) \\ \hline
1990 & Génome Humain / Web & \textbf{Scientifique} (projet connaissance) / \textbf{Technique} (Web) \\ \hline
1996 & Clonage Dolly & \textbf{Scientifique} (preuve totipotence) + \textbf{Technique} (méthode SCNT) \\ \hline
2003 & Génome séquencé & \textbf{Scientifique} (donnée brute) \\ \hline
2012 & Boson de Higgs & \textbf{Scientifique} (validation Modèle Standard) \\ \hline
2020 & Vaccins ARNm & \textbf{Technique} (plateforme vaccinale) issue \textbf{Scientifique} (biologie ARN) \\ \hline
\end{tabularx}
\end{center}

\medskip
\noindent \textbf{Phase 2 -- Frise chronologique collective (1945--2025)}

\begin{popfigure}
\begin{tikzpicture}[scale=0.9, every node/.style={font=\footnotesize}]
% Échelle : 80 ans = 18 cm => 1 an = 0.225 cm. Origine 1945 à x=0.
% Graduations tous les 10 ans (2.25 cm)
\draw[->, thick] (0,0) -- (18.5,0) node[right] {Temps};
\foreach \x/\y in {0/1945, 2.25/1955, 4.5/1965, 6.75/1975, 9/1985, 11.25/1995, 13.5/2005, 15.75/2015, 18/2025}
    \draw (\x,0.1) -- (\x,-0.1) node[below] {\y};

% Événements Mondiaux - Sciences (Bleu) : Au-dessus de l'axe (y>0)
% 1945 Fission militaire (x=0)
\draw[popblue, thick] (0, 1.2) node[above, align=center, popblue] {1945\\Fission\\militaire} -- (0, 0.2);
% 1953 ADN (x=1.8)
\draw[popblue, thick] (1.8, 1.6) node[above, align=center, popblue] {1953\\ADN} -- (1.8, 0.2);
% 1990 Génome Humain (x=10.125)
\draw[popblue, thick] (10.125, 1.2) node[above, align=center, popblue] {1990\\Génome\\Humain} -- (10.125, 0.2);
% 2003 Génome complet (x=13.05)
\draw[popblue, thick] (13.05, 1.6) node[above, align=center, popblue] {2003\\Génome\\complet} -- (13.05, 0.2);
% 2012 Boson Higgs (x=15.075)
\draw[popblue, thick] (15.075, 1.2) node[above, align=center, popblue] {2012\\Boson\\Higgs} -- (15.075, 0.2);
% 2020 Vaccin ARNm (x=16.875)
\draw[popblue, thick] (16.875, 1.6) node[above, align=center, popblue] {2020\\Vaccin\\ARNm} -- (16.875, 0.2);

% Événements Mondiaux - Techniques (Orange) : En-dessous de l'axe (y<0)
% 1957 Spoutnik (x=2.7)
\draw[poporange, thick] (2.7, -1.2) node[below, align=center, poporange] {1957\\Spoutnik} -- (2.7, -0.2);
% 1969 Lune (x=5.4)
\draw[poporange, thick] (5.4, -1.6) node[below, align=center, poporange] {1969\\Lune} -- (5.4, -0.2);
% 1971 Microproc (x=5.85)
\draw[poporange, thick] (5.85, -1.2) node[below, align=center, poporange] {1971\\Microproc.} -- (5.85, -0.2);
% 1973 Mobile 1G (x=6.3)
\draw[poporange, thick] (6.3, -1.6) node[below, align=center, poporange] {1973\\Mobile 1G} -- (6.3, -0.2);
% 1983 TCP/IP (x=8.55)
\draw[poporange, thick] (8.55, -1.2) node[below, align=center, poporange] {1983\\TCP/IP} -- (8.55, -0.2);
% 1990 Web (x=10.125)
\draw[poporange, thick] (10.125, -1.6) node[below, align=center, poporange] {1990\\Web} -- (10.125, -0.2);
% 1996 Dolly (x=11.475)
\draw[poporange, thick] (11.475, -1.2) node[below, align=center, poporange] {1996\\Dolly} -- (11.475, -0.2);
% 2007 Smartphone (x=13.95)
\draw[poporange, thick] (13.95, -1.6) node[below, align=center, poporange] {2007\\Smartphone} -- (13.95, -0.2);

% Événements Cameroun (Vert) : Au-dessus, plus haut (y>1.8)
% 1960 Indép. (x=3.375)
\draw[popgreen, thick] (3.375, 2.2) node[above, align=center, popgreen] {1960\\Indépendance} -- (3.375, 0.2);
% 2000 Libéral. Télécoms (x=12.375)
\draw[popgreen, thick] (12.375, 2.6) node[above, align=center, popgreen] {2000\\Libéralisation\\Télécoms} -- (12.375, 0.2);
% 2010 3G/4G (x=14.625)
\draw[popgreen, thick] (14.625, 2.2) node[above, align=center, popgreen] {2010\\Arrivée 3G/4G} -- (14.625, 0.2);

% Légende
\node[align=left, anchor=south west] at (0.5,-3.0) {
\textcolor{popblue}{\rule{4pt}{4pt}} Sciences \quad
\textcolor{poporange}{\rule{4pt}{4pt}} Techniques \quad
\textcolor{popgreen}{\rule{4pt}{4pt}} Cameroun
};
\end{tikzpicture}
\legende{Frise chronologique synthétique 1945--2025 : articulation Sciences (haut bleu), Techniques (bas orange), Jalons Cameroun (haut vert). Échelle respectée (1 an = 0.225 cm).}
\end{popfigure}

\medskip
\noindent \textbf{Phase 2 -- Diagramme en barres : Abonnés Mobile Cameroun (2000--2023)}

\begin{popfigure}
\begin{tikzpicture}
\begin{axis}[
    ybar,
    bar width=14pt,
    width=\linewidth,
    height=8cm,
    xlabel={Années},
    ylabel={Abonnés (millions)},
    xtick=data,
    xticklabels={2000,2005,2010,2015,2020,2023},
    ymin=0, ymax=30,
    ymajorgrids=true,
    grid style=dashed,
    legend style={at={(0.5,-0.15)},anchor=north,legend columns=-1},
    nodes near coords,
    nodes near coords align={vertical},
    every node near coord/.append style={font=\tiny, rotate=90, anchor=west, yshift=2pt}
]
\addplot[fill=popblue!60] coordinates {(0,0.15) (1,2.1) (2,10.5) (3,18.8) (4,24.5) (5,26.8)};
\legend{Abonnés Mobile}
\end{axis}
\end{tikzpicture}
\legende{Diagramme en barres : Évolution du nombre d'abonnés mobile au Cameroun (2000--2023). Source : ART/INS.}
\end{popfigure}

\medskip
\noindent \textbf{Analyse des trois phases de croissance :}
\begin{enumerate}
    \item \textbf{Phase 1 -- Démarrage lent (2000--2005) :} De 0,15 à 2,1 millions (+1300\%). Contexte : libéralisation du secteur (2000), arrivée de MTN (2000) et Orange (1999), coût élevé des terminaux et communications, couverture urbaine seulement.
    \item \textbf{Phase 2 -- Explosion exponentielle (2005--2015) :} De 2,1 à 18,8 millions (x9 en 10 ans). Facteurs : baisse drastique des prix (prix d'appel, terminaux low-cost), arrivée de la 3G (2010), forfaits prépayés adaptés, engouement jeunes, Mobile Money (dès 2011).
    \item \textbf{Phase 3 -- Maturation et saturation relative (2015--2023) :} De 18,8 à 26,8 millions (+43\% en 8 ans). Croissance ralentie. Taux de pénétration $\sim$ 95\% (multi-SIM). Enjeux : 4G/5G, data, inclusion financière, qualité de service, zones blanches.
\end{enumerate}

\medskip
\noindent \textbf{Phase 3 -- Proposition de paragraphe conclusif (Corrigé type)}

\medskip
\noindent \textbf{Thèse :} Oui, les révolutions scientifiques et techniques ont profondément transformé la vie des Camerounais depuis 1945, notamment par la diffusion massive des technologies de l'information et de la communication et les progrès de la santé, mais cette transformation reste inégalitaire du fait d'une fracture numérique persistante.

\medskip
\noindent \textbf{Arguments :} D'abord, la \textbf{révolution des télécommunications} a métamorphosé les usages quotidiens : le nombre d'abonnés mobiles est passé de \textbf{0,15 million en 2000 à 26,8 millions en 2023} (Dossier 3), soit une pénétration quasi universelle, favorisant l'émergence du \textbf{Mobile Money} qui compte \textbf{plus de 10 millions de comptes actifs en 2023} (Dossier 4, Extrait A) et assure l'inclusion financière des populations rurales non bancarisées. Ensuite, les \textbf{progrès de la biologie et de la médecine} (découverte de l'ADN 1953, vaccins ARNm 2020, télémédecine) ont contribué à l'allongement de l'\textbf{espérance de vie de 50,2 ans (2000) à 60,5 ans (2023)} (Dossier 3), tandis que des applications comme \emph{GiftedMom} améliorent le suivi maternel dans le Nord (Dossier 4, Extrait C). Enfin, l'\textbf{électrification} appuyée sur l'hydroélectricité (production passée de 2\,850 à 5\,200 GWh) soutient l'urbanisation et l'activité économique.

\medskip
\noindent \textbf{Nuance :} \textbf{Toutefois}, cette modernité est à géométrie variable. L'extrait B (2024) révèle une \textbf{fracture numérique territoriale majeure} : seulement \textbf{35\% des ménages ruraux} accèdent à l'internet haut débit contre \textbf{78\% à Yaoundé et Douala}, le coût du data demeurant un frein structurel pour les plus pauvres.

\medskip
\noindent \textbf{Conclusion :} Le bilan est donc contrasté : les RST ont indéniablement connecté et soigné le Cameroun, mais l'enjeu majeur des prochaines décennies réside dans la \textbf{résorption de la fracture numérique} pour faire de l'innovation un levier de développement inclusif et non un facteur d'exclusion.
\end{exoresolu}

\courssub{Bilan de l'activité}
Cette activité d'intégration a permis de :
\begin{itemize}
    \item \textbf{Structurer les connaissances :} La classification (Sciences vs Techniques) et la frise chronologique ont mis en évidence l'enchaînement cause/effet (Science $\to$ Technique) et l'accélération temporelle (cycles d'innovation raccourcis).
    \item \textbf{Quantifier le changement :} Le diagramme en barres a objectivé la « révolution mobile » camerounaise en trois phases lisibles (démarrage, explosion, maturation), reliant macro-données (ART) et micro-usages (Mobile Money).
    \item \textbf{Complexifier le regard :} La rédaction argumentée a forcé la nuance : le « oui » massif aux transformations (santé, communication, finance) ne masque pas la persistance des inégalités d'accès (fracture numérique territoriale/économique), notion clé du programme d'Histoire pour comprendre le monde contemporain.
    \item \textbf{Valider les savoir-faire :} Exploitation de sources variées (iconographie, statistiques, presse), construction d'outils graphiques normalisés, rédaction historique structurée (Thèse/Arguments/Nuance/Conclusion) -- compétences évaluées au Probatoire/Baccalauréat.
\end{itemize}

\begin{attention}[Pièges évités dans le corrigé]
\begin{itemize}
\item Ne pas confondre « découverte » (ADN 1953) et « innovation » (séquençage 2003, thérapie génique).
\item Ne pas oublier d'indiquer l'unité (millions) et l'échelle sur le graphique.
\item Ne pas lister les exemples sans les \emph{articuler} à la thèse (connecteurs logiques : « D'abord », « Ensuite », « Enfin »).
\item La nuance ne doit pas contredire la thèse mais la \emph{préciser} (transformé \emph{mais} inégalement).
\end{itemize}
\end{attention}

\newpage

\coursec{Méthodes et exercices résolus}

\coursec{Introduction : notions et contexte de l'accélération scientifique (1945-...)}

La fin de la Deuxième Guerre mondiale (1945) marque une rupture majeure dans l'histoire des sciences et des techniques. Le conflit a agi comme un formidable accélérateur : projets Manhattan (bombe atomique), radar, moteurs à réaction, cryptanalyse (ordinateurs naissants), production massive de pénicilline. Après 1945, cet élan se prolonge dans un contexte de \cle{guerre froide} (1947-1991), où la rivalité entre les \cle{États-Unis} et l'\cle{URSS} transforme la recherche scientifique en enjeu de puissance et de propagande. Parallèlement, la \cle{décolonisation} ouvre la voie à de nouveaux États (dont le Cameroun en 1960) qui doivent construire leur développement technique.

\begin{definition}[Révolutions scientifiques et techniques (RST)]
On appelle \cle{révolutions scientifiques et techniques} l'ensemble des découvertes fondamentales (physique, biologie) et des innovations appliquées (informatique, énergie, transports, agriculture) qui, depuis 1945, ont modifié en profondeur les conditions de vie, les structures économiques, les rapports de force géopolitiques et la représentation du monde. Elles se caractérisent par leur \cle{vitesse}, leur \cle{mondialisation} et leur \cle{caractère systémique} (interdépendance des domaines).
\end{definition}

\begin{savaistu}[Le Cameroun à l'heure des RST]
Dès l'indépendance (1960), le Cameroun s'insère dans ces dynamiques : construction de barrages hydroélectriques (Édéa, Song Loulou, Lagdo, Memve'ele), adoption de la téléphonie mobile (années 2000), programmes de vaccination élargie (PEV), recherche agronomique (IRAD, variétés améliorées de maïs, manioc, plantain). Le pays illustre à la fois les opportunités (inclusion financière par mobile money) et les défis (fracture numérique, dépendance technologique, pression sur l'environnement).
\end{savaistu}

\begin{aretenir}[Repères chronologiques majeurs (1945-2025)]
\begin{center}
\begin{tabularx}{\linewidth}{|c|X|}\hline
\textbf{Période} & \textbf{Événements emblématiques} \\ \hline
1945-1957 & Bombe atomique (1945), Transistor (1947), ADN (1953), Spoutnik (1957) \\ \hline
1958-1975 & Circuit intégré (1958), Gagarine (1961), Apollo 11 (1969), Microprocesseur (1971), Apollo-Soyouz (1975) \\ \hline
1976-1995 & PCR (1983), Web (1989), GSM (1991), Dolly (1996), ISS (1998) \\ \hline
1996-... & Génome humain (2003), iPhone (2007), CRISPR-Cas9 (2012), IA générative (2020s), Réutilisation fusées (SpaceX) \\ \hline
\end{tabularx}
\end{center}
\end{aretenir}

\begin{popfigure}
\begin{tikzpicture}[scale=0.95]
\draw[->,thick,popdark] (0,0) -- (15,0);
\foreach \x/\y/\label in {
  0.5/1945/Bombe A,
  2.0/1947/Transistor,
  3.5/1953/ADN,
  5.0/1957/Spoutnik,
  6.5/1961/Gagarine,
  8.5/1969/Lune,
  9.5/1971/µP,
  11.0/1983/PCR,
  12.0/1989/Web,
  13.0/2003/Génome,
  14.0/2012/CRISPR
}{
  \draw[thick,popdark] (\x,-0.1) -- (\x,0.1);
  \node[below,rotate=45,anchor=north east,popdark,font=\tiny] at (\x,-0.15) {\y};
  \node[above,align=center,popblue,font=\scriptsize] at (\x,0.3) {\label};
}
\node[above,popdark,font=\small\bfseries] at (7.5,1.2) {Frise synthétique des RST (1945-2025)};
\end{tikzpicture}
\legende{Frise chronologique des grandes étapes des révolutions scientifiques et techniques depuis 1945.}
\end{popfigure}

\coursec{Les grandes découvertes scientifiques : maîtrise de l'atome et conquête de l'espace}

\courssub{La maîtrise de l'atome : de la bombe à l'énergie civile}

Le projet \cle{Manhattan} (États-Unis, 1942-1945) aboutit aux premiers essais (Trinity, 16 juillet 1945) et aux bombardements d'\cle{Hiroshima} (6 août) et \cle{Nagasaki} (9 août 1945). Cette maîtrise de la fission nucléaire ouvre l'ère atomique : dissuasion nucléaire (course aux armements, crise de Cuba 1962), mais aussi applications civiles. La première centrale nucléaire industrielle s'ouvre à \cle{Calder Hall} (Royaume-Uni, 1956), suivie par la France (Marcoule, Chinon) et l'URSS (Obninsk 1954). La recherche sur la \cle{fusion} (projet \cle{ITER} à Cadarache, coopération internationale 35 pays) vise une énergie quasi inépuisable et peu radioactive.

\begin{aretenir}[Dates clés : L'atome]
\begin{itemize}
\item \textbf{1945} : Essai Trinity (USA), Hiroshima/Nagasaki — entrée dans l'ère nucléaire.
\item \textbf{1954} : Obninsk (URSS), première centrale nucléaire raccordée au réseau.
\item \textbf{1956} : Calder Hall (R.-U.), première centrale industrielle commerciale.
\item \textbf{1986} : Tchernobyl (URSS) — prise de conscience des risques majeurs.
\item \textbf{2007-...} : ITER (France) — démonstrateur de fusion thermonucléaire.
\end{itemize}
\end{aretenir}

\courssub{La conquête de l'espace : course à la Lune et applications}

La \cle{course à l'espace} est le théâtre le plus visible de la guerre froide. L'URSS prend l'avantage : \cle{Spoutnik 1} (4 octobre 1957, premier satellite), \cle{Youri Gagarine} (12 avril 1961, premier homme dans l'espace). Les États-Unis répondent par le programme \cle{Apollo} : \cle{Neil Armstrong} et \cle{Buzz Aldrin} marchent sur la Lune le \cle{21 juillet 1969} (Apollo 11). La détente permet la coopération : \cle{Apollo-Soyouz} (1975), puis \cle{Station spatiale internationale (ISS)} (1998-, 15 pays). Les retombées sont immenses : satellites de \cle{télécommunications} (TV, téléphone, GPS), \cle{météorologie}, \cle{observation de la Terre} (climat, ressources), \cle{navigation} (GPS américain, Galiléo européen, Glonass russe, Beidou chinois).

\begin{exemplebox}[Retombées spatiales au quotidien]
\begin{itemize}
\item \textbf{GPS / Galiléo :} géolocalisation précise (transport, agriculture de précision, secours).
\item \textbf{Météosat / satellites météo :} prévisions fiables, alertes cyclones, gestion de l'eau.
\item \textbf{Télécoms :} télévision par satellite, internet dans zones isolées, téléphonie maritime/aérienne.
\item \textbf{Observation (Spot, Sentinel, Landsat) :} suivi déforestation, urbanisation, rendements agricoles, catastrophes.
\end{itemize}
\end{exemplebox}

\begin{savaistu}[L'espace et le Cameroun]
Le Cameroun utilise les données spatiales pour la cartographie (CCT), la gestion des ressources forestières (satellites Sentinel), la météorologie (ASECNA) et les télécommunications (satellites Intelsat, puis fibre optique). Le pays a adhéré aux traités spatiaux de l'ONU (1967) et participe à l'Agence spatiale africaine (AfSA, 2023).
\end{savaistu}

\coursec{Les grandes découvertes scientifiques : révolution de la biologie et de la médecine}

\courssub{De l'ADN au génome humain : la révolution moléculaire}

La découverte de la structure en \cle{double hélice de l'ADN} par \cle{James Watson}, \cle{Francis Crick} (avec les données de \cle{Rosalind Franklin} et \cle{Maurice Wilkins}) en \cle{1953} fonde la biologie moléculaire. S'ensuivent le déchiffrage du \cle{code génétique} (années 1960, Nirenberg, Khorana), l'\cle{ADN recombinant} (1973, Cohen, Boyer — naissance des OGM et biotechnologies), la \cle{PCR} (\cle{Polymerase Chain Reaction}, \cle{Kary Mullis}, 1983 — amplification d'ADN, diagnostic, police scientifique). Le \cle{Projet Génome Humain} (1990-2003, consortium public + Celera) livre la séquence de référence en \cle{2003}. L'outil \cle{CRISPR-Cas9} (2012, \cle{Emmanuelle Charpentier}, \cle{Jennifer Doudna}) permet l'édition ciblée du génome (thérapies géniques, agriculture).

\begin{aretenir}[Dates clés : Biologie et Médecine]
\begin{itemize}
\item \textbf{1953} : Structure de l'ADN (Watson, Crick, Franklin, Wilkins).
\item \textbf{1955/1961} : Vaccins antipoliomyélitiques (Salk injecté, Sabin oral).
\item \textbf{1973} : ADN recombinant (début des biotechnologies/OGM).
\item \textbf{1978} : Premier « bébé-éprouvette » (FIV, Louise Brown, Steptoe/Edwards).
\item \textbf{1983} : PCR (Mullis) — révolution du diagnostic moléculaire.
\item \textbf{2003} : Achèvement du séquençage du génome humain.
\item \textbf{2012} : CRISPR-Cas9 (Charpentier, Doudna) — édition génomique.
\item \textbf{2020} : Vaccins à ARN messager (Covid-19, Karikó/Weissman) — nouvelle plateforme vaccinale.
\end{itemize}
\end{aretenir}

\courssub{Progrès médicaux et espérance de vie}

La généralisation des \cle{antibiotiques} (pénicilline post-1945), la \cle{vaccination} de masse (variole éradiquée 1980, polio quasi éradiquée), l'\cle{imagerie médicale} (scanner \cle{1972} Hounsfield/Cormack, IRM \cle{1977} Lauterbur/Mansfield), la \cle{chirurgie transplantatoire} (premier cœur \cle{Christiaan Barnard} 1967, rein, foie, moelle), la \cle{contraception} (pilule 1960) et la \cle{FIV} transforment la démographie : l'espérance de vie mondiale passe de \cle{$\sim$45 ans (1950)} à \cle{$\sim$73 ans (2022)}. Au Cameroun, elle passe de $\sim$40 ans (1960) à $\sim$60 ans (2023), grâce au PEV, à la lutte contre le paludisme (moustiquaires imprégnées, nouveaux traitements) et à l'amélioration de la santé maternelle.

\begin{savaistu}[Santé au Cameroun : défis et avancées]
Le Cameroun bénéficie des vaccins (PEV : BCG, DTC, Hépatite B, Rougeole, Fièvre jaune, VAR, Pneumocoque, Rotavirus, HPV), des traitements ARV (VIH/Sida), des ACT (paludisme). Défis persistants : mortalité maternelle/infantile, maladies non transmissibles (diabète, hypertension), accès aux soins en zone rurale, dépendance aux importations de médicaments (promotion de la pharmacopée locale et production locale).
\end{savaistu}

\coursec{Les grandes innovations techniques : informatique, internet et télécommunications}

\courssub{De l'ordinateur central au smartphone : la loi de Moore}

L'invention du \cle{transistor} (1947, \cle{Bardeen, Brattain, Shockley}, Bell Labs) remplace les tubes, puis le \cle{circuit intégré} (1958, \cle{Jack Kilby} Texas Instruments, \cle{Robert Noyce} Fairchild) et le \cle{microprocesseur} (1971, \cle{Intel 4004}) permettent la miniaturisation exponentielle (\cle{loi de Moore} : doublement des transistors tous les 2 ans). Naissent l'\cle{ordinateur personnel} (Altair 1975, Apple II 1977, IBM PC 1981), l'interface graphique (Xerox PARC, Macintosh 1984, Windows), puis le \cle{smartphone} (iPhone 2007, Android 2008) : ordinateur de poche connecté en permanence.

\begin{aretenir}[Dates clés : Informatique et Télécoms]
\begin{itemize}
\item \textbf{1947} : Transistor (Bell Labs).
\item \textbf{1958} : Circuit intégré (Kilby, Noyce).
\item \textbf{1969} : Arpanet (premier réseau à commutation de paquets, ancêtre d'Internet).
\item \textbf{1971} : Microprocesseur Intel 4004.
\item \textbf{1983} : TCP/IP (protocole standard d'Internet), DNS.
\item \textbf{1989} : World Wide Web (Tim Berners-Lee, CERN).
\item \textbf{1991} : Standard GSM (2G, téléphonie mobile numérique européenne).
\item \textbf{2007} : iPhone (smartphone tactile, App Store).
\item \textbf{2010s} : 4G, Cloud computing, Big Data, IA (Deep Learning).
\end{itemize}
\end{aretenir}

\courssub{Internet : du réseau militaire au Web planétaire}

\cle{Arpanet} (1969, DARPA, USA) relie 4 universités. L'adoption de \cle{TCP/IP} (1983) crée l'\cle{Internet} (réseau des réseaux). L'invention du \cle{World Wide Web} (1989, \cle{Tim Berners-Lee}, \cle{CERN}) : \cle{HTML}, \cle{HTTP}, \cle{URL} rend le réseau accessible au grand public (navigateurs Mosaic 1993, Netscape, IE, Chrome). La généralisation de la \cle{fibre optique} (années 1980-2000, débit Tbit/s), des câbles sous-marins (99\% du trafic intercontinental) et de la \cle{4G/5G} (débits mobiles > 100 Mbit/s à 1 Gbit/s) réalise le \cle{village planétaire}.

\begin{popfigure}
\begin{tikzpicture}
\begin{axis}[
  ybar, bar width=12pt,
  width=\linewidth, height=7cm,
  symbolic x coords={2000,2005,2010,2015,2020,2023},
  xtick=data,
  xlabel={Année}, ylabel={Internautes (milliards)},
  ymin=0, ymax=5.5,
  ymajorgrids=true, grid style=dashed,
  nodes near coords, nodes near coords align={vertical},
  every node near coord/.append style={font=\tiny, rotate=90, anchor=west, yshift=2pt},
  popblue!80!popdark,
  enlarge x limits=0.2
]
\addplot coordinates {
  (2000,0.4)
  (2005,1.0)
  (2010,2.0)
  (2015,3.2)
  (2020,4.6)
  (2023,5.3)
};
\end{axis}
\end{tikzpicture}
\legende{Évolution du nombre d'internautes dans le monde (source : UIT/World Bank).}
\end{popfigure}

\begin{methode}[Calculer un taux de variation (rappel)]
Pour mesurer l'évolution d'une grandeur entre deux dates :
\[ \text{Taux de variation (\%)} = \frac{V_{\text{finale}} - V_{\text{initiale}}}{V_{\text{initiale}}} \times 100 \]
\textbf{Exemple (2000-2023) :} $V_i = 0,4$ Md, $V_f = 5,3$ Md. $\frac{5,3 - 0,4}{0,4} \times 100 = \frac{4,9}{0,4} \times 100 = 12,25 \times 100 = 1\,225 \%$. Le nombre d'internautes a été multiplié par $\approx 13$ en 23 ans.
\end{methode}

\coursec{Les grandes innovations techniques : transports, énergie et agriculture}

\courssub{Transports : vitesse, masse et conteneurisation}

L'\cle{avion à réaction} (De Havilland Comet 1952, Boeing 707 1958) divise les temps de trajet intercontinentaux. Le \cle{TGV} (France, 1981, 270 km/h puis 320 km/h) et le \cle{Shinkansen} (Japon, 1964) révolutionnent le rail. La \cle{conteneurisation} (standard ISO, \cle{Malcom McLean}, 1956) standardise le fret maritime : coûts divisés par 20, temps de manutention divisé par 100, moteur de la \cle{mondialisation} des chaînes de valeur. L'\cle{espace} devient accessible : navette spatiale (1981-2011), puis \cle{New Space} (SpaceX Falcon 9 réutilisable 2015, Starship) abaisse le coût au kg en orbite.

\courssub{Énergie : nucléaire, fossiles et renouvelables}

La consommation mondiale d'énergie primaire est multipliée par \cle{3,5} (1950-2023). Les \cle{énergies fossiles} (pétrole, gaz, charbon) restent dominantes ($\sim$82\% en 2022) mais leur part baisse. Le \cle{nucléaire} (fission) fournit $\sim$10\% de l'électricité mondiale (faible CO2, déchets, risque). Les \cle{EnR} (hydraulique, éolien, solaire PV, biomasse) progressent vite : solaire PV $\times$ 1000 (2000-2023), éolien $\times$ 50. Le \cle{stockage} (batteries Li-ion, hydrogène vert, pompage-turbinage) et les \cle{réseaux intelligents} (smart grids) sont les clés de la transition.

\begin{savaistu}[Énergie au Cameroun]
Mix électrique : $\sim$60\% hydraulique (barrages : Edéa, Song Loulou, Lagdo, Memve'ele, Nachtigal 420 MW mis en service 2023-2024), $\sim$40\% thermique (fioul, gaz). Objectif : 3000 MW à l'horizon 2030. Potentiel solaire (Nord), éolien (Littoral), biomasse. Défis : accès rural (électrification $\sim$65\% national, $<$30\% rural), pertes réseau, coût du kWh.
\end{savaistu}

\courssub{Agriculture : révolution verte et biotechnologies}

La \cle{Révolution verte} (années 1960-1980, \cle{Norman Borlaug}, Prix Nobel 1970) : variétés naines à haut rendement (blé, riz IR8), engrais chimiques (Haber-Bosch), irrigation, mécanisation. Résultat : production céréalière mondiale $\times$ 3 (1960-2020), famine évitée en Asie/LatAm. Limites : pollution (nitrates, pesticides), érosion sols, perte biodiversité, dépendance intrants, inégalités (Afrique subsaharienne peu touchée). \cle{Deuxième révolution} : OGM (résistance insectes/herbicides, 1996-), agriculture de précision (GPS, drones, capteurs), \cle{CRISPR} (variétés améliorées sans transgénèse), agriculture urbaine/verticale.

\begin{exemplebox}[Agriculture au Cameroun : entre tradition et modernité]
\begin{itemize}
\item \textbf{IRAD} (Institut de recherche agricole) : variétés améliorées (maïs CMS, manioc TMS, plantain PITA), semences certifiées.
\item \textbf{Projets} : PIDMA, PADFA, ACEFA — intrants subventionnés, mécanisation (tracteurs, décortiqueuses), irrigation (bords de fleuve Logone, Bénoué).
\item \textbf{Enjeux} : autosuffisance (riz, maïs, poisson), transformation locale (cacao, café, coton, palmier à huile), adaptation climatique (variétés tolérantes sécheresse), jeunesse agricole.
\end{itemize}
\end{exemplebox}

\coursec{Bilan : les effets des RST sur le monde et le Cameroun}

\begin{aretenir}[Bilan synthétique : effets positifs et négatifs des RST]
\begin{center}
\begin{tabularx}{\linewidth}{|>{\bfseries}l|X|X|}\hline
\textbf{Domaine} & \textbf{Effets positifs (Progrès)} & \textbf{Effets négatifs / Risques (Défis)} \\ \hline
\textbf{Santé / Démographie} & Espérance de vie $\uparrow$, mortalité infantile $\downarrow$, éradication variole, contrôle maladies infectieuses, santé maternelle. & Résistance antibiotiques, pandémies (mobilité), vieillissement population, inégalités d'accès (Nord/Sud), bioéthique (CRISPR, GPA, fin de vie). \\ \hline
\textbf{Économie / Travail} & Productivité $\uparrow$, nouveaux métiers (IT, biotech, EnR), e-commerce, mobile money, télétravail, baisse coûts communication/transport. & Automatisation/IA $\rightarrow$ destruction emplois routiniers, précarisation (ubérisation), fracture numérique, concentration richesses (GAFAM, BATX), travail des enfants (minerais rares). \\ \hline
\textbf{Environnement / Climat} & Technologies bas carbone (EnR, nucléaire, batteries, efficacité énergétique), surveillance spatiale (climat, déforestation), agriculture de précision. & Réchauffement climatique (énergies fossiles), 6e extinction (habitats, pollution plastiques/chimiques), surexploitation ressources (eau, sols, minerais critiques), déchets nucléaires/électroniques. \\ \hline
\textbf{Géopolitique / Société} & Diffusion savoirs (Wikipédia, MOOCs), mobilisations citoyennes (réseaux sociaux), coopération spatiale (ISS), diplomatie scientifique. & Cyberguerre, désinformation, surveillance de masse, prolifération nucléaire, course aux armements (hypersoniques, espace), dépendance technologique (puces, cloud), fracture Nord/Sud. \\ \hline
\end{tabularx}
\end{center}
\end{aretenir}

\begin{savaistu}[Le Cameroun face aux RST : synthèse]
\begin{itemize}
\item \textbf{Atouts :} Jeunesse connectée (médiane 18 ans), ressources (minerais critiques : cobalt, lithium, terres rares), position géostratégique (golfe de Guinée), diaspora qualifiée.
\item \textbf{Leviers :} Stratégie \cle{Cameroun Numérique 2020/2030}, Zones économiques spéciales (Kribi, Limbé), formation ingénieurs (ENSP, Polytechnique, IAI), recherche (IRAD, IMPM, Universités).
\item \textbf{Vigilances :} Souveraineté numérique (données, cloud), transition écologique (forêt du bassin du Congo), inclusion sociale (femmes, zones rurales, handicap), gouvernance (corruption, instabilité NW/SW).
\end{itemize}
\end{savaistu}

\coursec{Fiches méthodes et exercices résolus}

\begin{methode}[Appliquer les notions de l'unité à une situation concrète de la vie courante]
Ce savoir-faire consiste à relier une grande découverte scientifique ou une innovation technique étudiée en classe à un usage quotidien (téléphone portable, vaccination, transport, énergie électrique). La démarche à suivre :
\begin{enumerate}
\item \textbf{Identifier} clairement la découverte ou l'innovation concernée (nom, date, acteur principal).
\item \textbf{Situer} cette innovation dans le contexte de l'accélération scientifique et technique après 1945 (guerre froide, course à l'espace, recherche médicale).
\item \textbf{Expliquer} le lien entre l'innovation et la situation concrète proposée : comment l'innovation transforme la vie quotidienne ?
\item \textbf{Illustrer} par un exemple précis, si possible camerounais (téléphonie mobile, barrages hydroélectriques, campagnes de vaccination).
\item \textbf{Conclure} en dégageant l'effet (positif ou négatif) de cette innovation sur la société.
\end{enumerate}
Réflexe à retenir : toujours répondre selon le schéma \cle{innovation} $\rightarrow$ \cle{usage concret} $\rightarrow$ \cle{conséquence sur la vie quotidienne}.
\end{methode}

\begin{exoresolu}[La téléphonie mobile au Cameroun : une application de la révolution des télécommunications]
Énoncé : À partir de vos connaissances sur les révolutions scientifiques et techniques depuis 1945, montrez comment la révolution des télécommunications transforme la vie quotidienne des Camerounais. Vous donnerez deux exemples concrets.
\tcblower
Solution détaillée :
\begin{enumerate}
\item \textbf{Identification de l'innovation.} La révolution des télécommunications repose sur les satellites artificiels (premier satellite \cle{Spoutnik} en 1957), les réseaux de téléphonie mobile et Internet (réseau \cle{Arpanet} en 1969, puis \cle{Web} développé par \cle{Tim Berners-Lee} au \cle{CERN} en 1989). Ces innovations appartiennent aux grandes innovations techniques de la période postérieure à 1945.
\item \textbf{Contexte.} Après la Deuxième Guerre mondiale, la rivalité entre les États-Unis et l'URSS (guerre froide) stimule la recherche spatiale et électronique. La miniaturisation des composants (transistor inventé en 1947, circuits intégrés, microprocesseur en 1971) rend possible la fabrication de téléphones portables à partir des années 1980.
\item \textbf{Lien avec la situation concrète.} Au Cameroun, la téléphonie mobile s'est diffusée massivement depuis les années 2000 (opérateurs \cle{MTN}, \cle{Orange}, \cle{Nexttel}). Elle s'appuie sur les réseaux cellulaires et les satellites de communication.
\item \textbf{Exemples concrets.}
\begin{itemize}
\item Le \emph{mobile money} (\cle{MTN Mobile Money}, \cle{Orange Money}) permet de payer, d'épargner et de transférer de l'argent sans compte bancaire, y compris dans les zones rurales : c'est une transformation directe de la vie économique quotidienne.
\item Internet mobile donne accès à l'information, à l'enseignement à distance (applications éducatives), aux réseaux sociaux et aux démarches administratives en ligne.
\end{itemize}
\item \textbf{Conclusion.} Issue de la course à l'espace et de la miniaturisation électronique après 1945, la révolution des télécommunications a profondément modifié la vie quotidienne des Camerounais : communication instantanée, inclusion financière, accès au savoir. Elle illustre parfaitement l'application des grandes innovations techniques à une situation concrète de la vie courante.
\end{enumerate}
\end{exoresolu}

\begin{methode}[Rédiger et justifier une démarche, une réponse ou une conclusion]
En Histoire, une réponse justifiée repose sur un raisonnement organisé et des preuves. La démarche à suivre :
\begin{enumerate}
\item \textbf{Analyser le sujet} : dégager les mots-clés, délimiter le sujet dans le temps (depuis 1945) et dans l'espace (monde, Cameroun).
\item \textbf{Mobiliser les connaissances} : lister les faits, dates, acteurs et notions du cours utiles au sujet.
\item \textbf{Construire un plan} : introduction (accroche, définition des termes, annonce du plan), développement en deux ou trois parties, conclusion.
\item \textbf{Justifier chaque affirmation} par un fait précis : date, nom de scientifique, événement (exemple : structure de l'ADN découverte par Watson et Crick en 1953).
\item \textbf{Rédiger une conclusion} qui répond exactement à la question posée et ouvre éventuellement sur une perspective.
\end{enumerate}
Réflexe à retenir : une affirmation sans \cle{preuve} (date, fait, acteur) ne vaut pas de point au Baccalauréat technique.
\end{methode}

\begin{exoresolu}[Sujet de dissertation : « La conquête de l'espace est-elle une conséquence de la guerre froide ? »]
Énoncé : Rédigez une introduction et une conclusion pour le sujet suivant : « La conquête de l'espace est-elle une conséquence de la guerre froide ? » Vous justifierez votre réponse par des faits précis.
\tcblower
Solution détaillée :

\textbf{Analyse du sujet.} Les mots-clés sont « conquête de l'espace » (exploration spatiale depuis 1957) et « guerre froide » (rivalité États-Unis/URSS de 1947 à 1991). Le sujet demande de justifier un lien de cause à effet.

\textbf{Introduction rédigée.} Après la Deuxième Guerre mondiale, le monde connaît une accélération sans précédent des découvertes scientifiques et des innovations techniques. Parmi elles, la conquête de l'espace occupe une place centrale : du premier satellite \cle{Spoutnik} en 1957 au premier homme sur la Lune en 1969, l'humanité franchit en douze ans des étapes décisives. Or cette épopée se déroule au cœur de la guerre froide, la rivalité idéologique, politique et technologique entre les États-Unis et l'URSS. On peut alors se demander dans quelle mesure la conquête de l'espace est une conséquence directe de cet affrontement. Nous verrons d'abord que la guerre froide est bien le moteur de la course à l'espace, puis que celle-ci dépasse progressivement ce cadre pour devenir une aventure scientifique et coopérative.

\textbf{Justifications attendues dans le développement (rappel).} \cle{Spoutnik} (1957) et \cle{Youri Gagarine}, premier homme dans l'espace (1961), sont des victoires de propagande soviétique ; la réponse américaine est le programme \cle{Apollo}, qui conduit \cle{Neil Armstrong} et \cle{Buzz Aldrin} sur la Lune le 21 juillet 1969. Les fusées spatiales dérivent des missiles balistiques militaires (V2 allemand récupéré après 1945).

\textbf{Conclusion rédigée.} La conquête de l'espace est donc bien, à l'origine, une conséquence de la guerre froide : c'est la compétition entre les deux supergrands qui finance et accélère les programmes spatiaux, chaque succès servant d'arme de propagande. Toutefois, elle débouche sur des retombées qui dépassent la rivalité Est-Ouest : satellites de télécommunications, d'observation météorologique et de navigation, coopérations internationales comme la mission \cle{Apollo-Soyouz} de 1975 puis la \cle{Station spatiale internationale}. Ainsi, née de la confrontation, l'aventure spatiale est devenue un instrument de progrès scientifique au service de toute l'humanité.
\end{exoresolu}

\begin{methode}[Construire une frise chronologique des grandes découvertes scientifiques]
Pour retenir et restituer les dates essentielles, la frise chronologique est un outil efficace. La démarche :
\begin{enumerate}
\item \textbf{Repérer} dans le cours les dates-clés des grandes découvertes (atome, espace, biologie).
\item \textbf{Classer} ces événements par ordre chronologique strict.
\item \textbf{Tracer} un axe gradué de 1945 à nos jours et y placer les événements avec un libellé court.
\item \textbf{Vérifier} la cohérence : chaque date doit correspondre au bon événement et au bon acteur.
\item \textbf{Exploiter} la frise pour répondre à une question : repérer les périodes d'accélération, relier les découvertes à leur contexte.
\end{enumerate}
Réflexe : au Bac, une date exacte associée au bon événement vaut des points ; une date approximative mal placée en fait perdre.
\end{methode}

\begin{exoresolu}[Frise des grandes découvertes scientifiques (1945-2003)]
Énoncé : Placez sur une frise chronologique les événements suivants, dans l'ordre : premier pas sur la Lune ; découverte de la structure de l'ADN ; première explosion atomique ; premier satellite artificiel ; séquençage du génome humain. Justifiez l'ordre retenu.
\tcblower
Solution détaillée :

\textbf{Classement chronologique.} On associe d'abord chaque événement à sa date :
\begin{itemize}
\item Première explosion atomique (essai \cle{Trinity}, États-Unis) : 16 juillet 1945 ;
\item Découverte de la structure en double hélice de l'ADN (\cle{Watson} et \cle{Crick}) : 1953 ;
\item Premier satellite artificiel (\cle{Spoutnik 1}, URSS) : 4 octobre 1957 ;
\item Premier pas sur la Lune (\cle{Neil Armstrong}, mission \cle{Apollo 11}) : 21 juillet 1969 ;
\item Achèvement du séquençage du génome humain (projet public \cle{Human Genome Project}) : 2003.
\end{itemize}

\textbf{Frise obtenue.}
\begin{popfigure}
\begin{tikzpicture}[scale=1]
\draw[->,thick,popdark] (0,0) -- (13.5,0);
\foreach \x/\y in {0.5/1945, 3.5/1953, 5.5/1957, 8.5/1969, 12.5/2003}{
  \draw[thick,popdark] (\x,-0.12) -- (\x,0.12);
  \node[below,popdark] at (\x,-0.15) {\y};
}
\node[above,align=center,popblue] at (0.5,0.2) {\footnotesize Bombe\\ \footnotesize atomique};
\node[above,align=center,poppurple] at (3.5,0.2) {\footnotesize Structure\\ \footnotesize de l'ADN};
\node[above,align=center,poporange] at (5.5,0.2) {\footnotesize Spoutnik 1};
\node[above,align=center,popgreen] at (8.5,0.2) {\footnotesize Homme sur\\ \footnotesize la Lune};
\node[above,align=center,poppink] at (12.5,0.2) {\footnotesize Génome\\ \footnotesize humain};
\end{tikzpicture}
\legende{Frise chronologique des grandes découvertes scientifiques de 1945 à 2003.}
\end{popfigure}

\textbf{Justification.} L'ordre retenu est strictement croissant : 1945 < 1953 < 1957 < 1969 < 2003. La maîtrise de l'atome ouvre la période (fin de la Deuxième Guerre mondiale) ; la biologie connaît sa révolution fondatrice en 1953 avec l'ADN ; la course à l'espace s'accélère entre 1957 et 1969 dans le contexte de la guerre froide ; la révolution génétique culmine en 2003 avec le séquençage complet du génome humain.

\textbf{Conclusion.} La frise montre une accélération continue des découvertes scientifiques depuis 1945, d'abord dominées par la physique (atome, espace), puis par la biologie (ADN, génome).
\end{exoresolu}

\begin{methode}[Commenter un document historique (texte ou image)]
Le commentaire de document est un exercice fréquent au Baccalauréat technique. La démarche en étapes :
\begin{enumerate}
\item \textbf{Présenter le document} : nature (texte, photographie, affiche), auteur, date, source, contexte.
\item \textbf{Décrire ou dégager les idées principales} : que voit-on ? que dit le texte ?
\item \textbf{Analyser} : dégager l'intention de l'auteur et la signification historique du document.
\item \textbf{Mettre en relation avec le cours} : relier le document aux notions étudiées (RST, guerre froide, décolonisation).
\item \textbf{Conclure} : apprécier l'intérêt et les limites du document comme source historique.
\end{enumerate}
Réflexe : ne jamais paraphraser le document ; il faut l'\cle{expliquer} à l'aide des connaissances du cours.
\end{methode}

\begin{exoresolu}[Commentaire d'une photographie : le premier pas sur la Lune (1969)]
Énoncé : On donne une photographie montrant l'astronaute Buzz Aldrin sur le sol lunaire, le 21 juillet 1969, lors de la mission Apollo 11. Présentez ce document, analysez sa signification historique et montrez son intérêt pour l'étude des révolutions scientifiques et techniques.
\tcblower
Solution détaillée :

\textbf{1. Présentation du document.} Il s'agit d'une photographie prise le 21 juillet 1969 par l'astronaute Neil Armstrong ; on y voit son coéquipier Buzz Aldrin en combinaison spatiale sur la surface de la Lune, devant le module lunaire de la mission américaine Apollo 11. C'est un document iconique diffusé dans le monde entier par la télévision.

\textbf{2. Description.} L'astronaute se déplace sur un sol gris et désertique, sans atmosphère ; le drapeau américain est planté à proximité. La combinaison spatiale, le module lunaire et les instruments scientifiques témoignent d'une technologie très avancée.

\textbf{3. Analyse.} Ce document illustre l'aboutissement de la course à l'espace engagée pendant la guerre froide : après les succès soviétiques (Spoutnik en 1957, Gagarine en 1961), les États-Unis remportent une victoire de propagande décisive en envoyant le premier homme sur la Lune. La photographie a une intention claire : démontrer la supériorité technologique américaine.

\textbf{4. Mise en relation avec le cours.} Apollo 11 résulte des grandes découvertes scientifiques (maîtrise des fusées issues de la recherche balistique, électronique embarquée) et des grandes innovations techniques (calculateurs, matériaux nouveaux, télécommunications permettant la retransmission en direct). Elle marque l'accélération scientifique et technique de l'après-1945.

\textbf{5. Conclusion.} Ce document est une source précieuse : il atteste d'un événement majeur du XXe siècle et de ses enjeux politiques. Sa limite est son caractère de propagande : il valorise un seul acteur (les États-Unis) et occulte la dimension collective et internationale de la recherche spatiale. Croisé avec d'autres sources, il reste un témoignage essentiel des révolutions scientifiques et techniques depuis 1945.
\end{exoresolu}

\begin{methode}[Exploiter des données chiffrées et calculer une évolution]
Certains sujets fournissent des statistiques (nombre d'internautes, production agricole, consommation d'énergie). La démarche :
\begin{enumerate}
\item \textbf{Lire} le tableau ou le graphique : titre, unités, dates, source.
\item \textbf{Repérer} les valeurs utiles au calcul demandé.
\item \textbf{Calculer} le taux de variation entre deux dates avec la formule :
\[ \text{Taux de variation (\%)} = \frac{V_{\text{finale}} - V_{\text{initiale}}}{V_{\text{initiale}}} \times 100 \]
\item \textbf{Interpréter} le résultat : forte ou faible croissance, en le reliant aux innovations techniques étudiées.
\item \textbf{Conclure} par une phrase qui relie le chiffre au cours.
\end{enumerate}
Réflexe : toujours écrire la formule, remplacer les valeurs, puis donner le résultat arrondi avec le symbole \%.
\end{methode}

\begin{exoresolu}[La diffusion d'Internet dans le monde]
Énoncé : Le nombre d'utilisateurs d'Internet dans le monde est passé d'environ 400 millions en 2000 à environ 4,9 milliards en 2021. Calculez le taux de variation entre 2000 et 2021, puis commentez le résultat en le reliant aux grandes innovations techniques.
\tcblower
Solution détaillée :

\textbf{1. Lecture des données.} Valeur initiale (2000) : $V_i = 400$ millions ; valeur finale (2021) : $V_f = 4\,900$ millions (4,9 milliards).

\textbf{2. Application de la formule.}
\begin{align*}
\text{Taux de variation} &= \frac{V_f - V_i}{V_i} \times 100 \\
&= \frac{4\,900 - 400}{400} \times 100 \\
&= \frac{4\,500}{400} \times 100 \\
&= 11,25 \times 100 \\
&= 1\,125\ \%
\end{align*}

\textbf{3. Interprétation.} Le nombre d'internautes a été multiplié par plus de 12 en vingt et un ans (une hausse de 1\,125 \%). Cette croissance spectaculaire s'explique par les grandes innovations techniques : invention du Web par Tim Berners-Lee en 1989, développement des ordinateurs personnels (microprocesseur de 1971), généralisation des téléphones intelligents et des réseaux mobiles, pose de câbles à fibre optique et lancement de satellites de communication.

\textbf{4. Conclusion.} Avec une progression de 1\,125 \% entre 2000 et 2021, Internet est devenu en une génération un outil planétaire : c'est l'illustration chiffrée de la révolution de l'informatique et des télécommunications, l'une des grandes innovations techniques de la période postérieure à 1945.
\end{exoresolu}

\begin{methode}[Rédiger un paragraphe argumenté sur les effets des RST]
Le paragraphe argumenté est une réponse courte (10 à 15 lignes) à une question précise. La démarche :
\begin{enumerate}
\item \textbf{Rédiger une phrase d'introduction} qui répond directement à la question (prise de position claire).
\item \textbf{Développer deux ou trois arguments}, chacun appuyé sur un fait précis (date, acteur, exemple).
\item \textbf{Nuancer} si nécessaire : effets positifs et effets négatifs (progrès médical mais bombe atomique ; productivité agricole mais pollution).
\item \textbf{Conclure} par une phrase de synthèse qui boucle le raisonnement.
\end{enumerate}
Réflexe : un paragraphe argumenté = une idée directrice + des preuves + une conclusion. Pas de liste de faits sans fil conducteur.
\end{methode}

\begin{exoresolu}[Les effets des révolutions scientifiques et techniques sur le monde et le Cameroun]
Énoncé : Rédigez un paragraphe argumenté répondant à la question : « Les révolutions scientifiques et techniques depuis 1945 ont-elles été bénéfiques pour le monde et pour le Cameroun ? »
\tcblower
Solution détaillée :

\textbf{Paragraphe rédigé.} Depuis 1945, les révolutions scientifiques et techniques ont profondément transformé le monde, avec des effets largement bénéfiques mais aussi des risques nouveaux, dont le Cameroun n'est pas exclu. Sur le plan positif, la révolution de la biologie et de la médecine a permis des progrès considérables : la découverte de la structure de l'ADN par Watson et Crick en 1953, le développement des vaccins (poliomyélite avec le vaccin de Salk en 1955), des antibiotiques et de l'imagerie médicale ont fait reculer les maladies et augmenté l'espérance de vie ; au Cameroun, les campagnes de vaccination et la lutte contre le paludisme en bénéficient directement. Les innovations techniques ont également transformé la vie quotidienne : l'informatique, Internet (Web créé en 1989) et la téléphonie mobile facilitent la communication, le commerce et l'accès au savoir, comme le montre l'essor du mobile money au Cameroun. Dans l'agriculture, la révolution verte (variétés améliorées, engrais, irrigation) a accru les rendements. Cependant, ces révolutions ont aussi des effets négatifs : la maîtrise de l'atome a conduit à la bombe atomique (Hiroshima et Nagasaki en 1945) et à la menace nucléaire ; l'industrialisation et la consommation d'énergies fossiles provoquent la pollution et le réchauffement climatique ; enfin, la fracture numérique creuse les inégalités entre pays riches et pays en développement. En définitive, les révolutions scientifiques et techniques depuis 1945 ont apporté des progrès majeurs au monde et au Cameroun, mais elles imposent un usage responsable pour limiter leurs dangers.

\textbf{Analyse de la copie.} Le paragraphe respecte la méthode : phrase d'introduction qui répond et nuance (« largement bénéfiques mais aussi des risques »), trois arguments justifiés par des faits datés (médecine, télécommunications, agriculture), une nuance (bombe atomique, pollution, fracture numérique) et une conclusion qui synthétise. Chaque affirmation est appuyée par une preuve : c'est ce qui rapporte des points au Baccalauréat technique.
\end{exoresolu}

\begin{attention}[Les 5 erreurs qui coûtent des points au Bac]
\begin{enumerate}
\item \textbf{Confondre les dates et les acteurs.} Ne pas attribuer la découverte de l'ADN (1953, Watson et Crick) à un autre savant, ni placer le premier pas sur la Lune (1969) avant Spoutnik (1957). Une frise chronologique fausse ruine toute la démonstration.
\item \textbf{Affirmer sans justifier.} Écrire « la guerre froide a provoqué la course à l'espace » sans citer Spoutnik, Gagarine ou Apollo 11 ne vaut rien : chaque affirmation doit être appuyée par un fait précis, daté et nommé.
\item \textbf{Paraphraser le document au lieu de le commenter.} Recopier le texte ou décrire l'image sans la relier au cours est la faute la plus fréquente : il faut toujours expliquer la signification historique du document à l'aide des notions de l'unité.
\item \textbf{Se tromper dans le calcul du taux de variation.} Erreurs classiques : diviser par la valeur finale au lieu de la valeur initiale, oublier de multiplier par 100, ou confondre le taux de variation avec la différence simple. Toujours écrire la formule avant de calculer.
\item \textbf{Sortir du sujet et de la période.} Traiter d'inventions antérieures à 1945 (machine à vapeur, télégraphe) ou rédiger une conclusion qui ne répond pas à la question posée : il faut rester dans le cadre « depuis la fin de la Deuxième Guerre mondiale » et conclure exactement sur le problème posé.
\end{enumerate}
\end{attention}

\newpage

\coursec{Exercices}

\courssub{Je vérifie mes connaissances}

\begin{exercice}[QCM : repérer les grandes étapes]
Pour chaque question, une seule réponse est exacte. Recopie la lettre de la bonne réponse.
\begin{enumerate}
\item La structure de l'ADN (double hélice) est découverte en 1953 par :
\begin{enumerate}
\item A. Fleming et Chain \quad B. Watson et Crick \quad C. Salk et Sabin
\end{enumerate}
\item Le premier homme à marcher sur la Lune, le 21 juillet 1969, est :
\begin{enumerate}
\item A. Youri Gagarine \quad B. Neil Armstrong \quad C. John Glenn
\end{enumerate}
\item La première utilisation militaire de la bombe atomique a lieu en :
\begin{enumerate}
\item A. 1939 \quad B. 1945 \quad C. 1962
\end{enumerate}
\item Le réseau Internet devient accessible au grand public dans les années :
\begin{enumerate}
\item A. 1960 \quad B. 1980 \quad C. 1990
\end{enumerate}
\end{enumerate}
\end{exercice}

\begin{exercice}[Vrai ou faux, en justifiant]
Pour chaque affirmation, réponds par \textbf{vrai} ou \textbf{faux}, puis justifie ta réponse en une ou deux phrases.
\begin{enumerate}
\item La pénicilline a été découverte après la Deuxième Guerre mondiale.
\item Youri Gagarine est le premier homme envoyé dans l'espace, en 1961.
\item La révolution verte a permis d'augmenter les rendements agricoles grâce à de nouvelles variétés de semences.
\item Le téléphone mobile s'est diffusé au Cameroun dès les années 1970.
\end{enumerate}
\end{exercice}

\begin{exercice}[Définitions]
Définis en deux ou trois phrases chacune des notions suivantes, en donnant pour chacune un exemple daté :
\begin{enumerate}
\item \cle{Révolution scientifique et technique} (RST) ;
\item \cle{Énergie nucléaire} ;
\item \cle{Révolution verte} ;
\item \cle{Mondialisation numérique}.
\end{enumerate}
\end{exercice}

\begin{exercice}[Calculs directs : lire une évolution]
Au Cameroun, le nombre d'abonnés au téléphone mobile est passé d'environ 2,3 millions en 2005 à environ 19,5 millions en 2019, pour une population d'environ 25,9 millions d'habitants en 2019.
\begin{enumerate}
\item Calcule l'augmentation du nombre d'abonnés entre 2005 et 2019.
\item Calcule le taux de variation en pourcentage entre 2005 et 2019 (arrondi à l'unité).
\item Calcule le taux de pénétration du mobile en 2019 (nombre d'abonnés pour 100 habitants), arrondi au dixième.
\end{enumerate}
\end{exercice}

\courssub{Je m'entraîne}

\begin{exercice}[Frise chronologique]
Replace sur une frise chronologique (1945--2000) les événements suivants, en précisant pour chacun le domaine concerné (atome, espace, médecine, informatique) : bombardement d'Hiroshima ; découverte de la structure de l'ADN ; premier vol spatial habité de Gagarine ; premier pas sur la Lune ; mise au point du vaccin contre la poliomyélite par Salk ; création du réseau ARPANET ; naissance du premier « bébé-éprouvette » (1978).
\end{exercice}

\begin{exercice}[Classer : découvertes ou innovations ?]
Recopie et complète le tableau ci-dessous en classant chaque élément dans la bonne colonne et en indiquant son domaine :
\begin{center}
\begin{tabularx}{\linewidth}{|l|X|X|}
\hline
\rowcolor{popblueL} \textbf{Élément} & \textbf{Découverte scientifique} & \textbf{Innovation technique} \\
\hline
Structure de l'ADN (1953) & & \\
\hline
Téléphone mobile & & \\
\hline
Vaccin antipoliomyélitique & & \\
\hline
Internet & & \\
\hline
Fission nucléaire maîtrisée & & \\
\hline
Ordinateur personnel & & \\
\hline
\end{tabularx}
\end{center}
\end{exercice}

\begin{exercice}[Construction graphique : l'espérance de vie au Cameroun]
Voici l'évolution de l'espérance de vie à la naissance au Cameroun (en années) :

\begin{center}
\begin{tabularx}{\linewidth}{|l|X|X|X|X|X|X|X|}
\hline
\rowcolor{popblueL} \textbf{Année} & 1960 & 1970 & 1980 & 1990 & 2000 & 2010 & 2020 \\
\hline
\textbf{Espérance de vie} & 42 & 46 & 51 & 54 & 52 & 56 & 60 \\
\hline
\end{tabularx}
\end{center}

\begin{enumerate}
\item Construis un diagramme en barres représentant cette évolution (1 cm pour 10 ans sur l'axe horizontal, 1 cm pour 5 années d'espérance de vie sur l'axe vertical).
\item Décris l'évolution générale entre 1960 et 2020, en relevant le léger recul autour de 2000.
\item Explique cette évolution par deux progrès médicaux étudiés dans la leçon (vaccination, antibiotiques, hygiène, imagerie médicale).
\end{enumerate}
\end{exercice}

\begin{exercice}[Question à réponse construite]
En une vingtaine de lignes, présente \textbf{trois grandes découvertes scientifiques} et \textbf{trois grandes innovations techniques} réalisées depuis 1945. Pour chacune, tu préciseras : la date, le domaine (physique, biologie, médecine, informatique, transports...) et une conséquence concrète sur la vie des hommes.
\end{exercice}

\courssub{Je me prépare au Bac}

\begin{exercice}[Dissertation type Baccalauréat technique]
\textbf{Sujet :} « Les révolutions scientifiques et techniques depuis 1945 : progrès ou menaces pour l'humanité ? »

Tu rédigeras une dissertation complète comportant :
\begin{enumerate}
\item une \textbf{introduction} (accroche, présentation du sujet, définition des RST, problématique, annonce du plan) ;
\item un \textbf{développement en deux parties} : première partie montrant que les RST constituent un progrès pour l'humanité (médecine, agriculture, communications, conquête spatiale) ; deuxième partie montrant leurs dangers et limites (arme nucléaire, pollutions, inégalités d'accès, dérives du numérique) ;
\item une \textbf{conclusion} répondant à la problématique et s'ouvrant sur l'avenir des sciences et techniques en Afrique et au Cameroun.
\end{enumerate}
Chaque argument devra être illustré par au moins un exemple précis et daté.
\end{exercice}

\begin{exercice}[Étude critique de documents]
\textbf{Document 1 :} Photographie (juillet 1969) : l'astronaute américain Buzz Aldrin, membre de la mission Apollo 11, pose le pied sur le sol lunaire ; le drapeau des États-Unis est planté à côté du module.

\textbf{Document 2 :} Abonnés au téléphone mobile au Cameroun (en millions) :

\begin{center}
\begin{tabularx}{\linewidth}{|l|X|X|X|X|X|}
\hline
\rowcolor{popblueL} \textbf{Année} & 2000 & 2005 & 2010 & 2015 & 2019 \\
\hline
\textbf{Abonnés (millions)} & 0,2 & 2,3 & 8,2 & 15,4 & 19,5 \\
\hline
\end{tabularx}
\end{center}

\textbf{Questions :}
\begin{enumerate}
\item Présente le document 1 (nature, date, auteur ou contexte, contenu).
\item À quelle grande étape des révolutions scientifiques et techniques ce document renvoie-t-il ? Cite deux autres étapes de la conquête de l'espace.
\item Décris l'évolution du nombre d'abonnés au mobile au Cameroun entre 2000 et 2019, en calculant le taux de variation global.
\item À quelle innovation technique cette évolution est-elle liée ? Rappelle ses grandes étapes (ordinateur, ARPANET, Internet, mobile).
\item Montre, par trois exemples concrets, comment les RST ont transformé la vie quotidienne des Camerounais (communication, commerce, santé, éducation).
\end{enumerate}
\end{exercice}

\begin{exercice}[Débat argumenté type Probatoire]
\textbf{Sujet :} « Le téléphone mobile est-il la plus grande innovation technique pour l'Afrique depuis 1945 ? »

\begin{enumerate}
\item Rédige une prise de position claire (oui ou non), annoncée dès l'introduction.
\item Justifie ta position par \textbf{trois arguments}, chacun appuyé sur un exemple camerounais ou africain précis : par exemple le mobile money (MTN Mobile Money, Orange Money), l'accès à l'information et à l'éducation, la santé à distance, ou au contraire la médecine (vaccins), l'énergie ou la révolution verte.
\item Intègre au moins un argument contraire au tien et réfute-le.
\item Conclus en proposant une condition pour que l'innovation technique profite à tous les Africains.
\end{enumerate}
Ton texte comportera entre vingt et trente lignes, rédigé et justifié.
\end{exercice}

\newpage

\coursec{Corrigés des exercices}

\coursec{Corrigés des activités — Leçon 05 : Les révolutions scientifiques et techniques depuis 1945}

\begin{corrige}[Exercice 1 — QCM : repérer les grandes étapes]
\begin{enumerate}
\item \textbf{B. Watson et Crick.} En 1953, James Watson et Francis Crick décrivent la structure en double hélice de l'ADN (avec les travaux de Rosalind Franklin). Fleming et Chain sont liés à la pénicilline ; Salk et Sabin aux vaccins contre la poliomyélite.
\item \textbf{B. Neil Armstrong.} Le 21 juillet 1969, lors de la mission Apollo 11, Neil Armstrong est le premier homme à marcher sur la Lune. Youri Gagarine est le premier homme dans l'espace (12 avril 1961), sans alunir.
\item \textbf{B. 1945.} Les 6 et 9 août 1945, les États-Unis bombardent Hiroshima et Nagasaki : c'est la première (et seule) utilisation militaire de la bombe atomique.
\item \textbf{C. 1990.} Né d'ARPANET (1969), le réseau Internet s'ouvre au grand public dans les années 1990, notamment grâce au Web inventé par Tim Berners-Lee en 1989.
\end{enumerate}
\textbf{Points de vigilance :} ne pas confondre Gagarine (premier homme dans l'espace) et Armstrong (premier homme sur la Lune) ; ne pas confondre la date de la bombe atomique (1945) avec la crise de Cuba (1962).
\end{corrige}

\begin{corrige}[Exercice 2 — Vrai ou faux, en justifiant]
\begin{enumerate}
\item \textbf{Faux.} La pénicilline est découverte par Alexander Fleming en 1928, donc \emph{avant} la Deuxième Guerre mondiale ; elle est toutefois produite industriellement pendant la guerre (à partir de 1943) pour soigner les blessés.
\item \textbf{Vrai.} Le 12 avril 1961, le Soviétique Youri Gagarine effectue le premier vol spatial habité à bord de Vostok 1 : il réalise un tour de la Terre en 108 minutes.
\item \textbf{Vrai.} La révolution verte (années 1960-1970) repose sur des variétés de blé et de riz à haut rendement (travaux de Norman Borlaug), associées aux engrais et à l'irrigation, ce qui a fortement augmenté les récoltes, notamment au Mexique et en Inde.
\item \textbf{Faux.} Le téléphone mobile ne se diffuse au Cameroun qu'à partir de la fin des années 1990 et surtout des années 2000 (environ 0,2 million d'abonnés en 2000, 19,5 millions en 2019). Dans les années 1970, seul le téléphone fixe, très limité, existe.
\end{enumerate}
\textbf{Points de vigilance :} une justification datée est exigée pour chaque item ; « vrai » ou « faux » seul ne rapporte pas de point.
\end{corrige}

\begin{corrige}[Exercice 3 — Définitions]
\begin{enumerate}
\item \textbf{Révolution scientifique et technique (RST) :} ensemble des transformations rapides et profondes des connaissances scientifiques et de leurs applications techniques depuis 1945, qui bouleversent la production, les communications et la vie quotidienne. \emph{Exemple :} la mise au point du premier ordinateur électronique ENIAC en 1946.
\item \textbf{Énergie nucléaire :} énergie libérée par la fission des noyaux d'atomes lourds (uranium), utilisée pour produire de l'électricité dans les centrales, mais aussi à des fins militaires. \emph{Exemple :} la première centrale nucléaire civile d'Obninsk (URSS, 1954).
\item \textbf{Révolution verte :} ensemble de techniques agricoles modernes (semences améliorées à haut rendement, engrais, pesticides, irrigation, mécanisation) mises en œuvre à partir des années 1960 pour augmenter la production alimentaire. \emph{Exemple :} les variétés de blé de Norman Borlaug diffusées en Inde à partir de 1965.
\item \textbf{Mondialisation numérique :} circulation instantanée et planétaire des informations, des biens culturels et financiers permise par l'informatique, Internet et les télécommunications. \emph{Exemple :} l'ouverture d'Internet au grand public dans les années 1990 et l'essor du mobile money en Afrique dans les années 2010.
\end{enumerate}
\textbf{Points de vigilance :} chaque définition doit comporter une date et un exemple précis, comme demandé par la consigne.
\end{corrige}

\begin{corrige}[Exercice 4 — Calculs directs : lire une évolution]
\begin{enumerate}
\item Augmentation du nombre d'abonnés entre 2005 et 2019 :
\[ 19,5 - 2,3 = 17,2 \text{ millions d'abonnés.} \]
\item Taux de variation entre 2005 et 2019 :
\[ \frac{19,5 - 2,3}{2,3} \times 100 = \frac{17,2}{2,3} \times 100 \approx 748 \%. \]
Le nombre d'abonnés a donc été multiplié par environ $8,5$ (soit une hausse d'environ 748 \%).
\item Taux de pénétration en 2019 (population estimée à 25,9 millions) :
\[ \frac{19,5}{25,9} \times 100 \approx 75,3 \text{ abonnés pour 100 habitants.} \]
\end{enumerate}
\textbf{Points de vigilance :} bien distinguer le taux de variation (748 \%) du coefficient multiplicateur (8,5) ; le taux de pénétration peut dépasser 100 \% dans certains pays car une même personne possède parfois plusieurs cartes SIM — ici il reste inférieur à 100 \%.
\end{corrige}

\begin{corrige}[Exercice 5 — Frise chronologique]
Ordre chronologique attendu, avec le domaine de chaque événement :
\begin{enumerate}
\item \textbf{1945 :} bombardement d'Hiroshima — domaine de l'\cle{atome} ;
\item \textbf{1953 :} découverte de la structure de l'ADN — domaine de la \cle{médecine}/biologie ;
\item \textbf{1955 :} vaccin contre la poliomyélite mis au point par Jonas Salk — domaine de la \cle{médecine} ;
\item \textbf{1961 :} premier vol spatial habité de Gagarine — domaine de l'\cle{espace} ;
\item \textbf{1969 :} premier pas sur la Lune (Apollo 11) et création du réseau ARPANET — domaines de l'\cle{espace} et de l'\cle{informatique} ;
\item \textbf{1978 :} naissance du premier « bébé-éprouvette » (Louise Brown, Grande-Bretagne) — domaine de la \cle{médecine}.
\end{enumerate}
\begin{popfigure}
\begin{tikzpicture}[scale=1.05]
\draw[->, thick, popdark] (0,0) -- (12.6,0) node[right]{\footnotesize temps};
\foreach \x/\y/\t/\d in {0.3/1945/Hiroshima/atome, 2.2/1953/ADN/médecine, 3.7/1955/Vaccin polio (Salk)/médecine, 5.4/1961/Gagarine/espace, 7.4/1969/Lune + ARPANET/espace--info., 9.6/1978/Bébé-éprouvette/médecine}{
  \draw[popblue, thick] (\x,0.08) -- (\x,-0.08);
  \node[below, font=\scriptsize, popdark] at (\x,-0.08) {\y};
  \node[above, font=\scriptsize, align=center, popink] at (\x,0.1) {\t\\\emph{\d}};
}
\end{tikzpicture}
\legende{Frise chronologique des grandes étapes des RST (1945--1978) : atome, médecine, espace, informatique.}
\end{popfigure}
\textbf{Points de vigilance :} l'année 1969 compte deux événements (Lune et ARPANET) ; respecter l'échelle régulière de la frise et indiquer le domaine pour chaque fait.
\end{corrige}

\begin{corrige}[Exercice 6 — Classer : découvertes ou innovations ?]
\begin{center}
\begin{tabularx}{\linewidth}{|l|X|X|}
\hline
\rowcolor{popblueL} \textbf{Élément} & \textbf{Découverte scientifique} & \textbf{Innovation technique} \\
\hline
Structure de l'ADN (1953) & Oui (biologie) & \\
\hline
Téléphone mobile & & Oui (télécommunications, années 1980-1990) \\
\hline
Vaccin antipoliomyélitique (1955) & Oui (médecine) & \\
\hline
Internet & & Oui (informatique/réseaux, années 1990) \\
\hline
Fission nucléaire maîtrisée (1942-1945) & Oui (physique) & \\
\hline
Ordinateur personnel (années 1970-1980) & & Oui (informatique) \\
\hline
\end{tabularx}
\end{center}
\textbf{Justification du critère :} une \cle{découverte scientifique} révèle une connaissance nouvelle sur la nature (structure de la matière, du vivant) ; une \cle{innovation technique} est la mise au point d'un objet ou d'un procédé utilisable par les hommes. Certains éléments relèvent des deux (le vaccin est aussi une innovation médicale) : on accepte la double appartenance si elle est justifiée.
\textbf{Points de vigilance :} la fission nucléaire est d'abord une découverte de physique (Hahn, 1938 ; première pile de Fermi, 1942) ; la bombe et la centrale sont ses applications techniques.
\end{corrige}

\begin{corrige}[Exercice 7 — Construction graphique : l'espérance de vie au Cameroun]
\begin{enumerate}
\item Diagramme en barres attendu (échelle : 1 cm pour 10 ans en abscisse, 1 cm pour 5 ans d'espérance de vie en ordonnée) :
\end{enumerate}
\begin{popfigure}
\begin{tikzpicture}
\begin{axis}[
  ybar, bar width=0.55cm,
  width=0.95\linewidth, height=6.5cm,
  ymin=0, ymax=70,
  ylabel={Espérance de vie (années)},
  xlabel={Année},
  symbolic x coords={1960,1970,1980,1990,2000,2010,2020},
  xtick=data,
  nodes near coords, every node near coord/.append style={font=\scriptsize},
  ymajorgrids=true, grid style={popline!40},
]
\addplot[fill=popblue!70, draw=popblue] coordinates {(1960,42) (1970,46) (1980,51) (1990,54) (2000,52) (2010,56) (2020,60)};
\end{axis}
\end{tikzpicture}
\legende{Espérance de vie à la naissance au Cameroun, 1960--2020 (en années). Source : données Banque mondiale, valeurs arrondies.}
\end{popfigure}
\begin{enumerate}
\setcounter{enumi}{1}
\item \textbf{Description :} l'espérance de vie augmente globalement, passant de 42 ans en 1960 à 60 ans en 2020, soit un gain de 18 ans en soixante ans. La hausse est régulière, avec toutefois un \textbf{léger recul autour de 2000} (de 54 ans en 1990 à 52 ans en 2000), lié notamment à l'épidémie de sida et aux difficultés sanitaires et économiques des années 1990.
\item \textbf{Explication :} cette progression s'explique d'abord par les \textbf{campagnes de vaccination} (poliomyélite, rougeole, fièvre jaune), qui ont fait reculer les maladies infantiles ; ensuite par la diffusion des \textbf{antibiotiques} et l'amélioration de l'\textbf{hygiène} (eau potable, assainissement), qui réduisent les infections. L'imagerie médicale et les progrès du diagnostic complètent ces avancées.
\end{enumerate}
\textbf{Points de vigilance :} axes gradués et légendés obligatoires ; la baisse de 2000 doit être explicitement relevée et expliquée, pas seulement constatée.
\end{corrige}

\begin{corrige}[Exercice 8 — Question à réponse construite]
\textbf{Corrigé rédigé (modèle de réponse attendue).}

Depuis 1945, les révolutions scientifiques et techniques ont transformé le monde. On peut citer trois grandes découvertes scientifiques et trois grandes innovations techniques.

Parmi les découvertes, la \textbf{maîtrise de la fission nucléaire} (physique, 1942-1945, travaux de Fermi puis projet Manhattan) a permis de produire une énergie considérable, utilisée dans les centrales électriques, mais aussi dans l'arme atomique. La \textbf{découverte de la structure de l'ADN} (biologie, 1953, Watson et Crick) a ouvert la voie à la génétique moderne : dépistage des maladies héréditaires, empreintes génétiques, biotechnologies. La mise au point du \textbf{vaccin contre la poliomyélite} (médecine, 1955, Jonas Salk) a permis de quasiment éliminer cette maladie paralysante dans le monde, sauvant des millions d'enfants.

Parmi les innovations techniques, l'\textbf{ordinateur} (informatique, ENIAC en 1946, puis ordinateur personnel dans les années 1970-1980) a automatisé les calculs et le traitement de l'information dans tous les secteurs. \textbf{Internet} (informatique, ARPANET en 1969, Web de Berners-Lee en 1989, ouvert au public dans les années 1990) a rendu l'information accessible instantanément à l'échelle planétaire. Enfin, le \textbf{téléphone mobile} (télécommunications, diffusé à partir des années 1990-2000) a révolutionné la communication : au Cameroun, les abonnés sont passés de 0,2 million en 2000 à 19,5 millions en 2019, facilitant les échanges, le commerce et le mobile money.

Ainsi, ces découvertes et innovations, réalisées en moins de soixante ans, ont profondément modifié la santé, l'énergie et les communications des hommes.

\textbf{Barème indicatif (sur 20) :} 3 découvertes datées et situées dans leur domaine (6 pts) ; 3 innovations datées et situées (6 pts) ; une conséquence concrète pour chacune (6 pts) ; cohérence et qualité de la rédaction (2 pts).
\textbf{Points de vigilance :} chaque élément cité doit comporter une date, un domaine et une conséquence ; ne pas confondre découverte (connaissance) et innovation (application technique).
\end{corrige}

\begin{corrige}[Exercice 9 — Dissertation type Baccalauréat technique]
\textbf{Introduction (modèle).} En juillet 1969, plus de 500 millions d'hommes regardent en direct Neil Armstrong marcher sur la Lune : image symbolique de la puissance acquise par la science depuis 1945. Depuis la fin de la Deuxième Guerre mondiale, l'humanité connaît une accélération sans précédent des révolutions scientifiques et techniques (RST), c'est-à-dire l'ensemble des découvertes fondamentales (atome, ADN, vaccins) et de leurs applications techniques (ordinateur, Internet, mobile) qui transforment la vie des sociétés. Mais ces progrès sont ambivalents : la même physique nucléaire éclaire les villes et a détruit Hiroshima. Dès lors, les RST depuis 1945 constituent-elles un progrès ou une menace pour l'humanité ? Nous montrerons d'abord qu'elles ont amélioré la condition humaine, puis qu'elles comportent des dangers réels qu'il faut maîtriser.

\textbf{Première partie : les RST, un progrès pour l'humanité.}
\begin{itemize}
\item \emph{Médecine :} les antibiotiques (pénicilline produite industriellement dès 1943), les vaccins (poliomyélite, Salk, 1955) et l'imagerie médicale ont fait reculer les maladies : l'espérance de vie au Cameroun passe de 42 ans en 1960 à 60 ans en 2020.
\item \emph{Agriculture :} la révolution verte (années 1960, semences de Borlaug) a multiplié les rendements et évité des famines en Inde et au Mexique.
\item \emph{Communications :} l'ordinateur (1946), Internet (années 1990) et le mobile (19,5 millions d'abonnés au Cameroun en 2019) rapprochent les hommes et dynamisent le commerce (mobile money).
\item \emph{Espace :} Gagarine (1961), Apollo 11 (1969) et les satellites ont permis les télécommunications, la météorologie et la navigation par GPS.
\end{itemize}

\textbf{Deuxième partie : les RST, une menace et des limites.}
\begin{itemize}
\item \emph{Arme nucléaire :} Hiroshima et Nagasaki (août 1945, plus de 200 000 morts) ; la course aux armements de la Guerre froide (crise de Cuba, 1962) a fait peser la menace d'une destruction de l'humanité.
\item \emph{Pollutions :} accidents nucléaires (Tchernobyl, 1986), pesticides de la révolution verte, réchauffement climatique lié aux énergies fossiles et à l'industrie.
\item \emph{Inégalités d'accès :} fracture numérique entre pays riches et pays pauvres, et au sein du Cameroun entre villes et campagnes.
\item \emph{Dérives du numérique :} surveillance, fausses informations, dépendance aux écrans ; questions éthiques de la génétique (bébé-éprouvette, 1978 ; clonage, brebis Dolly, 1996).
\end{itemize}

\textbf{Conclusion (modèle).} Les RST depuis 1945 sont donc à la fois un progrès indéniable — elles soignent, nourrissent et connectent les hommes — et une menace lorsqu'elles échappent au contrôle des sociétés. Tout dépend de l'usage qu'en font les hommes et les États. L'avenir de l'Afrique et du Cameroun dépendra de leur capacité à former des scientifiques, à développer leurs propres technologies (start-up numériques, énergies renouvelables) et à faire profiter tous leurs citoyens de ces progrès.

\textbf{Barème indicatif (sur 20) :} introduction complète (accroche, définition, problématique, plan) : 4 pts ; première partie argumentée avec exemples datés : 6 pts ; deuxième partie argumentée avec exemples datés : 6 pts ; conclusion répondant à la problématique avec ouverture : 2 pts ; expression et organisation : 2 pts.
\textbf{Points de vigilance :} la problématique doit être une question explicite ; chaque argument exige au moins un exemple précis et daté ; éviter le catalogue de faits sans analyse ; la conclusion doit trancher nuancément (progrès ET menaces selon les usages).
\end{corrige}

\begin{corrige}[Exercice 10 — Étude critique de documents]
\begin{enumerate}
\item \textbf{Présentation du document 1 :} il s'agit d'une photographie prise en juillet 1969 sur la Lune, lors de la mission américaine Apollo 11. On y voit l'astronaute Buzz Aldrin posant le pied sur le sol lunaire, à côté du module et du drapeau des États-Unis. Ce document est une source primaire, témoignage direct d'un événement historique mondial retransmis en direct à la télévision.
\item \textbf{Étape des RST :} ce document renvoie à la \cle{conquête de l'espace}, grande étape des révolutions scientifiques et techniques de l'après-1945. Deux autres étapes : le premier satellite artificiel Spoutnik 1 (URSS, 1957) et le premier vol spatial habité de Youri Gagarine (12 avril 1961). On peut aussi citer les stations spatiales et les satellites de télécommunications.
\item \textbf{Évolution des abonnés au mobile :} le nombre d'abonnés passe de 0,2 million en 2000 à 19,5 millions en 2019, soit une augmentation de $19,5 - 0,2 = 19,3$ millions. Taux de variation global :
\[ \frac{19,5 - 0,2}{0,2} \times 100 = \frac{19,3}{0,2} \times 100 = 9650 \%. \]
Le nombre d'abonnés a donc été multiplié par environ $97,5$ en dix-neuf ans : une croissance spectaculaire, particulièrement rapide entre 2005 et 2015.
\item \textbf{Innovation technique liée :} cette évolution est liée à la révolution de l'\cle{informatique et des télécommunications}. Ses grandes étapes : le premier ordinateur électronique ENIAC (1946) ; le réseau ARPANET (1969) ; l'ordinateur personnel (années 1970-1980) ; le Web de Tim Berners-Lee (1989) et l'ouverture d'Internet au public (années 1990) ; enfin la diffusion du téléphone mobile et des smartphones (années 1990-2000).
\item \textbf{Transformations de la vie quotidienne au Cameroun :}
\begin{itemize}
\item \emph{Communication :} les familles restent en contact entre villes et villages grâce au mobile, alors que le téléphone fixe était quasi inexistant ;
\item \emph{Commerce :} le mobile money (MTN Mobile Money, Orange Money) permet de payer, d'épargner et de transférer de l'argent sans compte bancaire ;
\item \emph{Santé et éducation :} prise de rendez-vous médicaux par téléphone, campagnes de sensibilisation sanitaire par SMS, accès des élèves et étudiants aux cours et informations en ligne.
\end{itemize}
\end{enumerate}
\textbf{Barème indicatif (sur 20) :} présentation du document (4 pts) ; étape des RST et deux exemples datés (4 pts) ; description chiffrée et calcul exact (4 pts) ; innovation et étapes (4 pts) ; trois exemples concrets camerounais (4 pts).
\textbf{Points de vigilance :} dans la présentation d'un document iconographique, toujours annoncer nature, date, contexte et contenu ; le taux de variation doit être calculé avec la formule et non estimé ; les exemples de la question 5 doivent être camerounais et concrets.
\end{corrige}

\begin{corrige}[Exercice 11 — Débat argumenté type Probatoire]
\textbf{Corrigé rédigé (modèle de prise de position positive).}

Depuis 1945, l'Afrique a bénéficié de nombreuses innovations techniques : vaccins, barrages, semences améliorées. Mais aucune ne s'est diffusée aussi vite et aussi largement que le téléphone mobile. J'estime donc que le mobile est bien la plus grande innovation technique pour l'Afrique depuis 1945, car il a transformé l'économie, l'accès à l'information et les services quotidiens.

Premier argument : le mobile a révolutionné l'économie africaine grâce au \emph{mobile money}. Au Cameroun, MTN Mobile Money et Orange Money permettent depuis les années 2010 à des millions de personnes sans compte bancaire de payer, d'épargner et de recevoir de l'argent : un commerçant de Douala peut être payé instantanément par un client de Ngaoundéré. Avec 19,5 millions d'abonnés en 2019 contre 0,2 million en 2000, le mobile est devenu l'outil économique le plus répandu du pays.

Deuxième argument : le mobile donne accès à l'\emph{information et à l'éducation}. Grâce aux smartphones et à Internet mobile, les élèves camerounais consultent des cours en ligne, les agriculteurs connaissent les prix des marchés avant de vendre leur cacao ou leur maïs, et les citoyens s'informent sans attendre la radio ou la télévision d'État.

Troisième argument : le mobile améliore la \emph{santé}. Les campagnes de vaccination sont annoncées par SMS, les malades des zones rurales contactent un médecin à distance, et les rappels de rendez-vous médicaux réduisent l'abandon des traitements.

On objectera que la médecine, avec les vaccins (poliomyélite, 1955) ou les antibiotiques, a sauvé plus de vies que le téléphone. Cet argument est sérieux : l'espérance de vie au Cameroun est passée de 42 ans en 1960 à 60 ans en 2020 grâce aux progrès médicaux. Mais on peut le réfuter : le mobile est précisément devenu l'instrument qui diffuse ces progrès médicaux (information sanitaire, télémédecine), et il agit en même temps sur l'économie, l'éducation et la vie sociale, ce qu'aucune autre innovation ne fait à cette échelle.

En conclusion, le téléphone mobile est bien la plus grande innovation technique pour l'Afrique depuis 1945, par sa rapidité de diffusion et la variété de ses usages. Toutefois, pour qu'il profite à tous les Africains, il faut une condition essentielle : réduire la fracture numérique en électrifiant les campagnes, en baissant le coût des données et en formant tous les citoyens, notamment les jeunes, aux outils numériques.

\textbf{Barème indicatif (sur 20) :} prise de position claire dès l'introduction (3 pts) ; trois arguments étayés par des exemples africains ou camerounais précis (9 pts) ; argument contraire intégré et réfuté (4 pts) ; conclusion avec condition formulée (2 pts) ; qualité de l'expression et respect de la longueur (2 pts).
\textbf{Points de vigilance :} la position (oui ou non) doit être annoncée explicitement et assumée jusqu'au bout ; un argument contraire non réfuté fait perdre des points ; les exemples doivent être précis et situés (noms, dates, chiffres) ; une position « non » correctement argumentée (vaccins, énergie, révolution verte) est tout aussi valable.
\end{corrige}

\newpage

\coursec{Fiche bilan}

\begin{popfigure}
\begin{center}
\begin{tikzpicture}[mindmap, grow cyclic, every node/.style={concept, text=white, font=\small\bfseries},
root concept/.append style={concept color=popdark, font=\small\bfseries},
level 1/.append style={level distance=3.4cm, sibling angle=60, font=\scriptsize},
level 2/.append style={level distance=2.1cm, sibling angle=45, font=\tiny}]
\node[root concept]{Révolutions scientifiques et techniques depuis 1945}
child[concept color=popblue]{ node{Maîtrise de l'atome}
  child[concept color=popblue!70]{ node{Bombe A 1945 ; bombe H 1952} }
  child[concept color=popblue!70]{ node{Centrales nucléaires civiles} }
}
child[concept color=poporange]{ node{Conquête de l'espace}
  child[concept color=poporange!70]{ node{Spoutnik 1957 ; Gagarine 1961} }
  child[concept color=poporange!70]{ node{Apollo 11 : Lune 1969} }
}
child[concept color=popgreen]{ node{Biologie et médecine}
  child[concept color=popgreen!70]{ node{ADN 1953 ; génie génétique} }
  child[concept color=popgreen!70]{ node{Vaccins, antibiotiques, greffes} }
}
child[concept color=poppurple]{ node{Informatique et internet}
  child[concept color=poppurple!70]{ node{Transistor 1947 ; micro-ordinateur} }
  child[concept color=poppurple!70]{ node{ARPANET 1969 ; Web 1991} }
}
child[concept color=poppink]{ node{Transports et énergie}
  child[concept color=poppink!70]{ node{Avions à réaction, TGV 1981} }
  child[concept color=poppink!70]{ node{Pétrole, barrages, solaire} }
}
child[concept color=popgold]{ node{Agriculture et effets}
  child[concept color=popgold!70]{ node{Révolution verte ; OGM} }
  child[concept color=popgold!70]{ node{Progrès et risques ; Cameroun} }
};
\end{tikzpicture}
\end{center}
\legende{Carte mentale de la leçon : les deux savoirs (grandes découvertes scientifiques, grandes innovations techniques) et leurs effets sur le monde et le Cameroun.}
\end{popfigure}

\begin{aretenir}[L'essentiel de la leçon]
\textbf{Définitions clés}
\begin{itemize}
\item \cle{Révolution scientifique} : bouleversement rapide et profond des connaissances qui transforme la vision du monde (maîtrise de l'atome, structure de l'ADN, conquête de l'espace).
\item \cle{Innovation technique} : application concrète d'une découverte scientifique à la production et à la vie quotidienne (transistor, ordinateur, internet, TGV).
\item \cle{Course à l'espace} : rivalité technologique entre les États-Unis et l'URSS pendant la Guerre froide (Spoutnik en 1957, Apollo 11 en 1969).
\item \cle{Révolution verte} : modernisation de l'agriculture (semences améliorées, engrais, irrigation, mécanisation) pour augmenter les rendements et lutter contre la famine.
\end{itemize}

\textbf{Repères chronologiques à connaître par cœur}
\begin{center}
\begin{tabularx}{\linewidth}{|l|X|}\hline
\rowcolor{popblueL} \textbf{Date} & \textbf{Événement} \\ \hline
1945 & Premières bombes atomiques (Hiroshima et Nagasaki) : début de l'ère nucléaire \\ \hline
1947 & Invention du transistor : naissance de l'électronique moderne \\ \hline
1953 & Découverte de la structure de l'ADN (Watson et Crick) \\ \hline
1957 & Spoutnik 1 : premier satellite artificiel (URSS) \\ \hline
1961 & Youri Gagarine : premier homme dans l'espace (URSS) \\ \hline
1969 & Apollo 11 : premiers hommes sur la Lune (États-Unis) ; naissance d'ARPANET, ancêtre d'internet \\ \hline
1981 & Mise en service du TGV en France : révolution des transports ferroviaires \\ \hline
1991 & Ouverture du Web au public (Tim Berners-Lee) : début de l'ère d'internet pour tous \\ \hline
\end{tabularx}
\end{center}

\textbf{Méthodes express}
\begin{itemize}
\item \emph{Commenter un document} : identifier la nature, l'auteur et la date ; dégager l'idée principale ; relier le document au contexte (Guerre froide, mondialisation) ; conclure sur les effets de la découverte ou de l'innovation.
\item \emph{Rédiger une conclusion} : rappeler le bilan (progrès et risques), puis ouvrir sur les défis actuels (environnement, inégalités, accès au numérique au Cameroun).
\item \emph{Appliquer à la vie courante} : citer un exemple concret (téléphonie mobile, vaccins, électricité des barrages) et expliquer la découverte scientifique qui le rend possible.
\end{itemize}
\end{aretenir}

\begin{aretenir}[Checklist avant le Bac]
$\square$ Je sais définir une révolution scientifique et une innovation technique et les distinguer.

$\square$ Je connais les dates clés : 1945 (atome), 1947 (transistor), 1953 (ADN), 1957 (Spoutnik), 1969 (Lune et ARPANET), 1991 (Web).

$\square$ Je sais expliquer le rôle de la Guerre froide dans la course à l'espace et à l'armement.

$\square$ Je peux citer deux applications pacifiques de l'énergie nucléaire (électricité, médecine) et leurs risques (accidents, déchets radioactifs).

$\square$ Je connais les grandes avancées médicales : antibiotiques, vaccins, greffes d'organes, imagerie médicale.

$\square$ Je sais décrire la filière : transistor $\rightarrow$ ordinateur $\rightarrow$ internet $\rightarrow$ téléphonie mobile.

$\square$ Je peux expliquer ce qu'est la révolution verte et ses limites (coût des engrais, dépendance, pollution de l'environnement).

$\square$ Je sais donner des exemples d'effets des révolutions scientifiques et techniques au Cameroun : barrages hydroélectriques (Songloulou, Édéa, Lagdo), téléphonie mobile, internet, campagnes de vaccination.

$\square$ Je sais construire un commentaire de document en trois parties (présentation, analyse, conclusion).

$\square$ Je sais rédiger une conclusion équilibrée : progrès, risques et défis pour l'avenir.
\end{aretenir}

\findecours

\end{document}
